\documentclass{beamer}
\usepackage{bm}
\usepackage{tabularx}
\setbeamertemplate{caption}[numbered]
\usepackage{xcolor}
\usepackage{multirow}
\usepackage{movie15}
\usepackage{Sweave}
%\usepackage{nameref}% http://ctan.org/pkg/nameref
%\newcommand{\currentchapter}{}
%\let\oldchapter\chapter
%\renewcommand{\chapter}[1]{\oldchapter{#1}\renewcommand{\currentchapter}{#1}}
\renewcommand{\thesection}{2.\arabic{section}}
\newcommand{\currentsection}{}
\let\oldsection\section
\renewcommand{\section}[1]{\oldsection{#1}\renewcommand{\currentsection}{#1}}

\newcommand{\currentsubsection}{}
\let\oldsubsection\subsection
\renewcommand{\subsection}[1]{\oldsubsection{#1}\renewcommand{\currentsubsection}{#1}}

\newcommand{\currentsubsubsection}{}
\let\oldsubsubsection\subsubsection
\renewcommand{\subsubsection}[1]{\oldsubsubsection{#1}\renewcommand{\currentsubsubsection}{#1}}


\newcommand{\boldbeta}{{\boldsymbol \beta}}
\newcommand{\boldtheta}{{\boldsymbol \theta}}
\newcommand{\boldmu}{{\boldsymbol \mu}}

\newenvironment{wideitemize}%
  {\begin{itemize}%
%    \setlength{\itemsep}{0pt}%
    \setlength{\parskip}{1em}}%
  {\end{itemize}}

%\usetheme{CambridgeUS}
\usetheme{Madrid}

%\title[Section~1.\insertsectionnumber \currentsection]{Chapter 1. Motivating Examples}
\title[\currentsection]{Chapter 2.Univariate Generalised Linear Models (GLM)}
\author[univariate GLM]{MAST90084 Statistical Modelling}
\vskip 1cm
\institute[Ch2]{Dennis Leung \\ \vskip 0.3cm SCHOOL OF MATHEMATICS AND STATISTICS\\
\vskip 0.3cm
   THE UNIVERSITY OF MELBOURNE}
\vskip 1cm
\date[Dennis Leung \hspace{8pt}]{}


\begin{document}

\frame{\titlepage}

%
%\begin{frame}
%
%\begin{itemize}
%\item Last time I did make a mistake when I explained (in handwriting) that in ANOVA with orthogonal design matrix, the $\Delta SS$ associated with  a particular covariate is invariant to the order in which the covariates show up in the ANOVA table. Apologize for the confusion. 
%
%\item A student was right about the orthogonal vectors in question need to be normalized before my argument can carry through....
%
%\item Let me fix this today. 
%\end{itemize}
%\end{frame}
%
%
%\begin{frame} {Regression in a LM with orthogonal design}
%\begin{itemize}
%
%\item Setup: We have  a response vector ${y}$ ($n \times 1$), and  $p$ covariate vectors $z_1, \dots z_p$ (each $n \times 1$).
%
%\item Also defined $u_i = z_i /\|z_i\|$ for $i =1, \dots, p$, the normalized versions. 
%
%\item Suppose we only pick a subset  $\{\tilde{z}_1, \dots, \tilde{z}_q\} \subset \{z_1, \dots z_p\}$, and let $\tilde{u}_1, \dots, \tilde{u}_q$ be the corresponding normalized version among the $u_i$'s. 
%
%
%\item Let $\tilde{Z} = [\tilde{z}_1 | ....|\tilde{z}_q ]$ be the design matrix that aligns these $\tilde{z}_i$'s together. 
%
%\item We know that the fitted vector $\tilde{y} =\tilde{Z} (\tilde{Z}' \tilde{Z})^{-1}\tilde{Z}' y = \sum_{i = 1}^q  \langle y , \tilde{u}_i \rangle\tilde{u}_i$. The last equality comes from the orthogonality of the $z_i$'s. 
%\end{itemize}
%
%
%\end{frame}
%
%
%\begin{frame}
%\begin{wideitemize}
%\item So now suppose we also regress $y$ on $\tilde{z}_1, \dots, \tilde{z}_q, z_a$, where $z_a$ is an additional element from the original set $z_1, \dots, z_p$. Also use $u_a$ to denote the normalized version for this $z_a$
%
%\item Analogously we will have fitted value vector $y_a = \sum_{i= 1}^q \langle y , \tilde{u}_i\rangle\tilde{u}_i + \langle y , {u}_a\rangle{u}_a$
%
%\item Now the sum of squared residuals from the first regression is \[\| y- \tilde{y}\|^2 = \langle  y - \tilde{y} , y- \tilde{y}\rangle = \langle y, y\rangle - \sum_{i = 1}^q \langle y , \tilde{u}_i  \rangle^2 \]
%
%\item Analogous, sum of squared residuals from the second regression is
%\[
%\| y- y_a\|^2  =  \langle y, y\rangle - \sum_{i = 1}^q \langle y , \tilde{u}_i  \rangle^2 - \langle y, {u}_a\rangle^2
%\]
%
%\end{wideitemize}
%
%\end{frame}
%
%\begin{frame}
%\begin{wideitemize}
%\item Now it is clear that the additional reduction in $SSE$ brought by $z_a$ is precisely $ \langle y, {u}_a\rangle^2
%$
%
%\item This is independent of what we have picked in our set $\tilde{z}_1, \dots \tilde{z}_q$ in the first place. 
%
%\item The upshot: the reduction in $SSE$ brought by each covariate is independent of the order the covariates are augmented to the rows of the ANOVA table, when we start with an orthogonal design. 
%
%\item But this is NOT true for ANODEV. 
%\end{wideitemize}


%\end{frame}





\begin{frame}[fragile] \frametitle{Outline}
\tableofcontents
\end{frame}
%
%\begin{frame}
%We will continue our discourse on  GLMs
%
%\begin{figure}
% \centering
%\includegraphics[width=0.5\textwidth]{NandW}
%% \label{figure:example}
%\end{figure}
%
%
%\end{frame}
%
%\begin{frame}[fragile] \frametitle{Outline}
%\tableofcontents
%\end{frame}
%%
%%\section{\S2.1 Univariate GLM}\label{sec:2.1}
%%%\subsection{\S2.1.1 Data}
%%%
%%%\begin{frame}{\S2.1.1 Data (ungrouped)}
%%%\begin{itemize}
%%%\item
%%%$y$, univariate response
%%%
%%%\item  $\textbf{x}=(x_1,\cdots,x_m)^T$, $m$ covariate variables
%%%
%%%\item
%%%A sample of $n$ individuals: 
%%%$$(y_i, \textbf{x}_i^T)=(y_i, x_{i1},\cdots, x_{im}), \quad i=1,2,\cdots,n$$
%%% in matrix form
%%%\begin{equation*}
%%%\textbf{y}=\left[\begin{array}{c} y_1 \\ \vdots \\ y_i \\ \vdots \\ y_n\end{array}\right], \quad
%%%X=\left[\begin{array}{ccc} x_{11} & \cdots & x_{1m} \\ \vdots & \cdots & \vdots \\
%%%x_{i1} & \cdots & x_{im} \\ \vdots & \cdots & \vdots \\ x_{n1} & \cdots & x_{nm} \end{array}
%%%\right]_{n\times m}
%%%\end{equation*}
%%%\end{itemize}
%%%\end{frame}
%%%
%%%\begin{frame}{\S2.1.1 Data (grouped)}
%%%\begin{itemize}
%%%\item 
%%%Sometimes two or more individuals have covariate values, i.e. some rows of $X$ are identical
%%%
%%%\item Individuals with  same $\bf x$ (say $n_i$ of them)  can be grouped. The group mean $\bar{y}$ responses is taken to be the group response. 
%%%
%%%
%%%\item $g = $ number of different ${\bf x}_i$
%%%\ \
%%%
%%%
%%%
%%%\begin{center}
%%%\begin{tabular}{cccc}
%%%& sizes & group means &  \\ \\
%%%Group 1 &  \multirow{5}{*}{$\left[\begin{array}{c} n_1 \\ \vdots \\ n_i \\ \vdots \\ n_g\end{array}\right]$,} 
%%%& \multirow{5}{*}{$\bar{\textbf{y}}=\left[\begin{array}{c} \bar{y}_1 \\ \vdots \\ \bar{y}_i \\ \vdots 
%%%\\ \bar{y}_g\end{array}\right]$,} & \multirow{5}{*}{$X=\left[\begin{array}{ccc} x_{11} & \cdots & x_{1m} \\ 
%%%\vdots & \cdots & \vdots \\ x_{i1} & \cdots & x_{im} \\ \vdots & \cdots & \vdots \\ x_{g1} & \cdots & x_{gm}
%%%\end{array}\right]_{g\times m}$} \\ $\vdots$ & & & \\ Group $i$ & & & \\ $\vdots$ & & & \\ Group $g$ & & &
%%%\end{tabular}
%%%\end{center}
%%%
%%%\end{itemize}
%%%\end{frame}
%%%
%%%\subsection{\S2.1.2 GLM}
%%%\begin{frame}{\S2.1.2 GLM (Its form)}
%%%\noindent
%%%\textbf{1. Distributional assumption}
%%%\begin{itemize}
%%%\item
%%%Given $\textbf{x}_i$'s, the $y_i$'s are (conditionally) independent.
%%%
%%%\item
%%%The (conditional) pdf of $y_i$ is
%%%$$f(y_i|\theta_i, \phi, \omega_i)=\exp\left\{\frac{y_i\theta_i-b(\theta_i)}{\phi}\cdot \omega_i +
%%%c(y_i,\phi, \omega_i)\right\}, \quad \mbox{where}$$
%%%\begin{itemize}
%%%\item
%%%$\theta_i$ -- \textit{natural} (or \textit{canonical}) parameter;
%%%\item
%%%$\phi$ --  \textit{scale} (or \textit{dispersion}) parameter;
%%%\item
%%%$b(\cdot)$ and $c(\cdot)$ --  specific functions determined by the actual distribution of $y_i$;
%%%\item
%%%$\omega_i$ is the \textit{weight}. 
%%%\begin{itemize}
%%%\item for ungrouped data, $\omega_i=1$ ; 
%%%\item for grouped data, $\omega_i$ equals $n_i$ when the group mean is the response, and equals $1/n_i$
%%%when group total is the response; 
%%%\end{itemize}
%%%\end{itemize}
%%%\end{itemize}
%%%\end{frame}
%%%
%%%\begin{frame}
%%%
%%%\begin{itemize}
%%%\item Jorgensen (1992) also calls this an exponential dispersion model
%%%\item $\mu_i = E(y_i|\textbf{x}_i)$ still denotes the mean
%%%\item We will model the dependence of $\mu_i $ on covariates via a link/response function
%%%%The (conditional) distribution of $y_i$ belongs to a \textit{simple exponential family} with  
%%%%$E(y_i|\textbf{x}_i)=\mu_i$, and possibly a common \textit{scale} (or \textit{dispersion}) parameter $\phi$.
%%%\end{itemize}
%%%
%%%\end{frame}
%%%
%%%
%%%\begin{frame}{\S2.1.2 GLM (Mean structure via link)}
%%%\noindent
%%%\textbf{2. Structural assumption}
%%%
%%%Mean $\mu_i$ is related to the \textit{linear predictor} $\eta_i=\textbf{z}_i^T\mbox{\boldmath $\beta$}$ by
%%%$$\mu_i=h(\eta_i)=h(\textbf{z}_i^T\mbox{\boldmath $\beta$}), \quad\mbox{respectively},\quad
%%%\eta_i=g(\mu_i)$$
%%%where
%%%\begin{itemize}
%%%\item
%%%$h(\cdot)$ is the \textit{response function};
%%%\item
%%%$g(\cdot)$ is the \textit{link function}, $g=h^{-1}$;
%%%\item
%%%$\mbox{\boldmath $\beta$}$ is a $p\times 1$ unknown \textit{parameter} vector;
%%%\item
%%%$\textbf{z}_i$ is a $p\times 1$ unknown \textit{design} vector, $\textbf{z}_i=\textbf{z}(\textbf{x}_i)$ (a transformation of $\textbf{x}_i$ ).
%%%\end{itemize}
%%%
%%%
%%%\end{frame}
%%%
%%%
%%%\begin{frame}
%%%\begin{itemize}
%%%
%%%\item
%%%Thus a  GLM has  3 components
%%%\begin{enumerate}
%%%\item
%%%type of the exponential family
%%%\item
%%%the response/link function
%%%\item
%%%the design vector 
%%%\end{enumerate}
%%%\item 
%%%Obviously, ${\bf z}(\textbf{x}_i)$ determines the mean $\mu_i$ via the link
%%%
%%%\item The exponential family construction must also have led to a relationship between $\mu_i$ and the exponential family parameters. What is this relationship?
%%%\end{itemize}
%%%
%%%\end{frame}
%%%
%%%\begin{frame}{\S2.1.2 GLM (Mean determines variance)}
%%%\noindent
%%%\textbf{Remark 1.}
%%%\begin{itemize}
%%%\item
%%%One can derive that
%%%$$\mu_i=b^\prime (\theta_i)=\frac{\partial b(\theta_i)}{\partial \theta_i} \quad
%%%\mbox{and}\quad \mbox{var} (y_i|\textbf{x}_i)\stackrel{\mbox{\scriptsize denoted}}{=}\sigma^2(\mu_i)=\frac{\phi}{\omega_i} v(\mu_i)$$
%%%with
%%%$v(\mu)=b^{\prime\prime}(\theta)=\frac{\partial^2 b(\theta)}{\partial \theta^2}$ being called the 
%%%\textit{variance function}.
%%%
%%%\item (see blackboard derivation)
%%%
%%%\item  $Note: $$b'$ is invertible mapping from $\theta$ to $\mu$, since $b''(\theta)$, being a variance, is $>0$.  (Hence the notation $v(\mu)$)
%%%\item
%%%$\Rightarrow$ Modeling the  mean  as $\mu=h(\textbf{z}^T\mbox{\boldmath $\beta$})$  also determines the variance structure
%%%
%%%\end{itemize}
%%%\end{frame}
%%%
%%%\begin{frame}
%%%\begin{itemize}
%%%\item
%%%different from linear models (LM) where $E(y|\textbf{x})$ and 
%%%$\mbox{var}(y|\textbf{x})$ are unrelated. 
%%%\item \textbf{quasi-likelihood} approach (Ch2.3) allows us to model the mean and variance structure separately. 
%%%%The relation between $\mu$ and $v(\mu)$ may also be specified through empirical evidence,
%%%%resulting in the so-called \textbf{quasi-likelihood} approach.
%%%\end{itemize}
%%%
%%%\end{frame}
%%%
%%%\begin{frame}{\S2.1.2 GLM (Natural/canonical link)}
%%%\noindent
%%%\textbf{Remark 2.}
%%%\begin{wideitemize}
%%%\item
%%%Choosing $h(\cdot)$ and $g(\cdot)$ should depend on the specific distribution of $y$.
%%%\item
%%%Note that assuming $\theta\equiv \theta(\mu)=\eta\equiv \textbf{z}^T\mbox{\boldmath $\beta$}$
%%%makes estimating $\mbox{\boldmath $\beta$}$ much easier.
%%%
%%%%Since $\mu=b^\prime(\theta)$ implies $\theta=(b^\prime)^{-1}(\mu)$, it follows that
%%%%$\eta=\theta(\mu)=(b^\prime)^{-1}(\mu)$ and $g(\mu)=(b^\prime)^{-1}(\mu)$ under this assumption.
%%%\item
%%%We call $\eta=(b^\prime)^{-1}(\mu)$ the \textit{natural} (or \textit{canonical}) link function,
%%%and $\mu=b^{\prime}(\eta)$ the \textit{natural response function}.
%%%%\item
%%%%It can be shown that the natural link function is
%%%%\begin{enumerate}
%%%%\item
%%%%$\eta=\mu$ if $y\sim \mbox{Normal}(\mu, \sigma^2)$; \textit{identity link};
%%%%\item
%%%%$\eta=\log\mu$ if $y\sim \mbox{Poisson}(\mu)$; \textit{log link};
%%%%\item
%%%%$\eta=\log\frac{p}{1-p}=\log\frac{\mu}{n_g-\mu}$ if $y\sim\mbox{Binomial}(n_g, p)$; \textit{logit (or logistic) link}.
%%%%\end{enumerate}
%%%\end{wideitemize}
%%%\end{frame}
%%%
%%%\begin{frame}{\S2.1.2 GLM (design vector specifics)}
%%%\noindent
%%%\textbf{Remark 3.}
%%%\begin{itemize}
%%%\item
%%%Concerning the design vector, an intercept term is often included
%%%$$\textbf{z}=(1,\textbf{z}^{*T})^T$$
%%%where $\textbf{z}^*=\textbf{z}^*(\textbf{x})$ is a function of $\textbf{x}$, including e.g. $\log(x_1)$,
%%%$x_2^2$, $\sin (x_m)$ etc. as well as $\textbf{x}$ itself.
%%%%\item
%%%%A categorical covariate with $k$ levels has to be coded by a $(k-1)\times 1$ dummy vector, or an equivalent
%%%%vector of size $q=k-1$ for it to be properly included in the design vector.
%%%\item A categorical covariate with $k$ levels can  be coded by a dummy vector ${\bf x} = (x^{(1)}, \dots, x^{(k-1)})$ of length $k-1$, where
%%%\[ x^{(j)} =
%%%  \begin{cases}
%%%    1      & \quad \text{if } j = 1, \dots, k-1\\
%%%   0  & \quad \text{if } j = k
%%%  \end{cases}
%%%\]
%%%(k-th category reserved as reference)
%%%\end{itemize}
%%%\end{frame}
%%%
%%%\begin{frame}{\S2.1.2 GLM (Grouped and ungrouped)}
%%%\noindent
%%%\textbf{Remark 4.}
%%%\begin{itemize}
%%%
%%%\item 
%%%Suppose ungrouped data 
%%%$(y_i, \textbf{x}_i^T)$, $i=1,\cdots, n$, are modelled by a GLM with
%%%$$E(y_i|\textbf{x}_i)=\mu_i=h(\textbf{z}^T\mbox{\boldmath $\beta$}), \quad
%%%\mbox{var}(y_i|\textbf{x}_i)=\phi v(\mu_i).$$
%%%\item
%%%If the data are grouped, so  $\bar{y}_i$ is  the response observation for group $i$ of
%%%size $n_i$ with covariate vector $\textbf{x}_i$, then
%%%% the distribution of $\bar{y}_i$ is within the exponential
%%%%family again, and
%%%$$E(\bar{y}_i|\textbf{x}_i)=\mu_i=h(\textbf{z}^T\mbox{\boldmath $\beta$}), \quad
%%%\mbox{var}(\bar{y}_i|\textbf{x}_i)=n_i^{-1}\phi v(\mu_i).$$
%%%So the same GLM applies to the grouped data as well if the variance function is properly adjusted 
%%%by a factor of $n_i^{-1}$, i.e. defining the weight $\omega_i=n_i$.
%%%\item By same token, if the group \emph{total} is used as response, one should define $\omega_i=1/n_i$.
%%%\end{itemize}
%%%\end{frame}
%%
%%%\begin{frame}{GLM  Common exponential families}
%%%
%%%Recall the table from last time
%%%\begin{table} [!htbp]
%%%\caption{Common exponential families} \label{tab2.1}
%%%\begin{center}
%%%{\tiny
%%%\begin{tabular}{|c|c|c|c|c|c|c|}
%%%\hline
%%%Distribution & $\theta (\mu)$ & $b(\theta)$ & $\phi$ & $E(y)=b^\prime(\theta)$ & $v(\mu)=b^{\prime\prime}(\theta)$
%%%& $\mbox{var}(y)=b^{\prime\prime}(\theta)\phi/\omega$ \\ \hline
%%%$\begin{array}{c} \mbox{Normal} \\ N(\mu,\sigma^2)\end{array}$ & $\mu$ & $\theta^2/2$ & $\sigma^2$ 
%%%& $\mu=\theta$ & 1 & $\sigma^2/\omega$ \\
%%%\hline
%%%$\begin{array}{c} \mbox{Bernoullil} \\ B(1,\pi)\end{array}$ & $\log\frac{\pi}{1\!-\!\pi}$ 
%%%& $\log (1\!+\!e^\theta)$ & 1 & $\pi=\frac{e^\theta}{1+e^\theta}$ &
%%%$\pi(1-\pi)$ & $\pi(1-\pi)/\omega$ \\ \hline
%%%$\begin{array}{c} \mbox{Poisson} \\ P(\lambda)\end{array}$ & $\log\lambda$ & $e^\theta$ 
%%%& 1 & $\lambda=e^\theta$ & $\lambda$ & $\lambda/\omega$ \\
%%%\hline
%%%$\begin{array}{c} \mbox{Gamma} \\ G(\mu,\nu)\end{array}$ & $-1/\mu$ & $-\log(-\theta)$ & $\nu^{-1}$ 
%%%& $\mu=-1/\theta$ & $\mu^2$ & $\mu^2\nu^{-1}/\omega$ \\ \hline
%%%$\begin{array}{c} \mbox{InverseGaussian} \\ IG(\mu,\sigma^2)\end{array}$ & $-1/2\mu^2$ & $-(-2\theta)^{1/2}$
%%%& $\sigma^2$ &  $\mu=(-2\theta)^{-1/2}$ & $\mu^3$ & $\mu^3\sigma^2/\omega$ \\ \hline
%%%\end{tabular}}
%%%\end{center}
%%%\end{table}
%%%%Remember that if $\theta = \eta = \textbf{z}^T\mbox{\boldmath $\beta$}$
%%%%\begin{itemize}
%%%%\item 2nd column gives the canonical link function
%%%%\item 3rd column gives the canonical response function
%%%%\end{itemize}
%%%\end{frame}
%%
%%
%%%\subsection{\S2.1.3 GLM for continuous responses}\label{sec:2.1.3}
%%%\begin{frame}{\large \S2.1.3 GLM for normal, Gamma \& inverse Gaussian responses }
%%%
%%%When $y\sim \mbox{Normal}(\mu,\sigma^2)$, $\mu=\eta=\textbf{z}^T\mbox{\boldmath $\beta$}$
%%%if the natural link is used. Also $\theta=\mu$, $b(\theta)=\frac{1}{2}\theta^2$, $\phi=\sigma^2$ and
%%%$v(\mu)=1$.
%%%
%%%\end{frame}
%%%
%%%\begin{frame}{\large \S2.1.3 GLM for normal, Gamma \& inverse Gaussian responses}
%%%When $y$ has a Gamma distribution (e.g. survival time, insurance claims, amount of rainfall, etc.), the pdf is
%%%$$f(y|\mu,\nu)=\frac{1}{\Gamma(\nu)}\left(\frac{\nu}{\mu}\right)^\nu y^{\nu-1}\exp\left(-\frac{\nu}{\mu}y\right),
%%%\quad y\geq 0,$$
%%%where $\mu>0$ is the mean parameter, and $\nu>0$ the shape parameter.
%%%\begin{itemize}
%%%\item
%%%The natural link is $\eta=-\theta=\mu^{-1}$. Linear predictor $\eta=\textbf{z}^T\mbox{\boldmath $\beta$}$.
%%%Hence $\mu^{-1}=\textbf{z}^T\mbox{\boldmath $\beta$}$.
%%%\item
%%%Alternatively, identical link $\eta=\mu$ and log-link $\eta=\log \mu$ can be used to avoid the restriction. (These will give positive mean $\mu$ so they are allowed)
%%%\end{itemize}
%%%\end{frame}
%%%
%%%\begin{frame}{\large \S2.1.3 GLM for normal, Gamma \& inverse Gaussian responses }
%%%\begin{itemize}
%%%\item
%%%When $y$ has an inverse Gaussian distribution $\mbox{IG}(\mu,\sigma^2)$ (useful for modelling life time etc.), the pdf is
%%%\begin{eqnarray*}
%%%f(y|\mu,\sigma^2) &\!\!\!=\!\!\!&\left(\frac{1}{2\pi\sigma^2 y^3}\right)^{1/2}\exp\left(-\frac{(y-\mu)^2}{2\sigma^2\mu^2y}
%%%\right), \quad\!\!\! y>0, \mu>0, \sigma^2>0 \\
%%%&\!\!\!=\!\!\!& \left(\frac{1}{2\pi\sigma^2 y^3}\right)^{1/2}\exp\left\{\frac{-y\frac{1}{2\mu^2}+\frac{1}{\mu}}{\sigma^2}
%%%-\frac{1}{2\sigma^2 y}\right\}
%%%\end{eqnarray*}
%%%\item
%%%Hence natural parameter $\theta=-\frac{1}{2\mu^2}$;  dispersion parameter $\phi=\sigma^2$.
%%%\begin{eqnarray*}
%%%b(\theta)=-\frac{1}{\mu}=-(-2\theta)^{1/2} \quad \Rightarrow \quad b^\prime(\theta)=(-2\theta)^{-1/2}=\mu \\
%%%\Rightarrow \; b^{\prime\prime}(\theta)\!=\!(-2\theta)^{-3/2}\!=\!\mu^3 \;\; \Rightarrow \;\;
%%%v(\mu)\!=\!\mu^3 \;\; \Rightarrow \;\; \mbox{Var}(y)\!=\!\sigma^2\mu^3
%%%\end{eqnarray*}
%%%The natural link is obtained by setting $\eta=\theta, \quad\Rightarrow\quad \eta=-\frac{1}{2\mu^2}$, i.e.
%%%$$\eta=\theta=(b^\prime)^{-1}(\mu)=-(2\mu^2)^{-1}.$$
%%%In practice $\eta=\mu^{-2}$ is used as the natural link.
%%%\end{itemize}
%%%\end{frame}
%%
%%\subsection{\S2.1.4 GLM for binary and binomial responses}\label{sec:2.1.4}
%%\begin{frame}{\S2.1.4 GLM for binary and binomial responses }
%%\begin{wideitemize}
%%\item
%%Suppose $y_i|\textbf{x}_i \sim \mbox{Binomial}(m_i,\pi_i)$ and $\bar{y}_i=\frac{y_i}{m_i}$ (scaled binomial) are used as the
%%response, $i=1,\cdots,n$. Then the pmf is {\footnotesize
%%$$f(y_i|m_i, \pi_i,\textbf{x}_i)\!=\!{m_i\choose y_i}\pi_i^{y_i}(1-\pi_i)^{m_i\!-\!y_i}\!=\!{m_i\choose y_i}\exp\left\{
%%m_i\bar{y}_i\log\frac{\pi_i}{1\!-\!\pi_i}+m_i\log (1\!-\!\pi_i)\right\}$$}
%%\item
%%Natural parameter $\theta_i=\log\frac{\pi_i}{1-\pi_i}$; dispersion $\phi=1$; weight $\omega_i=m_i$.
%%\item The scaled binomial distribution may treated as the average from grouped binary observations. (group size $ = m_i$ )
%%\end{wideitemize}
%%\end{frame}
%%
%%
%%
%%\begin{frame}
%%\begin{itemize}
%%\item
%%$\mu_i^*=E(y_i|\textbf{x}_i)=m_i\pi_i=m_i\frac{\exp(\theta_i)}{1+\exp(\theta_i)}$
%%{\small $$\Rightarrow \quad \mu_i=E(\bar{y}_i|\textbf{x}_i)=\pi_i=\frac{\exp(\theta_i)}{1+\exp(\theta_i)}$$}
%%\item
%%Also $b(\theta_i)=-\log (1-\pi_i)=\log (1+e^{\theta_i}) \quad \Rightarrow \quad 
%%b^\prime(\theta_i)=\frac{e^{\theta_i}}{1+e^{\theta_i}}=\pi_i$
%%{\small $$\Rightarrow\quad b^{\prime\prime}(\theta_i)=\frac{e^{\theta_i}}{(1+e^{\theta_i})^2}=\pi_i(1-\pi_i)
%%=v(\pi_i)$$}
%%$\Rightarrow\quad \mbox{Var}(\bar{y}_i|\textbf{x}_i)=\sigma^2(\mu_i)=\phi v(\mu_i)/\omega_i
%%=m_i^{-1}\pi_i(1-\pi_i)$.
%%
%%\end{itemize}
%%\end{frame}
%%
%%
%%
%%\begin{frame}{\S2.1.4 GLM for binary and binomial responses }
%%\begin{itemize}
%%\item
%%The natural link is the logit link
%%\[\eta_i=\theta_i=\log\frac{\pi_i}{1-\pi_i}=(b^\prime)^{-1}(\pi_i),\]
%%called the \textbf{logit} (or \textbf{logistic}) \textbf{link}, with the corresponding {\bf logistic} response function
%%\[
%%\pi_i = \frac{\exp(\eta_i)}{1 + \exp(\eta_i)}. 
%%\]
%%
%%This gives the well-known \textit{logit regression model}:
%%$\displaystyle \log\frac{\pi_i}{1-\pi_i}\!=\!\textbf{z}_i^T\!\mbox{\boldmath $\beta$}$
%%
%%
%%
%%
%%\end{itemize}
%%\end{frame}
%%
%%
%%
%%\begin{frame}
%%
%%Other models of the form 
%%\[
%%\pi = h(\eta)
%%\]
%%for  strictly monotonous distribution functions $h$  on the real line:
%%\begin{enumerate} 
%%\item
%%\textbf{Probit Model}:  \[\pi_i = \Phi(\textbf{z}_i^T\!\mbox{\boldmath $\beta$}), \] where $\Phi(\cdot)$ is the N(0,1) cdf. 
%%% $\eta_i=\Phi^{-1}(\pi_i)$, where 
%% 
%%% This results in the
%%%\textit{probit model} $$\Phi^{-1}(\pi_i)=\textbf{z}_i^T\!\mbox{\boldmath $\beta$}.$$
%%\item
%%\textbf{Complementary log-log Model}:
%%\[
%%\pi_i = h(\eta_i)\!=\!1\!-\!\exp(\!-\!\exp(\!\eta_i\!)\!)
%%\]
%%with link function 
%%$
%% \textbf{z}_i^T\!\mbox{\boldmath $\beta$}= \log (-\log(1-\pi_i)).
%%$
%%%
%%%
%%% $\eta_i=g(\pi_i)=\log (-\log(1-\pi_i))$, resulting in the
%%%\textit{complementary log-log model} 
%%%$$\log (-\log(1-\pi_i))=\textbf{z}_i^T\!\mbox{\boldmath $\beta$}.$$
%%%Note the corresponding \textbf{response function} is $h(\eta_i)\!=\!1\!-\!\exp(\!-\!\exp(\!\eta_i\!)\!)$.
%%\item
%%\textbf{Complementary log Model} 
%%\[
%% \pi_i = h(\eta_i)  = 
%%  \begin{cases} 
%%   1 - \exp(- \eta_i) & \text{if } \eta_i \geq 0 \\
%%   0       & \text{if } \eta_i < 0
%%  \end{cases}
%%\]
%%
%%%\[
%%%\pi_i = h(\eta_i) = 
%%%\eta_i=g(\pi_i)=-\log(1-\pi_i),
%%%\] 
%%%resulting in the
%%%\textit{complementary log model}
%%with link function $-\log(1-\pi_i)=\textbf{z}_i^T\!\mbox{\boldmath $\beta$}.$
%%%\item
%%%Any inverse cdf can serve as a link function for GLM with binomial responses.
%%\end{enumerate}
%%\end{frame}
%%
%%
%%\begin{frame}
%%
%%
%%
%%\begin{figure}
%% \centering
%%\includegraphics[width=1\textwidth]{BinaryResponseFunctions.png}
%% \label{figure:example}
%%\end{figure}
%%
%%\end{frame}
%%
%%\begin{frame}{\S2.1.4 Over-dispersion in binomial GLM}
%%\begin{itemize}
%%
%%\item In applications it often happens that $\mbox{var}(\bar{y}_i|\textbf{x}_i) > m_i^{-1}\pi_i(1-\pi_i)$,  making the binomial modeling not accurate. 
%%
%%\item Two main reasons for this  \textbf{over-dispersion}:
%%\begin{itemize}
%%\item
%%Unobserved heterogeneity, i.e. different Bernoulli trials in the data have different probabilities of ``success".
%%\item
%%Positive correlation between individual Bernoulli trials. (For example, different units in a group may be correlated, such as members in a household)
%%\end{itemize}
%%%
%%%\item
%%%Sometimes $\mbox{var}(\bar{y}_i|\textbf{x}_i) > m_i^{-1}\pi_i(1-\pi_i)$, so-called \textbf{over-dispersion},
%%%where $\bar{y}_i=m_i^{-1}y_i$ and $y_i$ is the number of ``successes" in $m_i$ Bernoulli trials.
%%%This implies the actual distribution of $y_i$ is not binomial. Two causes for this happening in practice.
%%%\begin{enumerate}
%%%\item
%%%Unobserved heterogeneity, i.e. different Bernoulli trials in the data have different probabilities of ``success".
%%%\item
%%%Positive correlation between individual Bernoulli trials.
%%%\end{enumerate}
%%
%%%\[\mbox{var}(\bar{y}_i|\textbf{x}_i) =\phi\cdot\frac{\pi_i(1-\pi_i)}{m_i},\quad
%%%\mbox{with $\phi>0$ an over-dispersion parameter.}
%%%\]
%%
%%%Since the actual distribution for $y_i$ is difficult to get in this situation, the simplest way to account
%%%for this over-dispersion is assuming
%%%$$\mbox{var}(\bar{y}_i|\textbf{x}_i) =\phi\cdot\frac{\pi_i(1-\pi_i)}{m_i},\quad
%%%\mbox{with $\phi>0$ an over-dispersion parameter.}$$
%%%\item
%%%\textbf{Remark}. Assuming a \textbf{beta-binomial distribution} for $y_i$ can be used to account for the
%%%over-dispersion.
%%\end{itemize}
%%\end{frame}
%%
%%\begin{frame}
%%\begin{itemize}
%%\item
%%One can account for this by assuming 
%%\[\mbox{var}(\bar{y}_i|\textbf{x}_i) =\phi\cdot\frac{\pi_i(1-\pi_i)}{m_i},\quad
%%\mbox{with $\phi>0$ an over-dispersion parameter.}
%%\]
%%with the over-dispersion parameter $\phi>1$
%%
%%\item Remark: Not same as setting $\phi \not= 1$ in the exponential family form of the binomial distribution! You wouldn't get a proper density this way. However, a genuine likelihood is not actually needed for inference; more on this later when we cover \emph{quasi-likelihood}. 
%%
%%%
%%%
%%%\item However, 
%%
%%
%%\item One  may assume  a \textbf{beta-binomial distribution}, which allow for variance of the form $\phi \frac{\pi_i(1-\pi_i)}{m_i}$
%%\end{itemize}
%%
%%\end{frame}
%%
%%\subsection{\S2.1.5 GLM for count responses}\label{sec:2.1.5}
%%\begin{frame}{\S2.1.5 GLM for count responses}
%%\begin{itemize}
%%\item
%%If $y_i|\textbf{x}_i \sim \mbox{Poisson}(\lambda=\mu_i)$ with $E(y_i|\textbf{x}_i)=\mu_i$
%%and $\mbox{var}(y_i|\textbf{x}_i)=\mu_i$.
%%
%%\item The natural/canonical parameter is
%%$\theta_i=\log\mu_i$
%%
%%\item This leads to the natural link  \textbf{log-link} $\eta_i=\log\mu_i$, resulting in
%%the \textbf{log linear model}:
%%\[
%%\log\mu_i=\eta_i=\textbf{z}_i^T\!\mbox{\boldmath $\beta$} \quad \iff \quad
%%\mu_i=\exp(\textbf{z}_i^T\!\mbox{\boldmath $\beta$})=\exp(\eta_i).
%%\]
%%
%%
%%\item Overdispersion can occur for similar reasons as for binomial data; can be accounted for
%%by introducing an over-dispersion parameter $\phi$ such that $\mbox{var}(y_i|\textbf{x}_i)=\phi\mu_i$.
%%
%%\item (a topic of quasi-likelihood later)
%%
%%\item Again, some distributions for count data do exist in the exponential family that allow for over-dispersion, such as the negative binomial distribution. 
%%\end{itemize}
%%\end{frame}
%%
%%\section{\S2.2 Likelihood Inference}\label{sec:2.2}
%%\begin{frame}{\S2.2 Likelihood Inference}
%%\begin{itemize}
%%\item
%%Regression analysis with GLMs is based on \textbf{likelihood maximization}. 
%%\item
%%It includes 
%%\begin{itemize}
%%\item parameter point estimation, 
%%\item confidence intervals, 
%%\item hypothesis testing,
%%\item  goodness-of-fit tests and 
%%\item  variable selection (or model selection).
%%
%%\end{itemize}
%%\item
%%For now we assume the model is completely and correctly specified, i.e., the GLM assumption is ``true".
%%
%%\item 
%%This can  be relaxed when we cover quasi-likelihood. 
%%\end{itemize}
%%\end{frame}
%\section{\S2.2 Likelihood Inference}\label{sec:2.2}
%%\subsection{\S2.2.1 Maximum likelihood estimation}\label{sec:2.2.1}
%%\begin{frame}{\S2.2.1 Maximum likelihood estimation}
%%\begin{itemize}
%%\item
%%Data:  response $y_i$, covariates $\textbf{x}_i$ and design
%%vector $\textbf{z}_i = {\bf z}_i(\textbf{x}_i)$, $i=1,2,\cdots, n$.
%%\item
%%Model:  GLM where $E(y_i|\textbf{x}_i)=\mu_i=h(\eta_i)=h(\textbf{z}_i^T\!\mbox{\boldmath $\beta$})$,
%%with response function $h(\cdot)$ and unknown coefficient $p\times 1$ vector $\mbox{\boldmath $\beta$}$. 
%%\item
%%Assumption:  design matrix $Z = Z_{n\times p}=(\textbf{z}_1,\cdots,\textbf{z}_n)^T$ has (full) rank $p$
%%and scale/dispersion parameter $\phi$ is
%%assumed known or otherwise consistently estimated. 
%%%\item
%%%Now the log-likelihood function is {\footnotesize
%%%\begin{eqnarray*}
%%%\ell(\mbox{\boldmath $\beta$})&=&\sum_{i=1}^n\ell_i(\mbox{\boldmath $\beta$})
%%%=\sum_{i=1}^n f(y_i|\theta_i,\phi,\omega_i) \\
%%%&=&\sum_{i=1}^n \frac{y_i\theta_i-b(\theta_i)}{\phi}\cdot\omega_i
%%%+\mbox{terms not involving $\theta_i$'s} \\
%%%&=& \sum_{i=1}^n \frac{y_i\theta (\mu_i)-b(\theta(\mu_i))}{\phi}\cdot\omega_i
%%%=\sum_{i=1}^n \frac{y_i\theta (h(\textbf{z}_i^T\!\mbox{\boldmath $\beta$}))-
%%%b(\theta(h(\textbf{z}_i^T\!\mbox{\boldmath $\beta$})))}{\phi}\cdot\omega_i
%%%\end{eqnarray*}}
%%\end{itemize}
%%
%%
%%\end{frame}
%%
%%\begin{frame}{\S2.2.1 Maximum likelihood estimation}
%%
%%The focus is estimating $\mbox{\boldmath $\beta$}$. 
%%Now the log-likelihood function is {\footnotesize
%%\begin{eqnarray*}
%%\ell(\mbox{\boldmath $\beta$})&=&\sum_{i=1}^n\ell_i(\mbox{\boldmath $\beta$})
%%=\sum_{i=1}^n f(y_i|\theta_i,\phi,\omega_i) \\
%%&=&\sum_{i=1}^n \frac{y_i\theta_i-b(\theta_i)}{\phi}\cdot\omega_i
%%+\mbox{terms not involving $\theta_i$'s} \\
%%&=& \sum_{i=1}^n \frac{y_i\theta (\mu_i)-b(\theta(\mu_i))}{\phi}\cdot\omega_i
%%=\sum_{i=1}^n \frac{y_i\theta (h(\textbf{z}_i^T\!\mbox{\boldmath $\beta$}))-
%%b(\theta(h(\textbf{z}_i^T\!\mbox{\boldmath $\beta$})))}{\phi}\cdot\omega_i
%%\end{eqnarray*}}
%%
%%\end{frame}
%%
%%\begin{frame}{\S2.2.1 Maximum likelihood estimation (2)}
%%\begin{itemize}
%%\item
%%The \textbf{score function} is the first derivative: {\footnotesize
%%\begin{eqnarray*}
%%\textbf{s}(\mbox{\boldmath $\beta$})&\!\!\!=\!\!\!& \frac{\partial \ell(\mbox{\boldmath $\beta$})}{\partial\mbox{\boldmath 
%%$\beta$}}=\sum_{i=1}^n\frac{\partial \ell_i(\mbox{\boldmath $\beta$})}{\partial\mbox{\boldmath $\beta$}}
%%=\sum_{i=1}^n  \frac{\partial \ell_i(\mbox{\boldmath $\beta$})}{\partial\theta_i}\cdot
%%\frac{\partial \theta_i}{\partial\mu_i}\cdot \frac{\partial \mu_i}{\partial\eta_i} \cdot
%%\frac{\partial \eta_i}{\partial\mbox{\boldmath $\beta$}}\\
%%&\!\!\!=\!\!\!&\sum_{i=1}^n \left(\frac{y_i-b^\prime(\theta_i)}{\phi}\cdot\omega_i\right)\cdot
%%\left[v(h(\textbf{z}_i^T\!\mbox{\boldmath $\beta$}))\right]^{-1}\cdot
%% \frac{\partial h(\eta_i)}{\partial\eta_i}\cdot \textbf{z}_i \\
%%&\!\!\!=\!\!\!& \sum_{i=1}^n \textbf{z}_i\cdot \frac{\partial h(\eta_i)}{\partial\eta_i}\cdot
%%\left[\frac{\phi v(h(\textbf{z}_i^T\!\mbox{\boldmath $\beta$}))}{\omega_i}\right]^{-1}\cdot
%%\left[y_i-\mu_i(\mbox{\boldmath $\beta$})\right] \\
%%&\!\!\!\stackrel{denoted}{=} \!\!\!& \sum_{i=1}^n \textbf{z}_i D_i(\mbox{\boldmath $\beta$})
%%\sigma_i^{-2}(\mbox{\boldmath $\beta$}) \left[y_i-\mu_i(\mbox{\boldmath $\beta$})\right] \\
%%&\!\!\!=\!\!\!& \left[\textbf{z}_1, \cdots, \textbf{z}_n\right]  \left[\!\begin{array}{ccc}
%%\!\!\!D_1(\mbox{\boldmath $\beta$})\!\! & & \\
%%& \!\!\!\ddots\!\!\! & \\ & & \!\!\!D_n(\mbox{\boldmath $\beta$})\!\!\!\end{array}\!\right] 
%%\left[\begin{array}{ccc}\!\!\!\!\sigma_1^2(\mbox{\boldmath $\beta$}) \!\!\! & & \\
%%&\!\!\!\!\ddots \!\!\!\!& \\ & & \!\!\!\!\sigma_n^2(\mbox{\boldmath $\beta$})\!\!\!\end{array}\right]^{\!\!-1}
%%\left[\begin{array}{c} \!\!\!y_1\!-\!\mu_1(\mbox{\boldmath $\beta$}) \!\!\!\! \\
%%\!\!\!\vdots \!\!\!\! \\ \!\!\! y_n\!-\!\mu_n(\mbox{\boldmath $\beta$})\!\!\!\!\end{array}\!\right] \\
%%&\!\!\!\stackrel{\mbox{matrix form}}{=} \!\!\!& Z^TD(\mbox{\boldmath $\beta$})\Sigma^{-1}(\mbox{\boldmath $\beta$})
%%\left[\textbf{y}-\mbox{\boldmath $\mu$}(\mbox{\boldmath $\beta$})\right]
%%\end{eqnarray*}}
%%\end{itemize}
%%\end{frame}
%%
%%\begin{frame}{\S2.2.1 Maximum likelihood estimation (3)}
%%Supporting arguments for the derivation of $s(\mbox{\boldmath $\beta$})$
%%\begin{itemize}
%%\item
%%$\displaystyle \mu_i=b^\prime(\theta_i) \quad \Rightarrow\quad \frac{\partial \mu_i}{\partial\theta_i}=
%%b^{\prime\prime}(\theta_i)=v(\mu_i)=v(h(\textbf{z}_i^T\!\mbox{\boldmath $\beta$}))$
%%\item
%%Also $\displaystyle \frac{\partial \theta_i}{\partial\mu_i}=\left(\frac{\partial \mu_i}{\partial\theta_i}\right)^{-1}$ 
%%and $\displaystyle \frac{\partial \eta_i}{\partial\mbox{\boldmath $\beta$}}=\textbf{z}_i$.
%%\item
%%{\footnotesize Denote $\displaystyle D_i(\mbox{\boldmath $\beta$})=\frac{\partial \mu_i}{\partial\eta_i}
%%=\frac{\partial h(\eta_i)}{\partial\eta_i}\mid_{\eta_i=\textbf{z}_i^T\!\mbox{\boldmath $\beta$}}$,\;
%%$\sigma_i^2(\mbox{\boldmath $\beta$})=\omega_i^{-1}\phi v(h(\textbf{z}_i^T\!\mbox{\boldmath $\beta$}))$
%%,
%%$i=1,\cdots, n$ (or $i=1,\cdots, g$ for grouped data).}
%%\item
%%{\footnotesize Denote $\displaystyle D(\mbox{\boldmath $\beta$})=\left[\begin{array}{ccc}
%%\!\!\!D_1(\mbox{\boldmath $\beta$})\!\! & & \\
%%& \!\!\!\ddots\!\!\! & \\ & & \!\!\!D_n(\mbox{\boldmath $\beta$})\!\!\!\end{array}\right]=\mbox{diag}
%%\{D_1(\mbox{\boldmath $\beta$}), \cdots, D_n(\mbox{\boldmath $\beta$})\}$}
%%\item
%%{\footnotesize Denote $\displaystyle \Sigma(\mbox{\boldmath $\beta$})=\left[\begin{array}{ccc}
%%\!\!\!\sigma_1^2(\mbox{\boldmath $\beta$})\!\! & & \\
%%& \!\!\!\ddots\!\!\! & \\ & & \!\!\!\sigma_n^2(\mbox{\boldmath $\beta$})\!\!\!\end{array}\right]=\mbox{diag}
%%\{\sigma_1^2(\mbox{\boldmath $\beta$}), \cdots, \sigma_n^2(\mbox{\boldmath $\beta$})\}$}
%%
%%
%%\end{itemize}
%%\end{frame}
%%
%%\begin{frame}{\S2.2.1 Maximum likelihood estimation (4)}
%%
%%The \textbf{Fisher information} matrix is the covariance of the score function at the point $\mbox{\boldmath $\beta$}$: 
%%\begin{eqnarray*}
%%F(\mbox{\boldmath $\beta$})&=& \mbox{cov}\left(\textbf{s}(\mbox{\boldmath $\beta$})\right)=
%%\sum_{i=1}^n \textbf{z}_i D_i^2(\mbox{\boldmath $\beta$})\sigma_i^{-2}(\mbox{\boldmath $\beta$}) 
%%\textbf{z}_i^T=\sum_{i=1}^n \textbf{z}_i w_i(\mbox{\boldmath $\beta$})\textbf{z}_i^T \\
%%&=& \left[\textbf{z}_1, \cdots, \textbf{z}_n\right]  \left[\begin{array}{ccc}
%%\!\!\!w_1(\mbox{\boldmath $\beta$})\!\! & & \\
%%& \!\!\!\ddots\!\!\! & \\ & & \!\!\!w_n(\mbox{\boldmath $\beta$})\!\!\!\end{array}\right] 
%%\left[\begin{array}{c} \textbf{z}_1^T \\ \vdots \\ \textbf{z}_n^T\end{array}\right]
%%=Z^TW(\mbox{\boldmath $\beta$})Z,
%%\end{eqnarray*}
%% where $w_i(\mbox{\boldmath $\beta$}):=D_i^2(\mbox{\boldmath $\beta$})\sigma_i^{-2}(\mbox{\boldmath $\beta$})$ and 
%%
%%\[ W(\mbox{\boldmath $\beta$})=\left[\begin{array}{ccc}
%%\!\!\!w_1(\mbox{\boldmath $\beta$})\!\! & & \\
%%& \!\!\!\ddots\!\!\! & \\ & & \!\!\!w_n(\mbox{\boldmath $\beta$})\!\!\!\end{array}\right]=\mbox{diag}
%%\{w_1(\mbox{\boldmath $\beta$}), \cdots, w_n(\mbox{\boldmath $\beta$})\}\]
%%
%%\end{frame}
%%
%%\begin{frame}{\S2.2.1 Maximum likelihood estimation (5)}
%%\begin{itemize}
%%\item
%%The \textbf{observed information} matrix with respect to $\mbox{\boldmath $\beta$}$ is {\footnotesize
%%\begin{eqnarray*}
%%F_{\mbox{obs}}(\mbox{\boldmath $\beta$})& \!\!\!= \!\!\!& -\frac{\partial^2\ell(\mbox{\boldmath $\beta$})}{
%%\partial \mbox{\boldmath $\beta$}\partial \mbox{\boldmath $\beta$}^T}=
%%-\frac{\partial \textbf{s}(\mbox{\boldmath $\beta$})}{\partial \mbox{\boldmath $\beta$}^T} \\
%%& \!\!\!= \!\!\!& -\sum_{i=1}^n \textbf{z}_i \frac{\partial \{D_i(\mbox{\boldmath $\beta$})\sigma_i^{-2}(
%%\mbox{\boldmath $\beta$})\} }{\partial \mbox{\boldmath $\beta$}^T}[y_i-\mu_i(\mbox{\boldmath $\beta$})]
%%+\sum_{i=1}^n \textbf{z}_i D_i(\mbox{\boldmath $\beta$})\sigma_i^{-2}(\mbox{\boldmath $\beta$}) \cdot
%%\frac{\partial\mu_i}{\partial \eta_i}\cdot \frac{\partial\eta_i}{\partial \mbox{\boldmath $\beta$}^T} \\
%%& \!\!\!= \!\!\!& -\sum_{i=1}^n \textbf{z}_i \frac{\partial}{\partial \mbox{\boldmath $\beta$}^T}
%%\{D_i(\mbox{\boldmath $\beta$})\sigma_i^{-2}(\mbox{\boldmath $\beta$})\}\cdot
%%[y_i-\mu_i(\mbox{\boldmath $\beta$})] +\sum_{i=1}^n \textbf{z}_i D_i^2(\mbox{\boldmath $\beta$})
%%\sigma_i^{-2}(\mbox{\boldmath $\beta$}) \cdot \textbf{z}_i^T
%%\end{eqnarray*}}
%%\item
%%Hence $E\left[F_{\mbox{obs}}(\mbox{\boldmath $\beta$})\right]=\sum_{i=1}^n \textbf{z}_i 
%%D_i^2(\mbox{\boldmath $\beta$})\sigma_i^{-2}(\mbox{\boldmath $\beta$}) \textbf{z}_i^T
%%=F(\mbox{\boldmath $\beta$})=\mbox{cov}\left(\textbf{s}(\mbox{\boldmath $\beta$})\right)$.
%%\end{itemize}
%%\end{frame}
%%
%%\begin{frame}{\S2.2.1 Maximum likelihood estimation (6)}
%%\begin{itemize}
%%\item
%%When natural link is used, $\textbf{s}(\mbox{\boldmath $\beta$})$ and $F(\mbox{\boldmath $\beta$})$ can be
%%simplified to
%%\begin{eqnarray*}
%%\textbf{s}(\mbox{\boldmath $\beta$})&= & \frac{1}{\phi}\sum_{i=1}^n \omega_i\textbf{z}_i
%%[y_i-\mu_i(\mbox{\boldmath $\beta$})]=\frac{1}{\phi}Z^T\Omega \left[\textbf{y}-\mbox{\boldmath $\mu$}
%%(\mbox{\boldmath $\beta$})\right] \\
%%F(\mbox{\boldmath $\beta$}) &= & \frac{1}{\phi}\sum_{i=1}^n \omega_iv(\mu_i(\mbox{\boldmath $\beta$}))
%%\textbf{z}_i\textbf{z}_i^T=\frac{1}{\phi}Z^T\Omega V(\mbox{\boldmath $\beta$})Z
%%\end{eqnarray*}
%%where $\Omega=\mbox{diag}\{\omega_1,\cdots,\omega_n\}$ and $V(\mbox{\boldmath $\beta$})
%%=\mbox{diag}\{v(\mu_1),\cdots, v(\mu_n)\}$.
%%\item
%%For natural links, $\theta_i=\eta_i=\textbf{z}_i^T\mbox{\boldmath $\beta$} \quad \Rightarrow\quad
%%\mu_i=b^\prime (\theta_i)=b^\prime(\textbf{z}_i^T\mbox{\boldmath $\beta$})\equiv 
%%h(\textbf{z}_i^T\mbox{\boldmath $\beta$})$
%%
%%$\Rightarrow \quad D_i(\mbox{\boldmath $\beta$})=\frac{\partial h(\eta_i)}{\partial \eta_i}
%%=b^{\prime\prime}(\textbf{z}_i^T\mbox{\boldmath $\beta$})=v(h(\textbf{z}_i^T\mbox{\boldmath $\beta$}))$
%%
%%$\Rightarrow \quad D_i(\mbox{\boldmath $\beta$})\sigma_i^{-2}(\mbox{\boldmath $\beta$})
%%=v(h(\textbf{z}_i^T\mbox{\boldmath $\beta$}))\cdot\frac{\omega_i}{\phi
%%v(h(\textbf{z}_i^T\mbox{\boldmath $\beta$}))}=\frac{\omega_i}{\phi}$ and
%%
%%$\Rightarrow \quad w_i(\mbox{\boldmath $\beta$})=D_i^2(\mbox{\boldmath $\beta$})
%%\sigma_i^{-2}(\mbox{\boldmath $\beta$})=D_i(\mbox{\boldmath $\beta$})\cdot \frac{\omega_i}{\phi}
%%=\frac{1}{\phi}\omega_i v(h(\textbf{z}_i^T\mbox{\boldmath $\beta$}))$.
%%\end{itemize}
%%\end{frame}
%%
%%
%%\begin{frame}
%%We'll continue our discussion on MLE estimate for GLMs. 
%%
%%\end{frame}
%%
%%
%%\begin{frame}{\S2.2.1 Numerical computation of the MLE $\hat{\mbox{\boldmath $\beta$}}$ by IRLS}
%%\begin{itemize}
%%\item
%%MLE $\hat{\mbox{\boldmath $\beta$}}$ is the one maximizing $\ell(\mbox{\boldmath $\beta$})$, and often 
%%is obtained by solving the \textbf{estimating equation} $\textbf{s}(\mbox{\boldmath $\beta$})=0$, which uses 
%%the iterative \textbf{Fisher scoring method}: {\scriptsize
%%\begin{eqnarray*}
%%\hat{\mbox{\boldmath $\beta$}}^{(k+1)}&\!\!\!\!\!\!=\!\!\!\!\!\! & \hat{\mbox{\boldmath $\beta$}}^{(k)}+
%%F^{-1}(\hat{\mbox{\boldmath $\beta$}}^{(k)})\cdot \textbf{s}(\hat{\mbox{\boldmath $\beta$}}^{(k)}) \\
%%&\!\!\!\!\!\!=\!\!\!\!\!\! & \hat{\mbox{\boldmath $\beta$}}^{(k)}+ 
%%\left[Z^TW(\hat{\mbox{\boldmath $\beta$}}^{(k)})Z\right]^{-1}
%% Z^TD(\hat{\mbox{\boldmath $\beta$}}^{(k)})\Sigma^{-1}(\hat{\mbox{\boldmath $\beta$}}^{(k)})
%%\left[\textbf{y}-\mbox{\boldmath $\mu$}(\hat{\mbox{\boldmath $\beta$}}^{(k)})\right] \\
%%&\!\!\!\!\!\!=\!\!\!\!\!\! & \hat{\mbox{\boldmath $\beta$}}^{(k)}+ 
%%\left[Z^TW(\hat{\mbox{\boldmath $\beta$}}^{(k)})Z\right]^{-1}
%% Z^TD^2(\hat{\mbox{\boldmath $\beta$}}^{(k)})\Sigma^{-1}(\hat{\mbox{\boldmath $\beta$}}^{(k)})
%%D^{-1}(\hat{\mbox{\boldmath $\beta$}}^{(k)}) 
%%\left[\textbf{y}-\mbox{\boldmath $\mu$}(\hat{\mbox{\boldmath $\beta$}}^{(k)})\right] \\
%%&\!\!\!\!\!\!=\!\!\!\!\!\! & \left[Z^TW(\hat{\mbox{\boldmath $\beta$}}^{(k)})Z\right]^{-1}\left\{
%%Z^TW(\hat{\mbox{\boldmath $\beta$}}^{(k)})Z \hat{\mbox{\boldmath $\beta$}}^{(k)}
%%+Z^TW(\hat{\mbox{\boldmath $\beta$}}^{(k)}) D^{-1}(\hat{\mbox{\boldmath $\beta$}}^{(k)})
%%\left[\textbf{y}-\mbox{\boldmath $\mu$}(\hat{\mbox{\boldmath $\beta$}}^{(k)})\right]\!\right\} \\
%%&\!\!\!\!\!\!=\!\!\!\!\!\!& \left[Z^TW(\hat{\mbox{\boldmath $\beta$}}^{(k)})Z\right]^{-1}\left\{
%%Z^TW(\hat{\mbox{\boldmath $\beta$}}^{(k)})\left[Z \hat{\mbox{\boldmath $\beta$}}^{(k)}
%%+ D^{-1}(\hat{\mbox{\boldmath $\beta$}}^{(k)})
%%\left[\textbf{y}-\mbox{\boldmath $\mu$}(\hat{\mbox{\boldmath $\beta$}}^{(k)})\right]\right]\right\} \\
%%&\!\!\!\!\!\!\stackrel{denoted}{=}\!\!\!\!\!\! & \left[Z^TW(\hat{\mbox{\boldmath $\beta$}}^{(k)})Z\right]^{-1}
%%Z^TW(\hat{\mbox{\boldmath $\beta$}}^{(k)})\cdot \tilde{\textbf{y}}^{(k)}
%%\end{eqnarray*}}
%%where $\tilde{\textbf{y}}^{(k)}=(\tilde{y}_1^{(k)},\cdots, \tilde{y}_n^{(k)})^T
%%=Z \hat{\mbox{\boldmath $\beta$}}^{(k)}+ D^{-1}(\hat{\mbox{\boldmath $\beta$}}^{(k)})
%%\left[\textbf{y}-\mbox{\boldmath $\mu$}(\hat{\mbox{\boldmath $\beta$}}^{(k)})\right]$
%%is the \textbf{working observation vector} at iteration $k$, $k=0,1,2,\cdots$.
%%\end{itemize}
%%\end{frame}
%%
%%
%%\begin{frame}{\S2.2.1 Implementing Fisher scoring for computing $\hat{\mbox{\boldmath $\beta$}}$}
%%\begin{enumerate}
%%\item
%%Start with an initial estimate $\hat{\mbox{\boldmath $\beta$}}^{(0)}$, often chosen as the unweighted
%%least squares (LS) estimate from $(g(y_i),\textbf{z}_i)$, $i=1,\cdots,n$, where $g(\cdot)$ is the link (modified if $g(y_i)$ is undefined, such as when $yi = 0$ in Poisson regression).
%%\item
%%For $k=0,1,2,\cdots$ compute 
%%$$\hat{\mbox{\boldmath $\beta$}}^{(k+1)}=\left[Z^TW(\hat{\mbox{\boldmath $\beta$}}^{(k)})Z\right]^{-1}
%%Z^TW(\hat{\mbox{\boldmath $\beta$}}^{(k)})\cdot \tilde{\textbf{y}}^{(k)}$$
%%until $\frac{\|\hat{\boldsymbol \beta}^{(k+1 )} - \hat{\boldsymbol \beta}^{(k)}\|}{\|\hat{\boldsymbol \beta}^{(k)}\|} \leq \epsilon $ for some prechosen small $\epsilon > 0$ (``convergence").
%%
%%\item It is hence called a  \textbf{iteratively reweighted least squares (IRLS) method}. 
%%\end{enumerate}
%%%Fisher scoring method in GLM estimation is also called the \textbf{iteratively reweighted least squares (IRLS) method}.
%%\end{frame}
%%
%%
%%\begin{frame}
%%
%%\begin{itemize}
%%
%%\item 
%%Each component of the working observation vector is simply a linearization of $g(y_i)$ at $\hat{\boldsymbol \beta}^{(k)}$
%%
%%\begin{align*}
%%g(y_i) &\approx g(\mu) +g'(\mu) (y_i - \mu) \\
%%&= {\bf z}_i^T{\boldsymbol \beta} +  \left(\frac{\partial h}{ \partial \eta} ({\bf z}_i^T{\boldsymbol \beta} )\right)^{-1}  (y_i - \mu_i ({\boldsymbol \beta}))
%%\end{align*}
%%(Generalized Linear Models, 2nd ed, McCullagh and Nelder, p.40)
%%
%%\item If use $F_{obs} (\cdot)$ instead of $F (\cdot)$, we get the standard Newton-Raphson procedure. ( These two procedures coincide if natural link function is used)
%%\end{itemize}
%%\end{frame}
%%
%%\begin{frame}{\S2.2.1 Asymptotic properties associated with the MLE $\hat{\mbox{\boldmath $\beta$}}$}
%%Under suitable ``regularity" assumptions, we have these nice properties:
%%\begin{enumerate}
%%\item
%%\textit{Asymptotic existence and uniqueless}. The MLE $\hat{\mbox{\boldmath $\beta$}}$ exists
%%and (locally) unique w/ probability tending to 1. 
%%\item
%%\textit{Consistency}. If $\mbox{\boldmath $\beta$}$ is the ``true'' value, then 
%%$\hat{\mbox{\boldmath $\beta$}}\rightarrow \mbox{\boldmath $\beta$}$ in probability 1 (\textit{weak consistency}),
%%or $\hat{\mbox{\boldmath $\beta$}}\rightarrow \mbox{\boldmath $\beta$}$ with probability 1 (or almost surely,
%%\textit{strong consistency}).
%%
%%%\item
%%%Estimating $\phi$: 
%%%$\displaystyle \hat{\phi}=\frac{1}{g-p}\sum_{i=1}^g\frac{(\bar{y}_i-\hat{\mu}_i)^2}{v(\hat{\mu}_i)/n_i}$
%%%is a consistent estimator of $\phi$. (This is true for grouped data only.)
%%\end{enumerate}
%%\end{frame}
%%
%%
%%\begin{frame}
%%{\S2.2.1 Asymptotic properties associated with the MLE $\hat{\mbox{\boldmath $\beta$}}$}
%%
%%\begin{enumerate}[3]
%%\item
%%\textit{Asymptotic normality}. $\hat{\mbox{\boldmath $\beta$}}\stackrel{a}{\sim} \mbox{Normal}
%%\left(\mbox{\boldmath $\beta$}, F^{-1}(\hat{\mbox{\boldmath $\beta$}})\right)$, equivalently
%%$$F^{1/2}(\hat{\mbox{\boldmath $\beta$}})\left(\hat{\mbox{\boldmath $\beta$}}-
%%\mbox{\boldmath $\beta$}\right)\stackrel{d}{\rightarrow} N(0,1) \quad\mbox{as $n\rightarrow\infty$}.$$
%%This implies $\mbox{cov}(\hat{\mbox{\boldmath $\beta$}})\stackrel{a}{=}A(\hat{\mbox{\boldmath $\beta$}})
%%= F^{-1}(\hat{\mbox{\boldmath $\beta$}})$. (The proof uses the property $\textbf{s}(\mbox{\boldmath $\beta$})
%%\stackrel{a}{\sim} \mbox{Normal}\left(0, F(\mbox{\boldmath $\beta$})\right)$.)
%%\end{enumerate}
%% \ \
%% 
%%(Now recall in the linear model case, the Fisher information is $\sigma^{-2} {\bf Z}^T {\bf Z}$)
%%
%%\end{frame}
%%
%%\begin{frame} {\S2.2.1 Estimation of $\phi$}
%%\begin{itemize}
%%\item We also note that $\phi$ only appears as a factor in $W$, and falls out of the expression
%%
%% \[\hat{\mbox{\boldmath $\beta$}}^{(k+1)}=\left[Z^TW(\hat{\mbox{\boldmath $\beta$}}^{(k)})Z\right]^{-1}
%%Z^TW(\hat{\mbox{\boldmath $\beta$}}^{(k)})\cdot \tilde{\textbf{y}}^{(k)}\]
%%\item But we need $\phi$ for computing the Fisher info for the limiting covariance of the MLE. 
%%
%%\item\[\hat{\phi}=\frac{1}{g-p}\sum_{i=1}^g\frac{(\bar{y}_i-\hat{\mu}_i)^2}{v(\hat{\mu}_i)/n_i}\]
%%is a consistent estimator of $\phi$. (This is true for grouped data only.)
%%\end{itemize}
%%\end{frame}
%%
%%\subsection{\S2.2.2 Linear hypothesis testing in GLM}
%%
%%\begin{frame}{\S2.2.2 Testing linear hypotheses}
%%\begin{itemize}
%%\item
%%Most testing problems concerning $\mbox{\boldmath $\beta$}$ have the form of \textit{linear hypothesis}:
%%$$H_0: C\mbox{\boldmath $\beta$}=\mbox{\boldmath $\xi$} \quad\mbox{vs.}\quad
%%H_1: C\mbox{\boldmath $\beta$}\neq\mbox{\boldmath $\xi$}$$
%%where $C$ is a known matrix of full row rank $s\leq p$, and $\mbox{\boldmath $\xi$}$ is a known vector.
%%\item
%%A special case is 
%%$$H_0: \mbox{\boldmath $\beta$}_r=0 \quad\mbox{vs.}\quad
%%H_1: \mbox{\boldmath $\beta$}_r\neq 0$$
%%where $\mbox{\boldmath $\beta$}_r$ is a sub-vector of $\mbox{\boldmath $\beta$}$. This is
%%equivalent to comparing the full model with one of its sub-models.
%%\end{itemize}
%%\end{frame}
%%
%%\begin{frame}{\S2.2.2 Likelihood ratio (LR), Wald and score tests (1)}
%%\begin{itemize}
%%\item
%%Three types of tests are available for testing $H_0: C\mbox{\boldmath $\beta$}=\mbox{\boldmath $\xi$}
%%\;\mbox{vs.}\; H_1: C\mbox{\boldmath $\beta$}\neq\mbox{\boldmath $\xi$}$:
%%
%%\begin{center}
%%\textbf{Likelihood ratio (LR) test}, \textbf{Wald test}, and \textbf{score test}.
%%\end{center}
%%\item
%%A special case is 
%%$$H_0: \mbox{\boldmath $\beta$}_r=0 \quad\mbox{vs.}\quad
%%H_1: \mbox{\boldmath $\beta$}_r\neq 0$$
%%where $\mbox{\boldmath $\beta$}_r$ is a sub-vector of $\mbox{\boldmath $\beta$}$. This is
%%equivalent to comparing the full model with one of its sub-models.
%%\end{itemize}
%%\end{frame}
%%
%%\begin{frame}{\S2.2.2 Likelihood ratio (LR), Wald and score tests (2)}
%%\begin{enumerate}
%%\item
%%\textbf{LR test statistic}:
%% \[\lambda=-2\left\{\ell(\tilde{\mbox{\boldmath $\beta$}}) -\ell(\hat{
%%\mbox{\boldmath $\beta$}})\right\},\] 
%%where $\tilde{\mbox{\boldmath $\beta$}}$ is the MLE of
%%$\mbox{\boldmath $\beta$}$ under $H_0$, and $\hat{\mbox{\boldmath $\beta$}}$ is the MLE of
%%$\mbox{\boldmath $\beta$}$ under $H_1$; also $\hat{\mbox{\boldmath $\beta$}}$ is the unrestricted MLE of
%%$\mbox{\boldmath $\beta$}$. 
%%
%%%
%%%$\lambda$ needs to be divided by $\hat{\phi}$ for over-dispersion models.
%%\item
%%\textbf{Wald test statistic}: \[W=\left(C\hat{\mbox{\boldmath $\beta$}}-\mbox{\boldmath $\xi$}\right)^T
%%\left[CF^{-1}(\hat{\mbox{\boldmath $\beta$}})C^T \right]^{-1}
%%\left(C\hat{\mbox{\boldmath $\beta$}}-\mbox{\boldmath $\xi$}\right),\] where $\hat{\mbox{\boldmath $\beta$}}$
%%is the unrestricted MLE of $\mbox{\boldmath $\beta$}$.
%%\item
%%\textbf{Score test statistic}: $U=\textbf{s}^T(\tilde{\mbox{\boldmath $\beta$}})F^{-1}(\tilde{\mbox{\boldmath $\beta$}})
%%\textbf{s}(\tilde{\mbox{\boldmath $\beta$}})$, where $\tilde{\mbox{\boldmath $\beta$}}$
%%is the MLE of $\mbox{\boldmath $\beta$}$ under $H_0$.
%%\end{enumerate}
%%%\begin{itemize}
%%%
%%%
%%%\end{itemize}
%%\end{frame}
%%
%%\begin{frame}{\S2.2.2 Likelihood ratio (LR), Wald and score tests (3)}
%%\begin{itemize}
%%\item  The  exponential model is assumed for the LR test; not defined for over-dispersion model such as the quasi likelihood to be covered. 
%%
%%\item 
%%Wald and score test statistics are properly defined for over-dispersion models, since they are based on the score ${\bf s}$ (More on this later again)
%%
%%\item Wald and score are quadratic approximation of the LR test (see blackboard derivation). 
%%
%%\item Wald and Score are computationally attractive; both eliminate the need to compute one of $\tilde{\boldsymbol \beta}$ or $\hat{\boldsymbol \beta}$. Wald is good when $H_1$ has been fitted, score is good when $H_0$ has been fitted. 
%%
%%
%%\end{itemize}
%%
%%\end{frame}
%%
%%\begin{frame} {\S2.2.2 Likelihood ratio (LR), Wald and score tests (4)}
%%\begin{itemize}
%%\item
%%Under \[H_0: C{\boldmath \beta}={\boldmath \xi} \text{ with } \text{rank}(C)=s, \] $\lambda, W$ and
%%$U$ all follow a $\chi^2(s)$ distribution asymptotically under `regularity conditions'.
%%%\item
%%%$\lambda/\hat{\phi}$, $W$ and $U$ can be used for \textbf{analysis of deviance} (i.e. comparing nested
%%%models) as a special case of testing linear hypotheses.
%%
%%\end{itemize}
%%
%%
%%\end{frame}
%
%\subsection{\S2.2.3 Deviance, residuals and goodness-of-fit measure in GLM} \label{sec:2.2.3}
%
%\begin{frame} {Model adequacy measure in GLM}
%\begin{itemize}
%\item First of all, testing model fit only makes sense when over-dispersion is not assumed, since we don't assume the exponential model form is true anyway. Again a topic of quasi-likelihood covered later.   
%
%\item In the sequel, ``goodness-of-fit" and ``model adequacy" mean the same thing.
%\end{itemize}
%\end{frame}
%
%\begin{frame}{\S2.2.3 Model adequacy measure in GLM }
%\begin{itemize}
%\item
%In linear model, goodness-of-fit  is measured by the sum of squared residuals (or errors)
%$\mbox{SSE}\!=\!\sum_{i=1}^n(y_i-\hat{\mu}_i)^2$. SSE is no longer a reasonable measure if $\mbox{Var}(y)$ 
%varies with $\mu$.
%\item
%A generalisation of SSE to GLM  is the \textbf{Pearson $\chi^2$ statistic}, taking the form
%$$\chi^2=\sum_{i=1}^g \frac{(\bar{y}_i-\hat{\mu}_i)^2}{v(\hat{\mu}_i)/n_i}$$
%where $\bar{y}_i$ is the group mean, $v(\hat{\mu}_i)$ is the estimated variance function, and $n_i$ is the size of group $i$.
%\end{itemize}
%\end{frame}
%
%
%
%\begin{frame}{\S2.2.3 Model adequacy measure in GLM }
%\begin{itemize}
%%\item
%%Pearson $\chi^2$ statistic is often an appropriate model adequacy measure (or discrepancy measure, or goodness-of-fit
%%measure).
%\item
%Asymptotically (i.e. all $n_i$'s are large),
% \[\chi^2/\phi \sim_a \chi^2(g-p)\] if the GLM model is true. $p = $  number of coefficient parameters in the GLM.
%\item
%If $n$ is large but $n_i$'s are small (in particular when $n_i=1$), it is dangerous to use the Pearson
%$\chi^2$ statistics to test model adequacy.
%
%\end{itemize}
%\end{frame}
%
%\begin{frame}{\S2.2.3 Model adequacy measure in GLM}
%\begin{itemize}
%\item
%Another model adequacy measure is the {\bf deviance}. 
%\item
%It is twice the difference between the maximum log-likelihood achievable in a saturated model and that in
%the current model, i.e.
%$$D^*(\textbf{y}_n;\mbox{\boldmath $\hat{\mu}$})=2\left[\ln L(\mbox{\boldmath $\theta$}^*,\phi\mid\textbf{y}_n)
%-\ln L(\mbox{\boldmath $\hat{\theta}$},\phi\mid\textbf{y}_n)\right]$$
%which is called the \textbf{scaled deviance}. Just a special case the likelihood ratio statistic. 
%
%% Here {\small \vspace{-3pt}
%%\begin{equation*}
%%L(\mbox{\boldmath $\theta$},\phi\mid\textbf{y}_n)=\exp\left\{\sum_{i=1}^n\frac{y_i\theta_i-b(\theta_i)}{\phi}\cdot
%%\omega_i +\sum_{i=1}^nc(y_i,\phi, \omega_i)\right\} \vspace{-3pt}
%%\label{eq-8.4}
%%\end{equation*}}
%%is the likelihood function for the data $\textbf{y}_n=(y_1,\cdots,y_n)^T$.
%\item Remark: we are having grouped data in mind here
%
%\item
%%$\mbox{\boldmath $\theta$}\!=\!(\theta_1,\cdots,\theta_n)^T$, $\mbox{\boldmath $\mu$}\!=\!(\mu_1,\cdots,\mu_n)^T
%%\!=\!(b^\prime(\theta_1),\cdots, b^\prime(\theta_n))^T$,
%
% $\mbox{\boldmath $\hat{\theta}$}$ is the MLE of
%$\mbox{\boldmath $\theta$}$ under the GLM
%\item $\mbox{\boldmath $\theta$}^*=(\theta_1^*,\cdots, \theta_g^*)^T$ is such that 
%\[
%(\bar{y}_1,\cdots,\bar{y}_g)^T
%=(b^\prime(\theta_1^*),\cdots,b^\prime(\theta_g^*))^T = ({\mu}^*_1, \dots, {\mu}^*_g).\]
%\end{itemize}
%\end{frame}
%
%\begin{frame}{\S2.2.3 Deviance measure in GLM (2)}
%\begin{itemize}
%\item
%We call $\mbox{\boldmath $\theta$}^*$ the MLE of $\mbox{\boldmath $\theta$}$ in the \textbf{saturated model}, a full model with $g$ parameters (as many parameters as there are observsations).
%\item
%In a saturated model, from the likelihood expression we see that in this case, the MLE of 
%$\theta_i$ needs to satisfy $\bar{y}_i=b^\prime(\theta_i^*) = {\mu}^*_i$.
%\item
%In GLM, $\theta_1,\cdots,\theta_g$ are determined by the regression parameters $\beta_1,\cdots,\beta_p$
%and the independent variables $x_1,\cdots,x_p$. So the MLE of $\beta_1,\cdots,\beta_p$ are to be obtained first, 
%the MLE $\mbox{\boldmath $\hat{\theta}$}$ are then obtained.
%\item
%%The saturated model provides the best fit to the data because it has the same number of independent parameters
%%and the observations.
%Maximum likelihood from the saturated model must be at least as large as that from a usual GLM.
%Therefore, the scaled deviance $D^*(\textbf{y}_n;\mbox{\boldmath $\hat{\mu}$})\geq 0$. 
%%and measures
%%the adequacy of fit of a GLM.
%\end{itemize}
%\end{frame}
%
%\begin{frame}{\S2.2.3 Deviance measure in GLM (3)}
%\begin{itemize}
%\item
%%Consider a random sample $\textbf{y}_n=(y_1,\cdots,y_n)^T$ and a GLM as being formulated so far.
%By convention
%$$D(\textbf{y}_n;\mbox{\boldmath $\hat{\mu}$}):=\phi D^*(\textbf{y}_n;\mbox{\boldmath $\hat{\mu}$})$$
%%and call $D(\textbf{y}_n;\mbox{\boldmath $\hat{\mu}$})$ 
%as the \textbf{deviance} of the GLM under consideration.
%\item
%Asymptotically, $$\phi^{-1}D(\textbf{y}_n;\mbox{\boldmath $\hat{\mu}$})\sim_a\chi^2(g-p)$$ 
%asymptotically when the GLM  for $Y$ is correct  + when group sizes $n_i$ are all large enough + ``regularity conditions". 
%\item
%But when $n_i$ is small (e.g. $n_i = 1$), use of the chi-square limit is dangerous. 
%\end{itemize}
%\end{frame}
%
%\begin{frame}{\S2.2.3 Deviances for several important distributions}
%\textbf{Deviances for several important distributions in GLM.}
%\begin{eqnarray*}
%\mbox{Normal}\!\!\! & N(\mu,\sigma^2): & D(\textbf{y}_n;\mbox{\boldmath $\hat{\mu}$})=
%\sum_{i=1}^n (y_i-\hat{\mu}_i)^2 \\
%\mbox{Poisson}\!\!\! &\mbox{Poi}(\mu): &  D(\textbf{y}_n;\mbox{\boldmath $\hat{\mu}$})=2\sum_{i=1}^n
%\left\{y_i\ln\frac{y_i}{\hat{\mu}_i}-(y_i-\hat{\mu}_i)\right\} \\
%\mbox{binomial}\!\!\! & \mbox{Bin}(m,\mu): & D(\textbf{y}_n;\mbox{\boldmath $\hat{\mu}$})=2\sum_{i=1}^n
%\left\{y_i\ln\frac{y_i}{\hat{\mu}_i}\!+\!(\!m_i\!-\!y_i\!)\ln\frac{m_i\!-\!y_i}{m_i\!-\!\hat{\mu}_i}\right\} \\
%\mbox{gamma} \!\!\!& \mbox{Gam}(\mu,\nu): & D(\textbf{y}_n;\mbox{\boldmath $\hat{\mu}$})=2\sum_{i=1}^n
%\left\{-\ln\frac{y_i}{\hat{\mu}_i}+\frac{y_i-\hat{\mu}_i}{\hat{\mu}_i}\right\}
%\end{eqnarray*}
%\end{frame}
%
%\begin{frame}{\S2.2.3 Analysis of deviance (1)}
%\begin{itemize}
%\item
%An analysis of variance (ANOVA) table can be constructed for each linear regression model as we have seen
%in a linear model course. 
%\item
%Similar to this, an \textbf{analysis of deviance (ANODEV)} table can be constructed for each GLM.
%\item
%Generally speaking, tests performed on ANOVA and ANODEV tables are likelihood ratio tests.
%\item
%However, there are some differences in interpreting ANOVA and ADODEV tables.
%\item
%Consider a $3\times 3$ factorial design, of two treatment factors A and B, with 2 replicates at each combination
%level of A and B. So there will be in total 18 observations for the response variable $y$. ($3\times 3$ factorial design means that each of treatment $A$ or $B$ has 3 levels)
%\end{itemize}
%\end{frame}
%
%\begin{frame}{\S2.2.3 Analysis of deviance (2)}
%If it is appropriate to fit the 18 $y$ observations by linear regression models, an ANOVA table 
%would typically look like the following:
%
%\begin{center} Analysis of Variance (ANOVA)\end{center}
%
%{\small
%\begin{tabular}{|cccccc|}\hline
%Model & df & SSE & $\Delta$SS & $\Delta$df & Effect Source \\ \hline
%1 & 17 & 1000 & - & -  & - \\ \hline
%A & 15 & 400 & $600=1000-400$ & $2=17-15$ & A \\ \hline
%A+B & 13 & 300 & $100=400-300$ & $2=15-13$ & B \\ \hline
%A+B+A:B & 9 & 100 & $200=300-100$ & $4=13-9$ & A:B \\ \hline 
%\end{tabular}}
%\begin{itemize}
%\item
%$\Delta$SS is the reduction of SSE as an additional term is added to the model.
%$\Delta\mbox{SS}\stackrel{d}{=}\chi^2(\Delta\mbox{df})$. Effects are assessed by
%$F$ tests.
%\item
%$\Delta$df is the reduction of the model degrees of freedom.
%\item
%The effects of A, B and A:B on $Y$ are orthogonal.
%\end{itemize}
%\end{frame}
%
%\begin{frame}{\S2.2.3 Analysis of deviance (3)}
%If it is more appropriate to fit the 18 $y$ observations by GLM, an ANODEV table would typically look like the following:
%
%\begin{center} Analysis of Deviance (ANODEV)\end{center}
%
%{\footnotesize
%\begin{tabular}{|cccccc|}\hline
%Model & df & Deviance & $\Delta D$ & $\Delta$df & Effect Source \\ \hline
%1 & 17 & 1200 & - & - & - \\ \hline
%A & 15 & 800 & $400\!=\!1200\!-\!800$ & $2\!=\!17\!-\!15$ & A ignoring B \\ \hline
%A+B & 13 & 500 & $300\!=\!800\!-\!500$ & $2\!=\!15\!-\!13$ & B eliminating A\\ \hline
%A+B+A:B & 9 & 200 & $300\!=\!500\!-\!200$ & $4\!=\!13\!-\!9$ & A:B eliminating \\ 
%& & & &  & A \& B \\ \hline 
%\end{tabular}}
%\begin{itemize}
%\item
%$\Delta D=D_0-D_1$ is the reduction of deviance between two nested models $M_0\subseteq M_1$.
%$\Delta D\stackrel{d}{\approx}\chi^2(\Delta\mbox{df})$ if $M_0$ is adequate.
%\item
%$\Delta$df is the reduction of the model degrees of freedom.
%\item
%Effects of A, B and A:B on $Y$ are not orthogonal, depend on their positions in the model, and
%are assessed by $\chi^2$ or $F$ tests.
%\end{itemize}
%\end{frame}
%
%%\section{\S8.4 Residuals}
%%\label{sec:residuals}
%
%\begin{frame}{\S2.2.3 GLM residuals (1)}
%\begin{itemize}
%\item
%In a linear regression model $y=E(y)+\varepsilon =\mu+\varepsilon=\sum_{j=1}^p x_j\beta_j+\varepsilon$,
%the residual is defined as $r_i=y_i-\hat{\mu}_i$ for an observation $y_i$, with $\hat{\mu}_i=\hat{y}_i$ being the 
%fitted value of $y_i$.
%The residual $r_i$ measures the goodness of fit to $y_i$ by the model if $\mbox{Var}(y)$ is constant.
%\item
%In GLM the residual of $y_i$ has to be defined differently to accommodate the cases
%where the variance function $v(\mu)$ is no longer constant.
%\item
%At least four types of residual are defined for GLM: \textbf{Pearson}, \textbf{deviance}, \textbf{working}, and
%\textbf{response}. All these residuals can be computed using the \texttt{R} function \texttt{residuals.glm()}.
%\end{itemize}
%\end{frame}
%
%\begin{frame}{\S2.2.3 GLM residuals (2)}
%\begin{itemize}
%\item
%\textbf{Pearson residual}: $r_{_{Pi}}=\frac{\bar{y}_i-\hat{\mu}_i}{\sqrt{v(\hat{\mu})/n_i}}$.
%
%Therefore $\sum_{i=1}^g r^2_{_{Pi}}=\sum_{i=1}^g \frac{(y_i-\hat{\mu}_i)^2}{v(\hat{\mu}_i)/n_i}=\chi^2$, the
%Pearson $\chi^2$ statistic. 
%\item
%\textbf{Deviance residual}: $r_{_{Di}}=\mbox{sign}(y_i-\hat{\mu}_i)\sqrt{d_i}$, \\ where $d_i$'s are such that
%$\displaystyle \sum_{i=1}^n d_i=D(\textbf{y}_n;\mbox{\boldmath $\hat{\mu}$})$. In particular, deviance residuals for
%the following distributions are:
%{\footnotesize
%\begin{eqnarray*}
%\mbox{Normal}\!\!\! & N(\mu,\sigma^2): & r_{_{Di}}\!=\!\mbox{sign}(y_i-\hat{\mu}_i)
%|y_i-\hat{\mu}_i|\!=\!y_i-\hat{\mu}_i \\
%\mbox{Poisson}\!\!\! &\mbox{Poi}(\mu): &  r_{_{Di}}\!=\!\mbox{sign}(y_i-\hat{\mu}_i)\sqrt{2\left\{y_i
%\ln\frac{y_i}{\hat{\mu}_i}-(y_i-\hat{\mu}_i)\right\}} \\
%\mbox{binomial}\!\!\! & \mbox{Bin}(m,\mu): & r_{_{Di}}\!=\!\mbox{sign}(y_i-\hat{\mu}_i)\sqrt{2
%\left\{y_i\ln\frac{y_i}{\hat{\mu}_i}\!+\!(\!m_i\!-\!y_i\!)\ln\frac{m_i\!-\!y_i}{m_i\!-\!\hat{\mu}_i}\right\}} \\
%\mbox{gamma} \!\!\!& \mbox{Gam}(\mu,\nu): & r_{_{Di}}\!=\!\mbox{sign}(y_i-\hat{\mu}_i)\sqrt{2
%\left\{-\ln\frac{y_i}{\hat{\mu}_i}+\frac{y_i-\hat{\mu}_i}{\hat{\mu}_i}\right\}}
%\end{eqnarray*}}
%\end{itemize}
%\end{frame}
%
%\begin{frame}{\S2.2.3 GLM residuals (3)}
%\begin{itemize}
%\item
%\textbf{Working residual}: $r_{_{Wi}}=(y_i-\hat{\mu}_i)\frac{d\eta_i}{d\mu_i}\mid_{\mu_i=\hat{\mu}_i}
%\approx g(y_i)-g(\hat{\mu}_i)$, \\
%where $g(\mu_i)$ is the link function and $\eta_i=\textbf{z}_i^T\mbox{\boldmath $\beta$}$ is the linear predictor in GLM.
%
%The working residual is used during computing the MLE of $\mbox{\boldmath $\beta$}$.
%\item
%\textbf{response residual}: $r_{_{Ri}}=y_i-\hat{\mu}_i$, \\
%which is not a right measure except when the GLM reduces to a linear model.
%\end{itemize}
%\end{frame}
%
%\begin{frame}[fragile] \frametitle{\S2.2.3 Example: Caesarian birth study (1)}
%For the data of Caesarian birth study we combine categories I and II of the response variable into
%one category \texttt{Infection}. Then fit logistic models to the reorganized data where the new response
%variable has only two categories \texttt{Infection} and \texttt{Non-infection}.
%
%\begin{scriptsize}
%\begin{Schunk}
%\begin{Sinput}
%> Caes.dat <- read.csv("D:/MAST90084/Example2-Caes.csv")
%> is.data.frame(Caes.dat)
%\end{Sinput} 
%\begin{Soutput}
%[1] TRUE
%\end{Soutput}
%\begin{Sinput}
%> Caes.dat
%\end{Sinput} 
%\begin{Soutput}
%  Infection total Antib RiskF Plan
%1         1    18     1     1    1
%2         0     2     1     0    1
%3        28    58     0     1    1
%4         8    40     0     0    1
%5        11    98     1     1    0
%6         0     0     1     0    0
%7        23    26     0     1    0
%8         0     9     0     0    0
%\end{Soutput}
%\end{Schunk}
%\end{scriptsize}
%\end{frame}
%
%\begin{frame}[fragile] \frametitle{\S2.2.3 Example: Caesarian birth study (2)}
%\begin{scriptsize}
%\begin{Schunk}
%\begin{Sinput}
%> Caes1=glm(Infection/total~Antib+RiskF+Plan,family=binomial, weight=total, data=Caes.dat)
%> summary(Caes1)
%\end{Sinput} 
%\begin{Soutput}
%Deviance Residuals: 
%       1         2         3         4         5         7         8  
% 0.26470  -0.15231  -0.78520   1.21563  -0.07162   1.49623  -2.56229  
%
%Coefficients:
%            Estimate Std. Error z value Pr(>|z|)    
%(Intercept)  -0.8207     0.4947  -1.659   0.0971 .  
%Antib        -3.2544     0.4813  -6.761 1.37e-11 ***
%RiskF         2.0299     0.4553   4.459 8.25e-06 ***
%Plan         -1.0720     0.4254  -2.520   0.0117 *  
%---
%Signif. codes:  0 '***' 0.001 '**' 0.01 '*' 0.05 '.' 0.1  ' ' 1
%
%(Dispersion parameter for binomial family taken to be 1)
%
%    Null deviance: 83.491  on 6  degrees of freedom
%Residual deviance: 10.997  on 3  degrees of freedom
%  (1 observation deleted due to missingness)
%AIC: 36.178
%Number of Fisher Scoring iterations: 5
%\end{Soutput}
%\end{Schunk}
%\end{scriptsize}
%\end{frame}
%
%\begin{frame}[fragile] \frametitle{\S2.2.3 Example: Caesarian birth study (3)}
%\begin{scriptsize}
%\begin{Schunk}
%\begin{Sinput}
%> anova(Caes1, test="Chi")
%\end{Sinput} 
%\begin{Soutput}
%Analysis of Deviance Table; Model: binomial, link: logit; Response: Infection/total
%Terms added sequentially (first to last)
%
%      Df Deviance Resid. Df Resid. Dev  Pr(>Chi)    
%NULL                      6     83.491              
%Antib  1   38.736         5     44.756 4.852e-10 ***
%RiskF  1   26.881         4     17.874 2.164e-07 ***
%Plan   1    6.878         3     10.997  0.008728 ** 
%\end{Soutput}
%\begin{Sinput}
%> resid(Caes1,type="deviance")
%\end{Sinput}
%\begin{Soutput}
%       1        2        3        4        5        7        8 
% 0.26470 -0.15231 -0.78520  1.21563 -0.07162  1.49623 -2.56229
%\end{Soutput}
%\begin{Sinput}
%> resid(Caes1,type="pearson")
%\end{Sinput}
%\begin{Soutput}
%       1        2        3        4        5        7        8 
% 0.27684 -0.10786 -0.78635  1.29465 -0.07141  1.38710 -1.99030
%\end{Soutput}
%\begin{Sinput}
%> sum(resid(Caes1,type="deviance")^2)    #[1] 10.99674
%> sum(resid(Caes1,type="pearson")^2)/(7-4)    #[1] 2.757724
%> 1-pchisq(sum(resid(Caes1,type="pearson")^2),3)   #[1] 0.04069089
%> 1-pchisq(sum(resid(Caes1,type="deviance")^2),3)   #[1] 0.01174351
%\end{Sinput}
%\end{Schunk}
%\end{scriptsize}
%\end{frame}
%
%\begin{frame}[fragile] \frametitle{\S2.2.3 Example: Caesarian birth study (4)}
%\begin{scriptsize}
%\begin{Schunk}
%\begin{Sinput}
%Caes2=glm(Infection/total~Antib+RiskF+Plan+RiskF:Plan,family=binomial, 
%                         weight=total, data=Caes.dat)
%summary(Caes2);  anova(Caes2,test="Chi")
%sum(resid(Caes2,type="pearson")^2)/(7-5)
%1-pchisq(sum(resid(Caes2,type="pearson")^2),2)
%1-pchisq(sum(resid(Caes2,type="deviance")^2),2)
%
%Caes3=glm(Infection/total~Antib+RiskF+Plan+RiskF:Plan,
%            family=binomial(link="probit"), weight=total, data=Caes.dat)
%summary(Caes3); anova(Caes3, test="Chi")
%
%Caes4=glm(Infection/total~Antib+RiskF+Plan+RiskF:Plan,
%           family=binomial(link="probit"), weight=total, data=Caes.dat[-6,])
%summary(Caes4); anova(Caes3, test="Chi")
%
%Caes5=glm(Infection/total~Antib+RiskF+Plan+RiskF:Plan,family=quasibinomial, 
%              weight=total, data=Caes.dat[-6,])
%summary(Caes5); anova(Caes5, test="F")
%
%sum(resid(Caes5,type="pearson")^2)/(7-5)
%
%Caes6=glm(Infection/total~Antib+RiskF+Plan+RiskF:Plan,family=quasibinomial, 
%           weight=total, data=Caes.dat, subset=c(1,2,3,4,5,7,8))
%summary(Caes6);  anova(Caes6, test="F")
%\end{Sinput} 
%\end{Schunk}
%\end{scriptsize}
%\end{frame}

\section{\S2.3 Some extensions of GLM: quasi-likelihood and Bayesian method}
\subsection{\S2.3.1 Quasi-likelihood models}
\begin{frame}{\S2.3.1 Quasi-likelihood models}
\begin{wideitemize}
\item In a GLM, the mean  necessarily implies a certain variance structure. 

\item e.g. Poisson model: we must have $\sigma^2 = \mu $, since $\phi$ is forced to be 1. 
% such as the Poisson model, 
%\[
%E(y|\textbf{x})=\mu=h(\textbf{z}^T\mbox{\boldmath $\beta$}), \ \ \mbox{var}(y|\textbf{x})=\phi v(\mu)/\omega, 
%\]
%so the aassociation between $E(y|\textbf{x})$ and  $\mbox{var}(y|\textbf{x})$ confined. May be insufficient for practical use.
%Assuming that the response variable $y$ in GLM belongs to a specific exponential family
%confines the association between $E(y|\textbf{x})=\mu=h(\textbf{z}^T\mbox{\boldmath $\beta$})$
%and $\mbox{var}(y|\textbf{x})=\phi v(\mu)/\omega$ to be of a specific type, which sometimes is
%insufficient for practical use.
\item
{\bf Quasi-likelihood models}: 

\begin{wideitemize}
\item 
Drop the exponential family (i.e. distributional) assumption\\
 $\Rightarrow$ separate the mean and
variance structures.



\item Only entails  mean and variance structural specifications (akin to structural assumptions in GLMs). 


\item Assuming that
\begin{enumerate}
\item the mean structure has been correctly specified, and
\item other appropriate regularity conditions are satisfied,
\end{enumerate}
parameters can  be estimated consistently,
and asymptotic inference is possible.

\end{wideitemize}
\end{wideitemize}
\end{frame}

\begin{frame} {Structural specifications}
%{\S2.3.1 Quasi-likelihood models}
\begin{wideitemize}
\item
The mean and variance specifications take the  forms
\begin{eqnarray}
E(y_i|\textbf{x}_i)&=&\mu_i=h(\textbf{z}_i^T\mbox{\boldmath $\beta$}) \label{mean_struct}\\
\mbox{Var}(y_i|\textbf{x}_i) &=&  \sigma^2(\mu_i)=\phi v(\mu_i)/\omega_i=\sigma^2_i(\mbox{\boldmath $\beta$})  \label{var_struct}
\end{eqnarray}
where $v(\cdot)$ is some variance function \underline{of choice}. 

\item {\bf Key feature}:  $\phi$ is not restricted to be particular number (unlike some exponential-family distributions such as Poisson where $\phi$ must be $1$). 



%\item But  $ \sigma_i^2 =  \sigma^2(\mu_i)= \phi v(\mu_i)/\omega_i$ is not necessarily correct(i.e. true). We call it a ``working" variance function. The true ( usually unknown) variance  $\sigma^2_{oi}(\textbf{x}_i) \equiv \mbox{Var}(y_i|\textbf{x}_i)$ can be different.
%


\item  Only the mean structure \eqref{mean_struct} is assumed to be correctly specified. 

\item The variance structure \eqref{var_struct} is not necessarily correct, i.e. it is a ``{\bf working}" variance function. The true  variance  structure
\[
\sigma^2_{oi}(\textbf{x}_i) \equiv \mbox{Var}(y_i|\textbf{x}_i)
\] can be different.

\end{wideitemize}
\end{frame}

\begin{frame} {Quasi-score function}

\begin{wideitemize}

\item \underline{Define} the {\bf quasi-score/generalized estimating }function (GEE):
\begin{align*}
\textbf{s}(\mbox{\boldmath $\beta$})= \sum_{i = 1}^n  {\bf s}_i({\boldsymbol \beta})  &\equiv \sum_{i=1}^n\textbf{z}_i D_i(\mbox{\boldmath $\beta$})
\sigma_i^{-2}(\mbox{\boldmath $\beta$})\left[y_i-\mu_i(\mbox{\boldmath $\beta$})\right] \\
&=Z^TD(\mbox{\boldmath $\beta$})\Sigma^{-1}(\mbox{\boldmath $\beta$})\left[\textbf{y}-
\mbox{\boldmath $\mu$}(\mbox{\boldmath $\beta$})\right],
\end{align*}

where 

\begin{wideitemize}
\item $D_i({\boldsymbol \beta})$ is  $h'(\cdot)$ evaluated at $\eta_i$
\item  $\sigma^2_i({\boldsymbol \beta})$ is the working variance for $i$
\item $D = \text{diag}(D_1, \dots, D_n)$
\item $\Sigma (\boldbeta) = \text{diag}(\sigma_1(\boldbeta)^2, \dots, \sigma_n(\boldbeta)^2)$
\item $Z^T = [{\bf z}_1, \dots, {\bf z}_n]$
\end{wideitemize}


\ \


\end{wideitemize}


\end{frame} 
% new frame
\begin{frame}{Quasi-score function}
\begin{wideitemize}
\item Very similar in form to the score function for a GLM 

\item However,  since the variance structure 
\[
\mbox{Var}(y_i|\textbf{x}_i)   =\phi v(\mu_i)/\omega_i
\]
can be of any wild form you \emph{specify}, it does not have  to be (and is generally not) the derivative of a proper log-likelihood function.  



\item \emph{Remark}:  some books filter the dispersion parameter $\phi$, which is embedded as a scale in $\Sigma$,  out of ${\bf s} ({\boldsymbol \beta})$

\end{wideitemize}

\end{frame}

%new frame
\begin{frame} {Estimation}
\begin{itemize}
\item The true $\boldbeta$ is estimated by ${\boldsymbol \hat{\boldsymbol \beta}}$, the {\bf quasi}-MLE (qMLE) that solves the \emph{estimating equation}:
\[
{\bf s}({\boldsymbol \beta}) = 0.
\] 

\ \
\item
$\hat{\mbox{\boldmath $\beta$}}$ can be found using a Newton-Raphson-like algorithm.

\ \
\item Turns out, the qMLE ${\boldsymbol \hat{\boldsymbol \beta}}$ can still consistently estimate $\boldbeta$ under regularity assumptions, despite the minimalist assumption of a correctly specified mean structure only.


\ \



\item Where does the estimation consistency come from?

\end{itemize}

\end{frame}

\begin{frame} {Estimation}
\begin{wideitemize}

\item Suppose $\boldbeta_0$ is the underlying true value for $\boldbeta$.

\item As a  \emph{random} vector-valued function in $\boldbeta$, each
\[
{\bf s}_i ({\boldsymbol \beta}) = \textbf{z}_i D_i(\mbox{\boldmath $\beta$})
\sigma_i^{-2}(\mbox{\boldmath $\beta$}) \left[y_i-\mu_i(\mbox{\boldmath $\beta$})\right], \]
 has expectation $0$ at  the true value of  $\boldsymbol \beta$, because {\bf the mean structure is assumed to be correct} and hence
 \[
 E[y_i - \mu_i (\boldbeta_0)] = 0.
 \]


\end{wideitemize}

\end{frame}

% new frame


\begin{frame} {Estimation}
\begin{wideitemize}



\item One may expect, under appropriate regularity conditions, by the law of large numbers, $n^{-1} \sum_{i = 1}^n {\bf s}_i ({\boldsymbol \beta})$ may converge to a  \emph{deterministic} function ${\bf s}_0 (\boldsymbol \beta)$ in $\boldsymbol \beta$, where ${\bf s}_0 $ is such that
\[
{\bf s}_0  (\boldbeta_0) = 0.
\] 


\item Hence ${\boldsymbol \hat{\boldsymbol \beta}}$ should be close to the true $\boldbeta_0$. 


\item   $\hat{\boldbeta}$ is also called a Z-estimator, for \emph{zeroing} out  the estimating equation. Refer to \cite[Chapter 5]{VDV1998}  for the rigorous theory. 

\end{wideitemize}

\end{frame}
% new frame
\begin{frame} {A historical note }


\begin{wideitemize}
\item Originally, \cite{wedderburn1974} did assume that the variance structure is also correctly specified. 

\item
A \textbf{quasi-likelihood} (more precisely, a quasi {\bf log}-likelihood) function can be constructed as
$\displaystyle Q(\mbox{\boldmath $\mu$}; \textbf{y})=\sum_{i=1}^n
\underbrace{\int_{y_i}^{\mu_i}\frac{y_i-t}{\omega_i^{-1} 
\phi v(t)}dt}_{\equiv Q_i (\mu_i, y_i)}$.

\item If you differentiate $Q({\boldsymbol \mu}; {\bf y})$ with respect to $\boldsymbol \beta$, 
\[
\frac{\partial Q({\boldsymbol \mu}; \textbf{y})}{\partial \boldbeta} = \sum_{i=1}^n \frac{\partial Q_i}{ \partial \mu_i} \frac{\partial \mu_i}{ \partial \boldbeta_i} = \sum_{i=1}^n
\underbrace{\textbf{z}_i D_i(\mbox{\boldmath $\beta$}) }_{\frac{\partial \mu_i}{\partial \boldbeta_i}}
 \cdot
  \underbrace{\frac{y_i - \mu_i}{ \phi v(\mu_i)/\omega_i}}_{ \frac{\partial Q_i}{ \partial \mu_i}  },
\]
which is  just ${\bf s}({\boldsymbol \beta})$.

\end{wideitemize}
\end{frame}

% another frame

% new frame
\begin{frame} {A historical note }


\begin{wideitemize}

\item 
Note $-2Q(\mbox{\boldmath $\mu$}; \textbf{y})$ is called the \textbf{scaled quasi-deviance}.
And the \textbf{quasi-deviance} is
\[
D(\textbf{y}, \mbox{\boldmath $\mu$})=-2\phi Q(\mbox{\boldmath $\mu$}; \textbf{y})
=-2\sum_{i=1}^n\int_{y_i}^{\mu_i}\frac{y_i-t}{\omega_i^{-1} v(t)}dt
\]

\item  Note that
\[
D({\bf y}, {\bf y}) = 0
\]
by setting $\boldmu = {\bf y}$.

\item  This   is analogous to the deviance in the full likelihood models, because the saturated model also estimates $\boldmu$ by ${\bf y}$ and  has deviance $0$. 



\item Different $v(\cdot)$ functions give different quasi log-likelihoods, and some correspond to real log-likelihoods for known distributions; see \cite[Chapter 9, Table 9.1]{GLM1989}.
\end{wideitemize}
\end{frame}

%
%\begin{frame}{\S2.3.1 Quasi-likelihood models (3)}
%\begin{itemize}
%\item
%The \textbf{quasi-score function} or \textbf{generalised estimating equation (GEE) function} is defined as
%$\displaystyle \textbf{s}(\mbox{\boldmath $\beta$})=\frac{\partial Q(\mbox{\boldmath $\mu$}; \textbf{y})}{\partial
%\mbox{\boldmath $\beta$}}$.
%\item
%It is not difficult to show that
%$$\textbf{s}(\mbox{\boldmath $\beta$})=\sum_{i=1}^n\textbf{z}_i D_i(\mbox{\boldmath $\beta$})
%\sigma_i^{-2}(\mbox{\boldmath $\beta$})\left[y_i-\mu_i(\mbox{\boldmath $\beta$})\right]
%=Z^TD(\mbox{\boldmath $\beta$})\Sigma^{-1}(\mbox{\boldmath $\beta$})\left[\textbf{y}-
%\mbox{\boldmath $\mu$}(\mbox{\boldmath $\beta$})\right]$$
%where $\mu_i(\mbox{\boldmath $\beta$})=h(\textbf{z}_i^T\mbox{\boldmath $\beta$})$ is assumed
%to be correctly specified, and $\sigma^2(\mu_i)=\phi v(h(\textbf{z}_i^T\mbox{\boldmath $\beta$}))/\omega_i$
%is only the \textit{working} one. Refer to page~21 for definitions of $D_i(\mbox{\boldmath $\beta$})$,
%$D(\mbox{\boldmath $\beta$})$ and $\Sigma(\mbox{\boldmath $\beta$})$.
%\item
%This $\textbf{s}(\mbox{\boldmath $\beta$})$ has the same form as we see in GLM with exponential
%family response, but has a generalised conception over that in GLM.
%\item
%Note in some papers/books, the dispersion parameter $\phi$ is filtered out of $\textbf{s}(\mbox{\boldmath $\beta$})$,
%which does not affect the estimation of $\mbox{\boldmath $\beta$}$.
%\end{itemize}
%\end{frame}

\begin{frame}{Information matrix and variance}
\begin{wideitemize}

\item
The \textbf{expected quasi-information} function is
\[
F(\boldbeta)\equiv E\left(-\frac{\partial \textbf{s}(\mbox{\boldmath $\beta$})}{\partial
\mbox{\boldmath $\beta$}^T}\right)=\sum_{i=1}^n\textbf{z}_i\textbf{z}_i^Tw_i(\mbox{\boldmath $\beta$})
=Z^TW(\mbox{\boldmath $\beta$})Z,
\]
where $W({\boldsymbol \beta}) = \text{diag}(w_1({\boldsymbol \beta}), \dots, w_n({\boldsymbol \beta}))$ for 
$
w_i ({\boldsymbol \beta}) \equiv D_i^2({\boldsymbol \beta}) \sigma_i^{-2}({\boldsymbol \beta})$.
 

\item
The variance of $\textbf{s}(\mbox{\boldmath $\beta$})$ is
$$V(\mbox{\boldmath $\beta$})\equiv\mbox{cov}(\textbf{s}(\mbox{\boldmath $\beta$}))=
\sum_{i=1}^n\textbf{z}_i\textbf{z}_i^TD_i^2(\mbox{\boldmath $\beta$})\cdot\frac{\sigma_{oi}^2(\textbf{x}_i)}{
\sigma_i^4(\mbox{\boldmath $\beta$})}.$$
\item
Both $E(\cdot)$ and $\text{cov} (\cdot)$ are with respect to the true distribution of the data $\textbf{y} = (y_1, \dots, y_n)$. 

\item In general, $F(\mbox{\boldmath $\beta$})\neq V(\mbox{\boldmath $\beta$})$ unless 
$\sigma_{oi}^2(\textbf{x}_i)=\sigma_i^2(\mbox{\boldmath $\beta$})$, i.e. the variance function 
is correctly specified.
\end{wideitemize}
\end{frame}

\begin{frame}{Covariance structure of $\hat{\boldsymbol \beta}$ }
\begin{itemize}
\item
Under regularity conditions, 
\[
\hat{\boldbeta} \stackrel{a}{\sim} N\left(\mbox{\boldmath $\beta$}, \hat{F}^{-1}\hat{V}
\hat{F}^{-1}\right),
\]
where 
\[\hat{V}=\sum_{i=1}^n\textbf{z}_i\textbf{z}_i^T 
D_i^2(\hat{\boldbeta})\cdot\frac{[y_i-h(\textbf{z}_i^T\hat{\boldbeta})]^2}{
\sigma_i^4(\hat{\boldbeta})}\]
and
\[
\hat{F} = F(\hat{\boldbeta})
\]
are estimates for  $V(\boldbeta)$ and $F(\boldbeta)$.

\ \

\item $\hat{F}^{-1}\hat{V}
\hat{F}^{-1}$ is famously known as a
\begin{itemize}
\item {\bf sandwich} estimator (for $\hat{V}$ being the ``ham"),  or 
\item {\bf robust} estimator (for its asymptotic validity regardless of whether $\text{Var}(y_i| {\bf x}_i)$ is correctly specified).
\end{itemize}




\end{itemize}
\end{frame}

\begin{frame}{Covariance structure of $\hat{\boldsymbol \beta}$ }

\begin{itemize}
\item By Taylor's approximation, we have
\[
0 = {\bf s} (\hat{\boldsymbol \beta}) \sim_a   {\bf s} ({\boldsymbol \beta})  - F({\boldsymbol \beta})(\hat{\boldsymbol \beta} - {\boldsymbol \beta}).
\]
(Apparently, we have replaced the Hessian by its expectation)

\ \

\item Hence,
\[
\text{cov } (\hat{\boldsymbol \beta} - {\boldsymbol \beta}) \approx F^{-1}({\boldsymbol \beta}) V(\boldbeta) F^{-1}({\boldsymbol \beta})
\]

\ \
\item However, if the variance structure happens to be  correctly specified, we have $F(\boldbeta)  = V(\boldbeta)$, which gives
\[
\text{cov } (\hat{\boldsymbol \beta} - {\boldsymbol \beta}) \approx F^{-1}({\boldsymbol \beta})
\]

\ \
\item In fact, this heuristics also explains the  covariance structure  $F^{-1}(\boldbeta)$ of the MLE $\hat{\boldbeta}$, where $F(\boldbeta)$ is the actual Fisher information under the full set of GLM assumptions.



\end{itemize}


\end{frame}




\begin{frame}{Covariance structure of $\hat{\boldsymbol \beta}$ (Summary)}
\begin{itemize}
\item If only the mean structure \eqref{mean_struct} is correctly specified:

\ \
\begin{wideitemize}
\item $\hat{\boldbeta} \sim_a N( \boldbeta, \hat{F}^{-1} \hat{V} \hat{F}^{-1})$ ({\bf robust/sandwich } covariance estimate)
\item By tracing where $\phi$ figures in $F$ and $V$, no $\phi$ appears in $\hat{F}^{-1} \hat{V} \hat{F}^{-1}$ due to cancellation. 
\item $\Rightarrow$  No need to estimate $\phi$ in practice! Note that our Pearson-residual-based estimator $\hat{\phi}$ only makes sense when the variance structure is correctly specified.
\end{wideitemize}

\ \

\item If  both the mean and variance structures \eqref{mean_struct} and \eqref{var_struct} are correctly specified:

\ \

\begin{wideitemize}
\item $\hat{\boldbeta} \sim_a N( \boldbeta, \hat{F}^{-1} )$ ({\bf naive}/{\bf model based } covariance estimate)

\item In this case, both $\hat{F}^{-1} $ and $\hat{F}^{-1} \hat{V} \hat{F}^{-1}$ are  (asymptotically) valid and they should also be close in value.  However, one may argue that $\hat{F}^{-1} $ is slightly favored  since less estimation is needed. 
\end{wideitemize}


\end{itemize}

\end{frame}

\begin{frame} {Efficiency (1) }



\begin{wideitemize}
\item Question: When we assume both the mean and variance structures are correctly specified, $F^{-1}$, the inverse of the quasi-information,  is a valid covariance matrix for $\hat{\boldbeta}$. But then, $F$ is also a proper Fisher information matrix when there exists a GLM  (a full-blown likelihood model) that admits the value of the  dispersion parameter $\phi$ embedded in $F$. Intuitively, shouldn't we lose some efficiency (i.e. having a smaller ``information matrix", in matrix sense) since now we've only correctly assumed the first two moments?  





\item Answer: Generally, we do lose efficiency. In particular, the quasi information $F$ is generally ``less"  (in matrix sense) than the Fisher information, $- E[\frac{\partial^2 \ell }{\partial \boldbeta^2}]$,   defined by the true (unknown) log-likelihood $\ell$. However, if $\ell$ {\bf happens to be an exponential family} log-likelihood, then we're lucky enough not to lose efficiency, because the qMLE coincides with the MLE which already exploits all distributional information by design.



\end{wideitemize}
\end{frame}

\begin{frame} {Efficiency (2)}
Under the minimal assumptions of a correctly specified mean structure and certain regularity conditions, one can show that 
\[
\left(\sum_{i=1}^n\textbf{z}_i\textbf{z}_i^TD_i^2(\mbox{\boldmath $\beta$})\cdot\frac{1}{
\sigma_{oi}^2(\mbox{\boldmath $\beta$})} \right)^{-1}.
\]
is a lower bound (in matrix sense) for the limiting-covariance of a quasi-likelihood type estimator; see Section 5  of \cite{GMT1984}  for a proof. 

\ \

Hence if we can correctly specify the variance structure, we can get $V^{-1}$ to be the matrix in display, and thus achieve efficiency. 






\end{frame} 

\begin{frame} {Testing in  quasi-likelihood model}
\begin{itemize}

\item
Testing $H_0: C\mbox{\boldmath $\beta$}=\mbox{\boldmath $\xi$}$ vs. $H_1:C\mbox{\boldmath $\beta$}
\neq \mbox{\boldmath $\xi$}$ is still possible by properly modified Wald or score statistics, e.g. a modified
Wald statistic is 
$$W_m=(C\hat{\boldbeta}-\mbox{\boldmath $\xi$})^T[C
\hat{F}^{-1}\hat{V}\hat{F}^{-1}C^T]^{-1}
(C\hat{\boldbeta}-\mbox{\boldmath $\xi$}),$$
with limiting distribution still being $\chi^2(r)$, $r=\mbox{Rank}(C)$.

\ \

\item
Quasi-LR (likelihood ratio) test statistic in general is not equivalent to modified Wald or modified
score test, and may not have a limiting $\chi^2$ distribution (Foutz \& Srivastava, 1977).


\end{itemize}


\end{frame}



\begin{frame} {Variance function with unknown parameter}

\begin{wideitemize}
\item So far,  the working variance function is a known function of the mean, $v(\mu_i)$.

\item Some authors (e.g. Nelder \& Pregibon, 1987) relaxed this and specify
\[
var(y_i | {\bf x}_i) = \phi v(\mu_i, \theta)/\omega_i,
\]
i.e., $v(\cdot)$ has an extra  parameter $\theta$ (not to be confused with the canonical parameters in the GLM!). 

\item We will investigate this in the example to follow. 

\end{wideitemize}
\end{frame}


\subsection{Example: Cell Differentiation Data}

%\begin{frame}{Agenda today}
%\begin{itemize}
%\item introduce \texttt{knitr}
%
%
%
%\item Look at cell differentiation data (Example 2.3, 2.6, 2.7) in F \& T. 
%
%%\item Finish Example 2.7 in F \& T on log linear Poisson model, with \texttt{R}. (Recall log linear link means $\mu = \exp(\eta)$)
%
%
%\item  Reference:  ``Modern Applied Statistics  With S"  by {\bf Venables and Ripley} ({\bf V \& R} ), downloadable from unimelb library. 
%
%\end{itemize}
%\end{frame}


\begin{frame}{Cell Differentiation data (F \& T, Example 2.3, 2.6, 2.7)}


%\item Study the immuno-activating effect of  two agents, TNF (tumor necrosis factor) and TFN (interferon) on inducing cell differentiation
%
%\item Log-linear Poisson mean model:
%\[
%\mu = E[y|TNF, IFN] = \exp( \beta_0 + \beta_1 TNF + \beta_2 IFN + \beta_3 TNF*IFN)
%\]
 \includegraphics[height=8cm]{celldiff}

\end{frame}

\begin{frame}

A handy reference  relevant to our R code associated with  these examples is 
\[
\text{``Modern Applied Statistics  With S"  by {\bf Venables and Ripley} ({\bf V \& R})},
\]
downloadable from unimelb library. 
\end{frame}

\begin{frame}{Cell Differentiation data (F \& T, Ex 2.3, 2.6, 2.7)}

\begin{itemize}
\item Study the immuno-activating effect of  two agents, TNF (tumor necrosis factor) and TFN (interferon) on inducing cell differentiation

\ \

\item Log-linear Poisson mean model:
\[
\mu = E[y|TNF, IFN] = \exp( \beta_0 + \beta_1 TNF + \beta_2 IFN + \beta_3 TNF*IFN)
\]

\ \

\item Doses of TNF  and IFN are treated as numbers (instead of factors)

\ \

\item Regression with both the  Poisson GLM (Ex. 2.3 and 2.6) and quasi-likelihood approaches (Ex. 2.7).


\end{itemize}

\end{frame}


%
%
%\begin{frame} {Recap a few quick facts first..}
%In the quasi-likelihood inference world, we only assume the mean and variance structure forms:
%\[
%E(y|{\bf x}) = \mu = h(z'\beta), \ \ Var(y|{\bf x}) = \sigma^2(\mu) = \phi v(\mu)
%\]
%\begin{wideitemize}
%\item Generally, for consistent $\beta$ estimate, we only need to have the mean structure correctly specified. In this case, 
%\begin{wideitemize}
%\item $\hat{\beta} \sim_a N( \beta, \hat{F}^{-1} \hat{V} \hat{F}^{-1})$ ({\bf robust/sandwich } covariance estimate)
%\end{wideitemize}
%\item But Wedderburn(1974) originally assumed both mean and variance structures are correctly specified. In this case
%\begin{wideitemize}
%\item $\hat{\beta} \sim_a N( \beta, \hat{F}^{-1} )$ ({\bf model based } covariance estimate)
%
%\item (In principle, both $\hat{F}^{-1} $ and $\hat{F}^{-1} \hat{V} \hat{F}^{-1}$ are valid and they should also be close in value.  However, one may argue that $\hat{F}^{-1} $ is slightly favored in this case since less estimation is needed. )
%\end{wideitemize}
%\end{wideitemize}
%\end{frame}

\begin{frame} { Poisson GLM (full likelihood) inference (Ex. 2.3 \& 2.6)}

\begin{itemize}
\item F \& T used  a $\chi^2(\underbrace{12}_{n - p})$ reference null to test the full Poisson GLM against the saturated model, and reject the Poisson GLM based on the deviance 142.4 . 


\ \
\item However, the dimension of the saturated model $n$ is the sample size (16) itself, so chi-squared asymptotic isn't necessarily valid (dimensionality growing with $n$). \\

\ \

$\Rightarrow$ take $\chi^2$ test with a grain of salt. (The biostatistician Johnathan Bartlett pointed out the same issue: \href{https://thestatsgeek.com/2014/04/26/deviance-goodness-of-fit-test-for-poisson-regression/}{\textcolor{red}{click link here}})

\ \

%\item Recall: for binomial data we also prefer that sample size is large and each group has enough data 

\item The $\chi^2$ asymptotic MAY be reliable for Poisson regression deviance when the mean for each data point $y_i$ is large enough; see the discussion of Bartlett above and $V\& R$ (p.187) for the detail. 

\end{itemize}


\end{frame}




%
%\begin{frame}
%We will continue our discourse on  GLMs
%
%\begin{figure}
% \centering
%\includegraphics[width=0.5\textwidth]{NandW}
%% \label{figure:example}
%\end{figure}
%
%
%\end{frame}
%
%\begin{frame}[fragile] \frametitle{Outline}
%\tableofcontents
%\end{frame}
%%
%%\section{\S2.1 Univariate GLM}\label{sec:2.1}
%%%\subsection{\S2.1.1 Data}
%%%
%%%\begin{frame}{\S2.1.1 Data (ungrouped)}
%%%\begin{itemize}
%%%\item
%%%$y$, univariate response
%%%
%%%\item  $\textbf{x}=(x_1,\cdots,x_m)^T$, $m$ covariate variables
%%%
%%%\item
%%%A sample of $n$ individuals: 
%%%$$(y_i, \textbf{x}_i^T)=(y_i, x_{i1},\cdots, x_{im}), \quad i=1,2,\cdots,n$$
%%% in matrix form
%%%\begin{equation*}
%%%\textbf{y}=\left[\begin{array}{c} y_1 \\ \vdots \\ y_i \\ \vdots \\ y_n\end{array}\right], \quad
%%%X=\left[\begin{array}{ccc} x_{11} & \cdots & x_{1m} \\ \vdots & \cdots & \vdots \\
%%%x_{i1} & \cdots & x_{im} \\ \vdots & \cdots & \vdots \\ x_{n1} & \cdots & x_{nm} \end{array}
%%%\right]_{n\times m}
%%%\end{equation*}
%%%\end{itemize}
%%%\end{frame}
%%%
%%%\begin{frame}{\S2.1.1 Data (grouped)}
%%%\begin{itemize}
%%%\item 
%%%Sometimes two or more individuals have covariate values, i.e. some rows of $X$ are identical
%%%
%%%\item Individuals with  same $\bf x$ (say $n_i$ of them)  can be grouped. The group mean $\bar{y}$ responses is taken to be the group response. 
%%%
%%%
%%%\item $g = $ number of different ${\bf x}_i$
%%%\ \
%%%
%%%
%%%
%%%\begin{center}
%%%\begin{tabular}{cccc}
%%%& sizes & group means &  \\ \\
%%%Group 1 &  \multirow{5}{*}{$\left[\begin{array}{c} n_1 \\ \vdots \\ n_i \\ \vdots \\ n_g\end{array}\right]$,} 
%%%& \multirow{5}{*}{$\bar{\textbf{y}}=\left[\begin{array}{c} \bar{y}_1 \\ \vdots \\ \bar{y}_i \\ \vdots 
%%%\\ \bar{y}_g\end{array}\right]$,} & \multirow{5}{*}{$X=\left[\begin{array}{ccc} x_{11} & \cdots & x_{1m} \\ 
%%%\vdots & \cdots & \vdots \\ x_{i1} & \cdots & x_{im} \\ \vdots & \cdots & \vdots \\ x_{g1} & \cdots & x_{gm}
%%%\end{array}\right]_{g\times m}$} \\ $\vdots$ & & & \\ Group $i$ & & & \\ $\vdots$ & & & \\ Group $g$ & & &
%%%\end{tabular}
%%%\end{center}
%%%
%%%\end{itemize}
%%%\end{frame}
%%%
%%%\subsection{\S2.1.2 GLM}
%%%\begin{frame}{\S2.1.2 GLM (Its form)}
%%%\noindent
%%%\textbf{1. Distributional assumption}
%%%\begin{itemize}
%%%\item
%%%Given $\textbf{x}_i$'s, the $y_i$'s are (conditionally) independent.
%%%
%%%\item
%%%The (conditional) pdf of $y_i$ is
%%%$$f(y_i|\theta_i, \phi, \omega_i)=\exp\left\{\frac{y_i\theta_i-b(\theta_i)}{\phi}\cdot \omega_i +
%%%c(y_i,\phi, \omega_i)\right\}, \quad \mbox{where}$$
%%%\begin{itemize}
%%%\item
%%%$\theta_i$ -- \textit{natural} (or \textit{canonical}) parameter;
%%%\item
%%%$\phi$ --  \textit{scale} (or \textit{dispersion}) parameter;
%%%\item
%%%$b(\cdot)$ and $c(\cdot)$ --  specific functions determined by the actual distribution of $y_i$;
%%%\item
%%%$\omega_i$ is the \textit{weight}. 
%%%\begin{itemize}
%%%\item for ungrouped data, $\omega_i=1$ ; 
%%%\item for grouped data, $\omega_i$ equals $n_i$ when the group mean is the response, and equals $1/n_i$
%%%when group total is the response; 
%%%\end{itemize}
%%%\end{itemize}
%%%\end{itemize}
%%%\end{frame}
%%%
%%%\begin{frame}
%%%
%%%\begin{itemize}
%%%\item Jorgensen (1992) also calls this an exponential dispersion model
%%%\item $\mu_i = E(y_i|\textbf{x}_i)$ still denotes the mean
%%%\item We will model the dependence of $\mu_i $ on covariates via a link/response function
%%%%The (conditional) distribution of $y_i$ belongs to a \textit{simple exponential family} with  
%%%%$E(y_i|\textbf{x}_i)=\mu_i$, and possibly a common \textit{scale} (or \textit{dispersion}) parameter $\phi$.
%%%\end{itemize}
%%%
%%%\end{frame}
%%%
%%%
%%%\begin{frame}{\S2.1.2 GLM (Mean structure via link)}
%%%\noindent
%%%\textbf{2. Structural assumption}
%%%
%%%Mean $\mu_i$ is related to the \textit{linear predictor} $\eta_i=\textbf{z}_i^T\mbox{\boldmath $\beta$}$ by
%%%$$\mu_i=h(\eta_i)=h(\textbf{z}_i^T\mbox{\boldmath $\beta$}), \quad\mbox{respectively},\quad
%%%\eta_i=g(\mu_i)$$
%%%where
%%%\begin{itemize}
%%%\item
%%%$h(\cdot)$ is the \textit{response function};
%%%\item
%%%$g(\cdot)$ is the \textit{link function}, $g=h^{-1}$;
%%%\item
%%%$\mbox{\boldmath $\beta$}$ is a $p\times 1$ unknown \textit{parameter} vector;
%%%\item
%%%$\textbf{z}_i$ is a $p\times 1$ unknown \textit{design} vector, $\textbf{z}_i=\textbf{z}(\textbf{x}_i)$ (a transformation of $\textbf{x}_i$ ).
%%%\end{itemize}
%%%
%%%
%%%\end{frame}
%%%
%%%
%%%\begin{frame}
%%%\begin{itemize}
%%%
%%%\item
%%%Thus a  GLM has  3 components
%%%\begin{enumerate}
%%%\item
%%%type of the exponential family
%%%\item
%%%the response/link function
%%%\item
%%%the design vector 
%%%\end{enumerate}
%%%\item 
%%%Obviously, ${\bf z}(\textbf{x}_i)$ determines the mean $\mu_i$ via the link
%%%
%%%\item The exponential family construction must also have led to a relationship between $\mu_i$ and the exponential family parameters. What is this relationship?
%%%\end{itemize}
%%%
%%%\end{frame}
%%%
%%%\begin{frame}{\S2.1.2 GLM (Mean determines variance)}
%%%\noindent
%%%\textbf{Remark 1.}
%%%\begin{itemize}
%%%\item
%%%One can derive that
%%%$$\mu_i=b^\prime (\theta_i)=\frac{\partial b(\theta_i)}{\partial \theta_i} \quad
%%%\mbox{and}\quad \mbox{var} (y_i|\textbf{x}_i)\stackrel{\mbox{\scriptsize denoted}}{=}\sigma^2(\mu_i)=\frac{\phi}{\omega_i} v(\mu_i)$$
%%%with
%%%$v(\mu)=b^{\prime\prime}(\theta)=\frac{\partial^2 b(\theta)}{\partial \theta^2}$ being called the 
%%%\textit{variance function}.
%%%
%%%\item (see blackboard derivation)
%%%
%%%\item  $Note: $$b'$ is invertible mapping from $\theta$ to $\mu$, since $b''(\theta)$, being a variance, is $>0$.  (Hence the notation $v(\mu)$)
%%%\item
%%%$\Rightarrow$ Modeling the  mean  as $\mu=h(\textbf{z}^T\mbox{\boldmath $\beta$})$  also determines the variance structure
%%%
%%%\end{itemize}
%%%\end{frame}
%%%
%%%\begin{frame}
%%%\begin{itemize}
%%%\item
%%%different from linear models (LM) where $E(y|\textbf{x})$ and 
%%%$\mbox{var}(y|\textbf{x})$ are unrelated. 
%%%\item \textbf{quasi-likelihood} approach (Ch2.3) allows us to model the mean and variance structure separately. 
%%%%The relation between $\mu$ and $v(\mu)$ may also be specified through empirical evidence,
%%%%resulting in the so-called \textbf{quasi-likelihood} approach.
%%%\end{itemize}
%%%
%%%\end{frame}
%%%
%%%\begin{frame}{\S2.1.2 GLM (Natural/canonical link)}
%%%\noindent
%%%\textbf{Remark 2.}
%%%\begin{wideitemize}
%%%\item
%%%Choosing $h(\cdot)$ and $g(\cdot)$ should depend on the specific distribution of $y$.
%%%\item
%%%Note that assuming $\theta\equiv \theta(\mu)=\eta\equiv \textbf{z}^T\mbox{\boldmath $\beta$}$
%%%makes estimating $\mbox{\boldmath $\beta$}$ much easier.
%%%
%%%%Since $\mu=b^\prime(\theta)$ implies $\theta=(b^\prime)^{-1}(\mu)$, it follows that
%%%%$\eta=\theta(\mu)=(b^\prime)^{-1}(\mu)$ and $g(\mu)=(b^\prime)^{-1}(\mu)$ under this assumption.
%%%\item
%%%We call $\eta=(b^\prime)^{-1}(\mu)$ the \textit{natural} (or \textit{canonical}) link function,
%%%and $\mu=b^{\prime}(\eta)$ the \textit{natural response function}.
%%%%\item
%%%%It can be shown that the natural link function is
%%%%\begin{enumerate}
%%%%\item
%%%%$\eta=\mu$ if $y\sim \mbox{Normal}(\mu, \sigma^2)$; \textit{identity link};
%%%%\item
%%%%$\eta=\log\mu$ if $y\sim \mbox{Poisson}(\mu)$; \textit{log link};
%%%%\item
%%%%$\eta=\log\frac{p}{1-p}=\log\frac{\mu}{n_g-\mu}$ if $y\sim\mbox{Binomial}(n_g, p)$; \textit{logit (or logistic) link}.
%%%%\end{enumerate}
%%%\end{wideitemize}
%%%\end{frame}
%%%
%%%\begin{frame}{\S2.1.2 GLM (design vector specifics)}
%%%\noindent
%%%\textbf{Remark 3.}
%%%\begin{itemize}
%%%\item
%%%Concerning the design vector, an intercept term is often included
%%%$$\textbf{z}=(1,\textbf{z}^{*T})^T$$
%%%where $\textbf{z}^*=\textbf{z}^*(\textbf{x})$ is a function of $\textbf{x}$, including e.g. $\log(x_1)$,
%%%$x_2^2$, $\sin (x_m)$ etc. as well as $\textbf{x}$ itself.
%%%%\item
%%%%A categorical covariate with $k$ levels has to be coded by a $(k-1)\times 1$ dummy vector, or an equivalent
%%%%vector of size $q=k-1$ for it to be properly included in the design vector.
%%%\item A categorical covariate with $k$ levels can  be coded by a dummy vector ${\bf x} = (x^{(1)}, \dots, x^{(k-1)})$ of length $k-1$, where
%%%\[ x^{(j)} =
%%%  \begin{cases}
%%%    1      & \quad \text{if } j = 1, \dots, k-1\\
%%%   0  & \quad \text{if } j = k
%%%  \end{cases}
%%%\]
%%%(k-th category reserved as reference)
%%%\end{itemize}
%%%\end{frame}
%%%
%%%\begin{frame}{\S2.1.2 GLM (Grouped and ungrouped)}
%%%\noindent
%%%\textbf{Remark 4.}
%%%\begin{itemize}
%%%
%%%\item 
%%%Suppose ungrouped data 
%%%$(y_i, \textbf{x}_i^T)$, $i=1,\cdots, n$, are modelled by a GLM with
%%%$$E(y_i|\textbf{x}_i)=\mu_i=h(\textbf{z}^T\mbox{\boldmath $\beta$}), \quad
%%%\mbox{var}(y_i|\textbf{x}_i)=\phi v(\mu_i).$$
%%%\item
%%%If the data are grouped, so  $\bar{y}_i$ is  the response observation for group $i$ of
%%%size $n_i$ with covariate vector $\textbf{x}_i$, then
%%%% the distribution of $\bar{y}_i$ is within the exponential
%%%%family again, and
%%%$$E(\bar{y}_i|\textbf{x}_i)=\mu_i=h(\textbf{z}^T\mbox{\boldmath $\beta$}), \quad
%%%\mbox{var}(\bar{y}_i|\textbf{x}_i)=n_i^{-1}\phi v(\mu_i).$$
%%%So the same GLM applies to the grouped data as well if the variance function is properly adjusted 
%%%by a factor of $n_i^{-1}$, i.e. defining the weight $\omega_i=n_i$.
%%%\item By same token, if the group \emph{total} is used as response, one should define $\omega_i=1/n_i$.
%%%\end{itemize}
%%%\end{frame}
%%
%%%\begin{frame}{GLM  Common exponential families}
%%%
%%%Recall the table from last time
%%%\begin{table} [!htbp]
%%%\caption{Common exponential families} \label{tab2.1}
%%%\begin{center}
%%%{\tiny
%%%\begin{tabular}{|c|c|c|c|c|c|c|}
%%%\hline
%%%Distribution & $\theta (\mu)$ & $b(\theta)$ & $\phi$ & $E(y)=b^\prime(\theta)$ & $v(\mu)=b^{\prime\prime}(\theta)$
%%%& $\mbox{var}(y)=b^{\prime\prime}(\theta)\phi/\omega$ \\ \hline
%%%$\begin{array}{c} \mbox{Normal} \\ N(\mu,\sigma^2)\end{array}$ & $\mu$ & $\theta^2/2$ & $\sigma^2$ 
%%%& $\mu=\theta$ & 1 & $\sigma^2/\omega$ \\
%%%\hline
%%%$\begin{array}{c} \mbox{Bernoullil} \\ B(1,\pi)\end{array}$ & $\log\frac{\pi}{1\!-\!\pi}$ 
%%%& $\log (1\!+\!e^\theta)$ & 1 & $\pi=\frac{e^\theta}{1+e^\theta}$ &
%%%$\pi(1-\pi)$ & $\pi(1-\pi)/\omega$ \\ \hline
%%%$\begin{array}{c} \mbox{Poisson} \\ P(\lambda)\end{array}$ & $\log\lambda$ & $e^\theta$ 
%%%& 1 & $\lambda=e^\theta$ & $\lambda$ & $\lambda/\omega$ \\
%%%\hline
%%%$\begin{array}{c} \mbox{Gamma} \\ G(\mu,\nu)\end{array}$ & $-1/\mu$ & $-\log(-\theta)$ & $\nu^{-1}$ 
%%%& $\mu=-1/\theta$ & $\mu^2$ & $\mu^2\nu^{-1}/\omega$ \\ \hline
%%%$\begin{array}{c} \mbox{InverseGaussian} \\ IG(\mu,\sigma^2)\end{array}$ & $-1/2\mu^2$ & $-(-2\theta)^{1/2}$
%%%& $\sigma^2$ &  $\mu=(-2\theta)^{-1/2}$ & $\mu^3$ & $\mu^3\sigma^2/\omega$ \\ \hline
%%%\end{tabular}}
%%%\end{center}
%%%\end{table}
%%%%Remember that if $\theta = \eta = \textbf{z}^T\mbox{\boldmath $\beta$}$
%%%%\begin{itemize}
%%%%\item 2nd column gives the canonical link function
%%%%\item 3rd column gives the canonical response function
%%%%\end{itemize}
%%%\end{frame}
%%
%%
%%%\subsection{\S2.1.3 GLM for continuous responses}\label{sec:2.1.3}
%%%\begin{frame}{\large \S2.1.3 GLM for normal, Gamma \& inverse Gaussian responses }
%%%
%%%When $y\sim \mbox{Normal}(\mu,\sigma^2)$, $\mu=\eta=\textbf{z}^T\mbox{\boldmath $\beta$}$
%%%if the natural link is used. Also $\theta=\mu$, $b(\theta)=\frac{1}{2}\theta^2$, $\phi=\sigma^2$ and
%%%$v(\mu)=1$.
%%%
%%%\end{frame}
%%%
%%%\begin{frame}{\large \S2.1.3 GLM for normal, Gamma \& inverse Gaussian responses}
%%%When $y$ has a Gamma distribution (e.g. survival time, insurance claims, amount of rainfall, etc.), the pdf is
%%%$$f(y|\mu,\nu)=\frac{1}{\Gamma(\nu)}\left(\frac{\nu}{\mu}\right)^\nu y^{\nu-1}\exp\left(-\frac{\nu}{\mu}y\right),
%%%\quad y\geq 0,$$
%%%where $\mu>0$ is the mean parameter, and $\nu>0$ the shape parameter.
%%%\begin{itemize}
%%%\item
%%%The natural link is $\eta=-\theta=\mu^{-1}$. Linear predictor $\eta=\textbf{z}^T\mbox{\boldmath $\beta$}$.
%%%Hence $\mu^{-1}=\textbf{z}^T\mbox{\boldmath $\beta$}$.
%%%\item
%%%Alternatively, identical link $\eta=\mu$ and log-link $\eta=\log \mu$ can be used to avoid the restriction. (These will give positive mean $\mu$ so they are allowed)
%%%\end{itemize}
%%%\end{frame}
%%%
%%%\begin{frame}{\large \S2.1.3 GLM for normal, Gamma \& inverse Gaussian responses }
%%%\begin{itemize}
%%%\item
%%%When $y$ has an inverse Gaussian distribution $\mbox{IG}(\mu,\sigma^2)$ (useful for modelling life time etc.), the pdf is
%%%\begin{eqnarray*}
%%%f(y|\mu,\sigma^2) &\!\!\!=\!\!\!&\left(\frac{1}{2\pi\sigma^2 y^3}\right)^{1/2}\exp\left(-\frac{(y-\mu)^2}{2\sigma^2\mu^2y}
%%%\right), \quad\!\!\! y>0, \mu>0, \sigma^2>0 \\
%%%&\!\!\!=\!\!\!& \left(\frac{1}{2\pi\sigma^2 y^3}\right)^{1/2}\exp\left\{\frac{-y\frac{1}{2\mu^2}+\frac{1}{\mu}}{\sigma^2}
%%%-\frac{1}{2\sigma^2 y}\right\}
%%%\end{eqnarray*}
%%%\item
%%%Hence natural parameter $\theta=-\frac{1}{2\mu^2}$;  dispersion parameter $\phi=\sigma^2$.
%%%\begin{eqnarray*}
%%%b(\theta)=-\frac{1}{\mu}=-(-2\theta)^{1/2} \quad \Rightarrow \quad b^\prime(\theta)=(-2\theta)^{-1/2}=\mu \\
%%%\Rightarrow \; b^{\prime\prime}(\theta)\!=\!(-2\theta)^{-3/2}\!=\!\mu^3 \;\; \Rightarrow \;\;
%%%v(\mu)\!=\!\mu^3 \;\; \Rightarrow \;\; \mbox{Var}(y)\!=\!\sigma^2\mu^3
%%%\end{eqnarray*}
%%%The natural link is obtained by setting $\eta=\theta, \quad\Rightarrow\quad \eta=-\frac{1}{2\mu^2}$, i.e.
%%%$$\eta=\theta=(b^\prime)^{-1}(\mu)=-(2\mu^2)^{-1}.$$
%%%In practice $\eta=\mu^{-2}$ is used as the natural link.
%%%\end{itemize}
%%%\end{frame}
%%
%%\subsection{\S2.1.4 GLM for binary and binomial responses}\label{sec:2.1.4}
%%\begin{frame}{\S2.1.4 GLM for binary and binomial responses }
%%\begin{wideitemize}
%%\item
%%Suppose $y_i|\textbf{x}_i \sim \mbox{Binomial}(m_i,\pi_i)$ and $\bar{y}_i=\frac{y_i}{m_i}$ (scaled binomial) are used as the
%%response, $i=1,\cdots,n$. Then the pmf is {\footnotesize
%%$$f(y_i|m_i, \pi_i,\textbf{x}_i)\!=\!{m_i\choose y_i}\pi_i^{y_i}(1-\pi_i)^{m_i\!-\!y_i}\!=\!{m_i\choose y_i}\exp\left\{
%%m_i\bar{y}_i\log\frac{\pi_i}{1\!-\!\pi_i}+m_i\log (1\!-\!\pi_i)\right\}$$}
%%\item
%%Natural parameter $\theta_i=\log\frac{\pi_i}{1-\pi_i}$; dispersion $\phi=1$; weight $\omega_i=m_i$.
%%\item The scaled binomial distribution may treated as the average from grouped binary observations. (group size $ = m_i$ )
%%\end{wideitemize}
%%\end{frame}
%%
%%
%%
%%\begin{frame}
%%\begin{itemize}
%%\item
%%$\mu_i^*=E(y_i|\textbf{x}_i)=m_i\pi_i=m_i\frac{\exp(\theta_i)}{1+\exp(\theta_i)}$
%%{\small $$\Rightarrow \quad \mu_i=E(\bar{y}_i|\textbf{x}_i)=\pi_i=\frac{\exp(\theta_i)}{1+\exp(\theta_i)}$$}
%%\item
%%Also $b(\theta_i)=-\log (1-\pi_i)=\log (1+e^{\theta_i}) \quad \Rightarrow \quad 
%%b^\prime(\theta_i)=\frac{e^{\theta_i}}{1+e^{\theta_i}}=\pi_i$
%%{\small $$\Rightarrow\quad b^{\prime\prime}(\theta_i)=\frac{e^{\theta_i}}{(1+e^{\theta_i})^2}=\pi_i(1-\pi_i)
%%=v(\pi_i)$$}
%%$\Rightarrow\quad \mbox{Var}(\bar{y}_i|\textbf{x}_i)=\sigma^2(\mu_i)=\phi v(\mu_i)/\omega_i
%%=m_i^{-1}\pi_i(1-\pi_i)$.
%%
%%\end{itemize}
%%\end{frame}
%%
%%
%%
%%\begin{frame}{\S2.1.4 GLM for binary and binomial responses }
%%\begin{itemize}
%%\item
%%The natural link is the logit link
%%\[\eta_i=\theta_i=\log\frac{\pi_i}{1-\pi_i}=(b^\prime)^{-1}(\pi_i),\]
%%called the \textbf{logit} (or \textbf{logistic}) \textbf{link}, with the corresponding {\bf logistic} response function
%%\[
%%\pi_i = \frac{\exp(\eta_i)}{1 + \exp(\eta_i)}. 
%%\]
%%
%%This gives the well-known \textit{logit regression model}:
%%$\displaystyle \log\frac{\pi_i}{1-\pi_i}\!=\!\textbf{z}_i^T\!\mbox{\boldmath $\beta$}$
%%
%%
%%
%%
%%\end{itemize}
%%\end{frame}
%%
%%
%%
%%\begin{frame}
%%
%%Other models of the form 
%%\[
%%\pi = h(\eta)
%%\]
%%for  strictly monotonous distribution functions $h$  on the real line:
%%\begin{enumerate} 
%%\item
%%\textbf{Probit Model}:  \[\pi_i = \Phi(\textbf{z}_i^T\!\mbox{\boldmath $\beta$}), \] where $\Phi(\cdot)$ is the N(0,1) cdf. 
%%% $\eta_i=\Phi^{-1}(\pi_i)$, where 
%% 
%%% This results in the
%%%\textit{probit model} $$\Phi^{-1}(\pi_i)=\textbf{z}_i^T\!\mbox{\boldmath $\beta$}.$$
%%\item
%%\textbf{Complementary log-log Model}:
%%\[
%%\pi_i = h(\eta_i)\!=\!1\!-\!\exp(\!-\!\exp(\!\eta_i\!)\!)
%%\]
%%with link function 
%%$
%% \textbf{z}_i^T\!\mbox{\boldmath $\beta$}= \log (-\log(1-\pi_i)).
%%$
%%%
%%%
%%% $\eta_i=g(\pi_i)=\log (-\log(1-\pi_i))$, resulting in the
%%%\textit{complementary log-log model} 
%%%$$\log (-\log(1-\pi_i))=\textbf{z}_i^T\!\mbox{\boldmath $\beta$}.$$
%%%Note the corresponding \textbf{response function} is $h(\eta_i)\!=\!1\!-\!\exp(\!-\!\exp(\!\eta_i\!)\!)$.
%%\item
%%\textbf{Complementary log Model} 
%%\[
%% \pi_i = h(\eta_i)  = 
%%  \begin{cases} 
%%   1 - \exp(- \eta_i) & \text{if } \eta_i \geq 0 \\
%%   0       & \text{if } \eta_i < 0
%%  \end{cases}
%%\]
%%
%%%\[
%%%\pi_i = h(\eta_i) = 
%%%\eta_i=g(\pi_i)=-\log(1-\pi_i),
%%%\] 
%%%resulting in the
%%%\textit{complementary log model}
%%with link function $-\log(1-\pi_i)=\textbf{z}_i^T\!\mbox{\boldmath $\beta$}.$
%%%\item
%%%Any inverse cdf can serve as a link function for GLM with binomial responses.
%%\end{enumerate}
%%\end{frame}
%%
%%
%%\begin{frame}
%%
%%
%%
%%\begin{figure}
%% \centering
%%\includegraphics[width=1\textwidth]{BinaryResponseFunctions.png}
%% \label{figure:example}
%%\end{figure}
%%
%%\end{frame}
%%
%%\begin{frame}{\S2.1.4 Over-dispersion in binomial GLM}
%%\begin{itemize}
%%
%%\item In applications it often happens that $\mbox{var}(\bar{y}_i|\textbf{x}_i) > m_i^{-1}\pi_i(1-\pi_i)$,  making the binomial modeling not accurate. 
%%
%%\item Two main reasons for this  \textbf{over-dispersion}:
%%\begin{itemize}
%%\item
%%Unobserved heterogeneity, i.e. different Bernoulli trials in the data have different probabilities of ``success".
%%\item
%%Positive correlation between individual Bernoulli trials. (For example, different units in a group may be correlated, such as members in a household)
%%\end{itemize}
%%%
%%%\item
%%%Sometimes $\mbox{var}(\bar{y}_i|\textbf{x}_i) > m_i^{-1}\pi_i(1-\pi_i)$, so-called \textbf{over-dispersion},
%%%where $\bar{y}_i=m_i^{-1}y_i$ and $y_i$ is the number of ``successes" in $m_i$ Bernoulli trials.
%%%This implies the actual distribution of $y_i$ is not binomial. Two causes for this happening in practice.
%%%\begin{enumerate}
%%%\item
%%%Unobserved heterogeneity, i.e. different Bernoulli trials in the data have different probabilities of ``success".
%%%\item
%%%Positive correlation between individual Bernoulli trials.
%%%\end{enumerate}
%%
%%%\[\mbox{var}(\bar{y}_i|\textbf{x}_i) =\phi\cdot\frac{\pi_i(1-\pi_i)}{m_i},\quad
%%%\mbox{with $\phi>0$ an over-dispersion parameter.}
%%%\]
%%
%%%Since the actual distribution for $y_i$ is difficult to get in this situation, the simplest way to account
%%%for this over-dispersion is assuming
%%%$$\mbox{var}(\bar{y}_i|\textbf{x}_i) =\phi\cdot\frac{\pi_i(1-\pi_i)}{m_i},\quad
%%%\mbox{with $\phi>0$ an over-dispersion parameter.}$$
%%%\item
%%%\textbf{Remark}. Assuming a \textbf{beta-binomial distribution} for $y_i$ can be used to account for the
%%%over-dispersion.
%%\end{itemize}
%%\end{frame}
%%
%%\begin{frame}
%%\begin{itemize}
%%\item
%%One can account for this by assuming 
%%\[\mbox{var}(\bar{y}_i|\textbf{x}_i) =\phi\cdot\frac{\pi_i(1-\pi_i)}{m_i},\quad
%%\mbox{with $\phi>0$ an over-dispersion parameter.}
%%\]
%%with the over-dispersion parameter $\phi>1$
%%
%%\item Remark: Not same as setting $\phi \not= 1$ in the exponential family form of the binomial distribution! You wouldn't get a proper density this way. However, a genuine likelihood is not actually needed for inference; more on this later when we cover \emph{quasi-likelihood}. 
%%
%%%
%%%
%%%\item However, 
%%
%%
%%\item One  may assume  a \textbf{beta-binomial distribution}, which allow for variance of the form $\phi \frac{\pi_i(1-\pi_i)}{m_i}$
%%\end{itemize}
%%
%%\end{frame}
%%
%%\subsection{\S2.1.5 GLM for count responses}\label{sec:2.1.5}
%%\begin{frame}{\S2.1.5 GLM for count responses}
%%\begin{itemize}
%%\item
%%If $y_i|\textbf{x}_i \sim \mbox{Poisson}(\lambda=\mu_i)$ with $E(y_i|\textbf{x}_i)=\mu_i$
%%and $\mbox{var}(y_i|\textbf{x}_i)=\mu_i$.
%%
%%\item The natural/canonical parameter is
%%$\theta_i=\log\mu_i$
%%
%%\item This leads to the natural link  \textbf{log-link} $\eta_i=\log\mu_i$, resulting in
%%the \textbf{log linear model}:
%%\[
%%\log\mu_i=\eta_i=\textbf{z}_i^T\!\mbox{\boldmath $\beta$} \quad \iff \quad
%%\mu_i=\exp(\textbf{z}_i^T\!\mbox{\boldmath $\beta$})=\exp(\eta_i).
%%\]
%%
%%
%%\item Overdispersion can occur for similar reasons as for binomial data; can be accounted for
%%by introducing an over-dispersion parameter $\phi$ such that $\mbox{var}(y_i|\textbf{x}_i)=\phi\mu_i$.
%%
%%\item (a topic of quasi-likelihood later)
%%
%%\item Again, some distributions for count data do exist in the exponential family that allow for over-dispersion, such as the negative binomial distribution. 
%%\end{itemize}
%%\end{frame}
%%
%%\section{\S2.2 Likelihood Inference}\label{sec:2.2}
%%\begin{frame}{\S2.2 Likelihood Inference}
%%\begin{itemize}
%%\item
%%Regression analysis with GLMs is based on \textbf{likelihood maximization}. 
%%\item
%%It includes 
%%\begin{itemize}
%%\item parameter point estimation, 
%%\item confidence intervals, 
%%\item hypothesis testing,
%%\item  goodness-of-fit tests and 
%%\item  variable selection (or model selection).
%%
%%\end{itemize}
%%\item
%%For now we assume the model is completely and correctly specified, i.e., the GLM assumption is ``true".
%%
%%\item 
%%This can  be relaxed when we cover quasi-likelihood. 
%%\end{itemize}
%%\end{frame}
%\section{\S2.2 Likelihood Inference}\label{sec:2.2}
%%\subsection{\S2.2.1 Maximum likelihood estimation}\label{sec:2.2.1}
%%\begin{frame}{\S2.2.1 Maximum likelihood estimation}
%%\begin{itemize}
%%\item
%%Data:  response $y_i$, covariates $\textbf{x}_i$ and design
%%vector $\textbf{z}_i = {\bf z}_i(\textbf{x}_i)$, $i=1,2,\cdots, n$.
%%\item
%%Model:  GLM where $E(y_i|\textbf{x}_i)=\mu_i=h(\eta_i)=h(\textbf{z}_i^T\!\mbox{\boldmath $\beta$})$,
%%with response function $h(\cdot)$ and unknown coefficient $p\times 1$ vector $\mbox{\boldmath $\beta$}$. 
%%\item
%%Assumption:  design matrix $Z = Z_{n\times p}=(\textbf{z}_1,\cdots,\textbf{z}_n)^T$ has (full) rank $p$
%%and scale/dispersion parameter $\phi$ is
%%assumed known or otherwise consistently estimated. 
%%%\item
%%%Now the log-likelihood function is {\footnotesize
%%%\begin{eqnarray*}
%%%\ell(\mbox{\boldmath $\beta$})&=&\sum_{i=1}^n\ell_i(\mbox{\boldmath $\beta$})
%%%=\sum_{i=1}^n f(y_i|\theta_i,\phi,\omega_i) \\
%%%&=&\sum_{i=1}^n \frac{y_i\theta_i-b(\theta_i)}{\phi}\cdot\omega_i
%%%+\mbox{terms not involving $\theta_i$'s} \\
%%%&=& \sum_{i=1}^n \frac{y_i\theta (\mu_i)-b(\theta(\mu_i))}{\phi}\cdot\omega_i
%%%=\sum_{i=1}^n \frac{y_i\theta (h(\textbf{z}_i^T\!\mbox{\boldmath $\beta$}))-
%%%b(\theta(h(\textbf{z}_i^T\!\mbox{\boldmath $\beta$})))}{\phi}\cdot\omega_i
%%%\end{eqnarray*}}
%%\end{itemize}
%%
%%
%%\end{frame}
%%
%%\begin{frame}{\S2.2.1 Maximum likelihood estimation}
%%
%%The focus is estimating $\mbox{\boldmath $\beta$}$. 
%%Now the log-likelihood function is {\footnotesize
%%\begin{eqnarray*}
%%\ell(\mbox{\boldmath $\beta$})&=&\sum_{i=1}^n\ell_i(\mbox{\boldmath $\beta$})
%%=\sum_{i=1}^n f(y_i|\theta_i,\phi,\omega_i) \\
%%&=&\sum_{i=1}^n \frac{y_i\theta_i-b(\theta_i)}{\phi}\cdot\omega_i
%%+\mbox{terms not involving $\theta_i$'s} \\
%%&=& \sum_{i=1}^n \frac{y_i\theta (\mu_i)-b(\theta(\mu_i))}{\phi}\cdot\omega_i
%%=\sum_{i=1}^n \frac{y_i\theta (h(\textbf{z}_i^T\!\mbox{\boldmath $\beta$}))-
%%b(\theta(h(\textbf{z}_i^T\!\mbox{\boldmath $\beta$})))}{\phi}\cdot\omega_i
%%\end{eqnarray*}}
%%
%%\end{frame}
%%
%%\begin{frame}{\S2.2.1 Maximum likelihood estimation (2)}
%%\begin{itemize}
%%\item
%%The \textbf{score function} is the first derivative: {\footnotesize
%%\begin{eqnarray*}
%%\textbf{s}(\mbox{\boldmath $\beta$})&\!\!\!=\!\!\!& \frac{\partial \ell(\mbox{\boldmath $\beta$})}{\partial\mbox{\boldmath 
%%$\beta$}}=\sum_{i=1}^n\frac{\partial \ell_i(\mbox{\boldmath $\beta$})}{\partial\mbox{\boldmath $\beta$}}
%%=\sum_{i=1}^n  \frac{\partial \ell_i(\mbox{\boldmath $\beta$})}{\partial\theta_i}\cdot
%%\frac{\partial \theta_i}{\partial\mu_i}\cdot \frac{\partial \mu_i}{\partial\eta_i} \cdot
%%\frac{\partial \eta_i}{\partial\mbox{\boldmath $\beta$}}\\
%%&\!\!\!=\!\!\!&\sum_{i=1}^n \left(\frac{y_i-b^\prime(\theta_i)}{\phi}\cdot\omega_i\right)\cdot
%%\left[v(h(\textbf{z}_i^T\!\mbox{\boldmath $\beta$}))\right]^{-1}\cdot
%% \frac{\partial h(\eta_i)}{\partial\eta_i}\cdot \textbf{z}_i \\
%%&\!\!\!=\!\!\!& \sum_{i=1}^n \textbf{z}_i\cdot \frac{\partial h(\eta_i)}{\partial\eta_i}\cdot
%%\left[\frac{\phi v(h(\textbf{z}_i^T\!\mbox{\boldmath $\beta$}))}{\omega_i}\right]^{-1}\cdot
%%\left[y_i-\mu_i(\mbox{\boldmath $\beta$})\right] \\
%%&\!\!\!\stackrel{denoted}{=} \!\!\!& \sum_{i=1}^n \textbf{z}_i D_i(\mbox{\boldmath $\beta$})
%%\sigma_i^{-2}(\mbox{\boldmath $\beta$}) \left[y_i-\mu_i(\mbox{\boldmath $\beta$})\right] \\
%%&\!\!\!=\!\!\!& \left[\textbf{z}_1, \cdots, \textbf{z}_n\right]  \left[\!\begin{array}{ccc}
%%\!\!\!D_1(\mbox{\boldmath $\beta$})\!\! & & \\
%%& \!\!\!\ddots\!\!\! & \\ & & \!\!\!D_n(\mbox{\boldmath $\beta$})\!\!\!\end{array}\!\right] 
%%\left[\begin{array}{ccc}\!\!\!\!\sigma_1^2(\mbox{\boldmath $\beta$}) \!\!\! & & \\
%%&\!\!\!\!\ddots \!\!\!\!& \\ & & \!\!\!\!\sigma_n^2(\mbox{\boldmath $\beta$})\!\!\!\end{array}\right]^{\!\!-1}
%%\left[\begin{array}{c} \!\!\!y_1\!-\!\mu_1(\mbox{\boldmath $\beta$}) \!\!\!\! \\
%%\!\!\!\vdots \!\!\!\! \\ \!\!\! y_n\!-\!\mu_n(\mbox{\boldmath $\beta$})\!\!\!\!\end{array}\!\right] \\
%%&\!\!\!\stackrel{\mbox{matrix form}}{=} \!\!\!& Z^TD(\mbox{\boldmath $\beta$})\Sigma^{-1}(\mbox{\boldmath $\beta$})
%%\left[\textbf{y}-\mbox{\boldmath $\mu$}(\mbox{\boldmath $\beta$})\right]
%%\end{eqnarray*}}
%%\end{itemize}
%%\end{frame}
%%
%%\begin{frame}{\S2.2.1 Maximum likelihood estimation (3)}
%%Supporting arguments for the derivation of $s(\mbox{\boldmath $\beta$})$
%%\begin{itemize}
%%\item
%%$\displaystyle \mu_i=b^\prime(\theta_i) \quad \Rightarrow\quad \frac{\partial \mu_i}{\partial\theta_i}=
%%b^{\prime\prime}(\theta_i)=v(\mu_i)=v(h(\textbf{z}_i^T\!\mbox{\boldmath $\beta$}))$
%%\item
%%Also $\displaystyle \frac{\partial \theta_i}{\partial\mu_i}=\left(\frac{\partial \mu_i}{\partial\theta_i}\right)^{-1}$ 
%%and $\displaystyle \frac{\partial \eta_i}{\partial\mbox{\boldmath $\beta$}}=\textbf{z}_i$.
%%\item
%%{\footnotesize Denote $\displaystyle D_i(\mbox{\boldmath $\beta$})=\frac{\partial \mu_i}{\partial\eta_i}
%%=\frac{\partial h(\eta_i)}{\partial\eta_i}\mid_{\eta_i=\textbf{z}_i^T\!\mbox{\boldmath $\beta$}}$,\;
%%$\sigma_i^2(\mbox{\boldmath $\beta$})=\omega_i^{-1}\phi v(h(\textbf{z}_i^T\!\mbox{\boldmath $\beta$}))$
%%,
%%$i=1,\cdots, n$ (or $i=1,\cdots, g$ for grouped data).}
%%\item
%%{\footnotesize Denote $\displaystyle D(\mbox{\boldmath $\beta$})=\left[\begin{array}{ccc}
%%\!\!\!D_1(\mbox{\boldmath $\beta$})\!\! & & \\
%%& \!\!\!\ddots\!\!\! & \\ & & \!\!\!D_n(\mbox{\boldmath $\beta$})\!\!\!\end{array}\right]=\mbox{diag}
%%\{D_1(\mbox{\boldmath $\beta$}), \cdots, D_n(\mbox{\boldmath $\beta$})\}$}
%%\item
%%{\footnotesize Denote $\displaystyle \Sigma(\mbox{\boldmath $\beta$})=\left[\begin{array}{ccc}
%%\!\!\!\sigma_1^2(\mbox{\boldmath $\beta$})\!\! & & \\
%%& \!\!\!\ddots\!\!\! & \\ & & \!\!\!\sigma_n^2(\mbox{\boldmath $\beta$})\!\!\!\end{array}\right]=\mbox{diag}
%%\{\sigma_1^2(\mbox{\boldmath $\beta$}), \cdots, \sigma_n^2(\mbox{\boldmath $\beta$})\}$}
%%
%%
%%\end{itemize}
%%\end{frame}
%%
%%\begin{frame}{\S2.2.1 Maximum likelihood estimation (4)}
%%
%%The \textbf{Fisher information} matrix is the covariance of the score function at the point $\mbox{\boldmath $\beta$}$: 
%%\begin{eqnarray*}
%%F(\mbox{\boldmath $\beta$})&=& \mbox{cov}\left(\textbf{s}(\mbox{\boldmath $\beta$})\right)=
%%\sum_{i=1}^n \textbf{z}_i D_i^2(\mbox{\boldmath $\beta$})\sigma_i^{-2}(\mbox{\boldmath $\beta$}) 
%%\textbf{z}_i^T=\sum_{i=1}^n \textbf{z}_i w_i(\mbox{\boldmath $\beta$})\textbf{z}_i^T \\
%%&=& \left[\textbf{z}_1, \cdots, \textbf{z}_n\right]  \left[\begin{array}{ccc}
%%\!\!\!w_1(\mbox{\boldmath $\beta$})\!\! & & \\
%%& \!\!\!\ddots\!\!\! & \\ & & \!\!\!w_n(\mbox{\boldmath $\beta$})\!\!\!\end{array}\right] 
%%\left[\begin{array}{c} \textbf{z}_1^T \\ \vdots \\ \textbf{z}_n^T\end{array}\right]
%%=Z^TW(\mbox{\boldmath $\beta$})Z,
%%\end{eqnarray*}
%% where $w_i(\mbox{\boldmath $\beta$}):=D_i^2(\mbox{\boldmath $\beta$})\sigma_i^{-2}(\mbox{\boldmath $\beta$})$ and 
%%
%%\[ W(\mbox{\boldmath $\beta$})=\left[\begin{array}{ccc}
%%\!\!\!w_1(\mbox{\boldmath $\beta$})\!\! & & \\
%%& \!\!\!\ddots\!\!\! & \\ & & \!\!\!w_n(\mbox{\boldmath $\beta$})\!\!\!\end{array}\right]=\mbox{diag}
%%\{w_1(\mbox{\boldmath $\beta$}), \cdots, w_n(\mbox{\boldmath $\beta$})\}\]
%%
%%\end{frame}
%%
%%\begin{frame}{\S2.2.1 Maximum likelihood estimation (5)}
%%\begin{itemize}
%%\item
%%The \textbf{observed information} matrix with respect to $\mbox{\boldmath $\beta$}$ is {\footnotesize
%%\begin{eqnarray*}
%%F_{\mbox{obs}}(\mbox{\boldmath $\beta$})& \!\!\!= \!\!\!& -\frac{\partial^2\ell(\mbox{\boldmath $\beta$})}{
%%\partial \mbox{\boldmath $\beta$}\partial \mbox{\boldmath $\beta$}^T}=
%%-\frac{\partial \textbf{s}(\mbox{\boldmath $\beta$})}{\partial \mbox{\boldmath $\beta$}^T} \\
%%& \!\!\!= \!\!\!& -\sum_{i=1}^n \textbf{z}_i \frac{\partial \{D_i(\mbox{\boldmath $\beta$})\sigma_i^{-2}(
%%\mbox{\boldmath $\beta$})\} }{\partial \mbox{\boldmath $\beta$}^T}[y_i-\mu_i(\mbox{\boldmath $\beta$})]
%%+\sum_{i=1}^n \textbf{z}_i D_i(\mbox{\boldmath $\beta$})\sigma_i^{-2}(\mbox{\boldmath $\beta$}) \cdot
%%\frac{\partial\mu_i}{\partial \eta_i}\cdot \frac{\partial\eta_i}{\partial \mbox{\boldmath $\beta$}^T} \\
%%& \!\!\!= \!\!\!& -\sum_{i=1}^n \textbf{z}_i \frac{\partial}{\partial \mbox{\boldmath $\beta$}^T}
%%\{D_i(\mbox{\boldmath $\beta$})\sigma_i^{-2}(\mbox{\boldmath $\beta$})\}\cdot
%%[y_i-\mu_i(\mbox{\boldmath $\beta$})] +\sum_{i=1}^n \textbf{z}_i D_i^2(\mbox{\boldmath $\beta$})
%%\sigma_i^{-2}(\mbox{\boldmath $\beta$}) \cdot \textbf{z}_i^T
%%\end{eqnarray*}}
%%\item
%%Hence $E\left[F_{\mbox{obs}}(\mbox{\boldmath $\beta$})\right]=\sum_{i=1}^n \textbf{z}_i 
%%D_i^2(\mbox{\boldmath $\beta$})\sigma_i^{-2}(\mbox{\boldmath $\beta$}) \textbf{z}_i^T
%%=F(\mbox{\boldmath $\beta$})=\mbox{cov}\left(\textbf{s}(\mbox{\boldmath $\beta$})\right)$.
%%\end{itemize}
%%\end{frame}
%%
%%\begin{frame}{\S2.2.1 Maximum likelihood estimation (6)}
%%\begin{itemize}
%%\item
%%When natural link is used, $\textbf{s}(\mbox{\boldmath $\beta$})$ and $F(\mbox{\boldmath $\beta$})$ can be
%%simplified to
%%\begin{eqnarray*}
%%\textbf{s}(\mbox{\boldmath $\beta$})&= & \frac{1}{\phi}\sum_{i=1}^n \omega_i\textbf{z}_i
%%[y_i-\mu_i(\mbox{\boldmath $\beta$})]=\frac{1}{\phi}Z^T\Omega \left[\textbf{y}-\mbox{\boldmath $\mu$}
%%(\mbox{\boldmath $\beta$})\right] \\
%%F(\mbox{\boldmath $\beta$}) &= & \frac{1}{\phi}\sum_{i=1}^n \omega_iv(\mu_i(\mbox{\boldmath $\beta$}))
%%\textbf{z}_i\textbf{z}_i^T=\frac{1}{\phi}Z^T\Omega V(\mbox{\boldmath $\beta$})Z
%%\end{eqnarray*}
%%where $\Omega=\mbox{diag}\{\omega_1,\cdots,\omega_n\}$ and $V(\mbox{\boldmath $\beta$})
%%=\mbox{diag}\{v(\mu_1),\cdots, v(\mu_n)\}$.
%%\item
%%For natural links, $\theta_i=\eta_i=\textbf{z}_i^T\mbox{\boldmath $\beta$} \quad \Rightarrow\quad
%%\mu_i=b^\prime (\theta_i)=b^\prime(\textbf{z}_i^T\mbox{\boldmath $\beta$})\equiv 
%%h(\textbf{z}_i^T\mbox{\boldmath $\beta$})$
%%
%%$\Rightarrow \quad D_i(\mbox{\boldmath $\beta$})=\frac{\partial h(\eta_i)}{\partial \eta_i}
%%=b^{\prime\prime}(\textbf{z}_i^T\mbox{\boldmath $\beta$})=v(h(\textbf{z}_i^T\mbox{\boldmath $\beta$}))$
%%
%%$\Rightarrow \quad D_i(\mbox{\boldmath $\beta$})\sigma_i^{-2}(\mbox{\boldmath $\beta$})
%%=v(h(\textbf{z}_i^T\mbox{\boldmath $\beta$}))\cdot\frac{\omega_i}{\phi
%%v(h(\textbf{z}_i^T\mbox{\boldmath $\beta$}))}=\frac{\omega_i}{\phi}$ and
%%
%%$\Rightarrow \quad w_i(\mbox{\boldmath $\beta$})=D_i^2(\mbox{\boldmath $\beta$})
%%\sigma_i^{-2}(\mbox{\boldmath $\beta$})=D_i(\mbox{\boldmath $\beta$})\cdot \frac{\omega_i}{\phi}
%%=\frac{1}{\phi}\omega_i v(h(\textbf{z}_i^T\mbox{\boldmath $\beta$}))$.
%%\end{itemize}
%%\end{frame}
%%
%%
%%\begin{frame}
%%We'll continue our discussion on MLE estimate for GLMs. 
%%
%%\end{frame}
%%
%%
%%\begin{frame}{\S2.2.1 Numerical computation of the MLE $\hat{\mbox{\boldmath $\beta$}}$ by IRLS}
%%\begin{itemize}
%%\item
%%MLE $\hat{\mbox{\boldmath $\beta$}}$ is the one maximizing $\ell(\mbox{\boldmath $\beta$})$, and often 
%%is obtained by solving the \textbf{estimating equation} $\textbf{s}(\mbox{\boldmath $\beta$})=0$, which uses 
%%the iterative \textbf{Fisher scoring method}: {\scriptsize
%%\begin{eqnarray*}
%%\hat{\mbox{\boldmath $\beta$}}^{(k+1)}&\!\!\!\!\!\!=\!\!\!\!\!\! & \hat{\mbox{\boldmath $\beta$}}^{(k)}+
%%F^{-1}(\hat{\mbox{\boldmath $\beta$}}^{(k)})\cdot \textbf{s}(\hat{\mbox{\boldmath $\beta$}}^{(k)}) \\
%%&\!\!\!\!\!\!=\!\!\!\!\!\! & \hat{\mbox{\boldmath $\beta$}}^{(k)}+ 
%%\left[Z^TW(\hat{\mbox{\boldmath $\beta$}}^{(k)})Z\right]^{-1}
%% Z^TD(\hat{\mbox{\boldmath $\beta$}}^{(k)})\Sigma^{-1}(\hat{\mbox{\boldmath $\beta$}}^{(k)})
%%\left[\textbf{y}-\mbox{\boldmath $\mu$}(\hat{\mbox{\boldmath $\beta$}}^{(k)})\right] \\
%%&\!\!\!\!\!\!=\!\!\!\!\!\! & \hat{\mbox{\boldmath $\beta$}}^{(k)}+ 
%%\left[Z^TW(\hat{\mbox{\boldmath $\beta$}}^{(k)})Z\right]^{-1}
%% Z^TD^2(\hat{\mbox{\boldmath $\beta$}}^{(k)})\Sigma^{-1}(\hat{\mbox{\boldmath $\beta$}}^{(k)})
%%D^{-1}(\hat{\mbox{\boldmath $\beta$}}^{(k)}) 
%%\left[\textbf{y}-\mbox{\boldmath $\mu$}(\hat{\mbox{\boldmath $\beta$}}^{(k)})\right] \\
%%&\!\!\!\!\!\!=\!\!\!\!\!\! & \left[Z^TW(\hat{\mbox{\boldmath $\beta$}}^{(k)})Z\right]^{-1}\left\{
%%Z^TW(\hat{\mbox{\boldmath $\beta$}}^{(k)})Z \hat{\mbox{\boldmath $\beta$}}^{(k)}
%%+Z^TW(\hat{\mbox{\boldmath $\beta$}}^{(k)}) D^{-1}(\hat{\mbox{\boldmath $\beta$}}^{(k)})
%%\left[\textbf{y}-\mbox{\boldmath $\mu$}(\hat{\mbox{\boldmath $\beta$}}^{(k)})\right]\!\right\} \\
%%&\!\!\!\!\!\!=\!\!\!\!\!\!& \left[Z^TW(\hat{\mbox{\boldmath $\beta$}}^{(k)})Z\right]^{-1}\left\{
%%Z^TW(\hat{\mbox{\boldmath $\beta$}}^{(k)})\left[Z \hat{\mbox{\boldmath $\beta$}}^{(k)}
%%+ D^{-1}(\hat{\mbox{\boldmath $\beta$}}^{(k)})
%%\left[\textbf{y}-\mbox{\boldmath $\mu$}(\hat{\mbox{\boldmath $\beta$}}^{(k)})\right]\right]\right\} \\
%%&\!\!\!\!\!\!\stackrel{denoted}{=}\!\!\!\!\!\! & \left[Z^TW(\hat{\mbox{\boldmath $\beta$}}^{(k)})Z\right]^{-1}
%%Z^TW(\hat{\mbox{\boldmath $\beta$}}^{(k)})\cdot \tilde{\textbf{y}}^{(k)}
%%\end{eqnarray*}}
%%where $\tilde{\textbf{y}}^{(k)}=(\tilde{y}_1^{(k)},\cdots, \tilde{y}_n^{(k)})^T
%%=Z \hat{\mbox{\boldmath $\beta$}}^{(k)}+ D^{-1}(\hat{\mbox{\boldmath $\beta$}}^{(k)})
%%\left[\textbf{y}-\mbox{\boldmath $\mu$}(\hat{\mbox{\boldmath $\beta$}}^{(k)})\right]$
%%is the \textbf{working observation vector} at iteration $k$, $k=0,1,2,\cdots$.
%%\end{itemize}
%%\end{frame}
%%
%%
%%\begin{frame}{\S2.2.1 Implementing Fisher scoring for computing $\hat{\mbox{\boldmath $\beta$}}$}
%%\begin{enumerate}
%%\item
%%Start with an initial estimate $\hat{\mbox{\boldmath $\beta$}}^{(0)}$, often chosen as the unweighted
%%least squares (LS) estimate from $(g(y_i),\textbf{z}_i)$, $i=1,\cdots,n$, where $g(\cdot)$ is the link (modified if $g(y_i)$ is undefined, such as when $yi = 0$ in Poisson regression).
%%\item
%%For $k=0,1,2,\cdots$ compute 
%%$$\hat{\mbox{\boldmath $\beta$}}^{(k+1)}=\left[Z^TW(\hat{\mbox{\boldmath $\beta$}}^{(k)})Z\right]^{-1}
%%Z^TW(\hat{\mbox{\boldmath $\beta$}}^{(k)})\cdot \tilde{\textbf{y}}^{(k)}$$
%%until $\frac{\|\hat{\boldsymbol \beta}^{(k+1 )} - \hat{\boldsymbol \beta}^{(k)}\|}{\|\hat{\boldsymbol \beta}^{(k)}\|} \leq \epsilon $ for some prechosen small $\epsilon > 0$ (``convergence").
%%
%%\item It is hence called a  \textbf{iteratively reweighted least squares (IRLS) method}. 
%%\end{enumerate}
%%%Fisher scoring method in GLM estimation is also called the \textbf{iteratively reweighted least squares (IRLS) method}.
%%\end{frame}
%%
%%
%%\begin{frame}
%%
%%\begin{itemize}
%%
%%\item 
%%Each component of the working observation vector is simply a linearization of $g(y_i)$ at $\hat{\boldsymbol \beta}^{(k)}$
%%
%%\begin{align*}
%%g(y_i) &\approx g(\mu) +g'(\mu) (y_i - \mu) \\
%%&= {\bf z}_i^T{\boldsymbol \beta} +  \left(\frac{\partial h}{ \partial \eta} ({\bf z}_i^T{\boldsymbol \beta} )\right)^{-1}  (y_i - \mu_i ({\boldsymbol \beta}))
%%\end{align*}
%%(Generalized Linear Models, 2nd ed, McCullagh and Nelder, p.40)
%%
%%\item If use $F_{obs} (\cdot)$ instead of $F (\cdot)$, we get the standard Newton-Raphson procedure. ( These two procedures coincide if natural link function is used)
%%\end{itemize}
%%\end{frame}
%%
%%\begin{frame}{\S2.2.1 Asymptotic properties associated with the MLE $\hat{\mbox{\boldmath $\beta$}}$}
%%Under suitable ``regularity" assumptions, we have these nice properties:
%%\begin{enumerate}
%%\item
%%\textit{Asymptotic existence and uniqueless}. The MLE $\hat{\mbox{\boldmath $\beta$}}$ exists
%%and (locally) unique w/ probability tending to 1. 
%%\item
%%\textit{Consistency}. If $\mbox{\boldmath $\beta$}$ is the ``true'' value, then 
%%$\hat{\mbox{\boldmath $\beta$}}\rightarrow \mbox{\boldmath $\beta$}$ in probability 1 (\textit{weak consistency}),
%%or $\hat{\mbox{\boldmath $\beta$}}\rightarrow \mbox{\boldmath $\beta$}$ with probability 1 (or almost surely,
%%\textit{strong consistency}).
%%
%%%\item
%%%Estimating $\phi$: 
%%%$\displaystyle \hat{\phi}=\frac{1}{g-p}\sum_{i=1}^g\frac{(\bar{y}_i-\hat{\mu}_i)^2}{v(\hat{\mu}_i)/n_i}$
%%%is a consistent estimator of $\phi$. (This is true for grouped data only.)
%%\end{enumerate}
%%\end{frame}
%%
%%
%%\begin{frame}
%%{\S2.2.1 Asymptotic properties associated with the MLE $\hat{\mbox{\boldmath $\beta$}}$}
%%
%%\begin{enumerate}[3]
%%\item
%%\textit{Asymptotic normality}. $\hat{\mbox{\boldmath $\beta$}}\stackrel{a}{\sim} \mbox{Normal}
%%\left(\mbox{\boldmath $\beta$}, F^{-1}(\hat{\mbox{\boldmath $\beta$}})\right)$, equivalently
%%$$F^{1/2}(\hat{\mbox{\boldmath $\beta$}})\left(\hat{\mbox{\boldmath $\beta$}}-
%%\mbox{\boldmath $\beta$}\right)\stackrel{d}{\rightarrow} N(0,1) \quad\mbox{as $n\rightarrow\infty$}.$$
%%This implies $\mbox{cov}(\hat{\mbox{\boldmath $\beta$}})\stackrel{a}{=}A(\hat{\mbox{\boldmath $\beta$}})
%%= F^{-1}(\hat{\mbox{\boldmath $\beta$}})$. (The proof uses the property $\textbf{s}(\mbox{\boldmath $\beta$})
%%\stackrel{a}{\sim} \mbox{Normal}\left(0, F(\mbox{\boldmath $\beta$})\right)$.)
%%\end{enumerate}
%% \ \
%% 
%%(Now recall in the linear model case, the Fisher information is $\sigma^{-2} {\bf Z}^T {\bf Z}$)
%%
%%\end{frame}
%%
%%\begin{frame} {\S2.2.1 Estimation of $\phi$}
%%\begin{itemize}
%%\item We also note that $\phi$ only appears as a factor in $W$, and falls out of the expression
%%
%% \[\hat{\mbox{\boldmath $\beta$}}^{(k+1)}=\left[Z^TW(\hat{\mbox{\boldmath $\beta$}}^{(k)})Z\right]^{-1}
%%Z^TW(\hat{\mbox{\boldmath $\beta$}}^{(k)})\cdot \tilde{\textbf{y}}^{(k)}\]
%%\item But we need $\phi$ for computing the Fisher info for the limiting covariance of the MLE. 
%%
%%\item\[\hat{\phi}=\frac{1}{g-p}\sum_{i=1}^g\frac{(\bar{y}_i-\hat{\mu}_i)^2}{v(\hat{\mu}_i)/n_i}\]
%%is a consistent estimator of $\phi$. (This is true for grouped data only.)
%%\end{itemize}
%%\end{frame}
%%
%%\subsection{\S2.2.2 Linear hypothesis testing in GLM}
%%
%%\begin{frame}{\S2.2.2 Testing linear hypotheses}
%%\begin{itemize}
%%\item
%%Most testing problems concerning $\mbox{\boldmath $\beta$}$ have the form of \textit{linear hypothesis}:
%%$$H_0: C\mbox{\boldmath $\beta$}=\mbox{\boldmath $\xi$} \quad\mbox{vs.}\quad
%%H_1: C\mbox{\boldmath $\beta$}\neq\mbox{\boldmath $\xi$}$$
%%where $C$ is a known matrix of full row rank $s\leq p$, and $\mbox{\boldmath $\xi$}$ is a known vector.
%%\item
%%A special case is 
%%$$H_0: \mbox{\boldmath $\beta$}_r=0 \quad\mbox{vs.}\quad
%%H_1: \mbox{\boldmath $\beta$}_r\neq 0$$
%%where $\mbox{\boldmath $\beta$}_r$ is a sub-vector of $\mbox{\boldmath $\beta$}$. This is
%%equivalent to comparing the full model with one of its sub-models.
%%\end{itemize}
%%\end{frame}
%%
%%\begin{frame}{\S2.2.2 Likelihood ratio (LR), Wald and score tests (1)}
%%\begin{itemize}
%%\item
%%Three types of tests are available for testing $H_0: C\mbox{\boldmath $\beta$}=\mbox{\boldmath $\xi$}
%%\;\mbox{vs.}\; H_1: C\mbox{\boldmath $\beta$}\neq\mbox{\boldmath $\xi$}$:
%%
%%\begin{center}
%%\textbf{Likelihood ratio (LR) test}, \textbf{Wald test}, and \textbf{score test}.
%%\end{center}
%%\item
%%A special case is 
%%$$H_0: \mbox{\boldmath $\beta$}_r=0 \quad\mbox{vs.}\quad
%%H_1: \mbox{\boldmath $\beta$}_r\neq 0$$
%%where $\mbox{\boldmath $\beta$}_r$ is a sub-vector of $\mbox{\boldmath $\beta$}$. This is
%%equivalent to comparing the full model with one of its sub-models.
%%\end{itemize}
%%\end{frame}
%%
%%\begin{frame}{\S2.2.2 Likelihood ratio (LR), Wald and score tests (2)}
%%\begin{enumerate}
%%\item
%%\textbf{LR test statistic}:
%% \[\lambda=-2\left\{\ell(\tilde{\mbox{\boldmath $\beta$}}) -\ell(\hat{
%%\mbox{\boldmath $\beta$}})\right\},\] 
%%where $\tilde{\mbox{\boldmath $\beta$}}$ is the MLE of
%%$\mbox{\boldmath $\beta$}$ under $H_0$, and $\hat{\mbox{\boldmath $\beta$}}$ is the MLE of
%%$\mbox{\boldmath $\beta$}$ under $H_1$; also $\hat{\mbox{\boldmath $\beta$}}$ is the unrestricted MLE of
%%$\mbox{\boldmath $\beta$}$. 
%%
%%%
%%%$\lambda$ needs to be divided by $\hat{\phi}$ for over-dispersion models.
%%\item
%%\textbf{Wald test statistic}: \[W=\left(C\hat{\mbox{\boldmath $\beta$}}-\mbox{\boldmath $\xi$}\right)^T
%%\left[CF^{-1}(\hat{\mbox{\boldmath $\beta$}})C^T \right]^{-1}
%%\left(C\hat{\mbox{\boldmath $\beta$}}-\mbox{\boldmath $\xi$}\right),\] where $\hat{\mbox{\boldmath $\beta$}}$
%%is the unrestricted MLE of $\mbox{\boldmath $\beta$}$.
%%\item
%%\textbf{Score test statistic}: $U=\textbf{s}^T(\tilde{\mbox{\boldmath $\beta$}})F^{-1}(\tilde{\mbox{\boldmath $\beta$}})
%%\textbf{s}(\tilde{\mbox{\boldmath $\beta$}})$, where $\tilde{\mbox{\boldmath $\beta$}}$
%%is the MLE of $\mbox{\boldmath $\beta$}$ under $H_0$.
%%\end{enumerate}
%%%\begin{itemize}
%%%
%%%
%%%\end{itemize}
%%\end{frame}
%%
%%\begin{frame}{\S2.2.2 Likelihood ratio (LR), Wald and score tests (3)}
%%\begin{itemize}
%%\item  The  exponential model is assumed for the LR test; not defined for over-dispersion model such as the quasi likelihood to be covered. 
%%
%%\item 
%%Wald and score test statistics are properly defined for over-dispersion models, since they are based on the score ${\bf s}$ (More on this later again)
%%
%%\item Wald and score are quadratic approximation of the LR test (see blackboard derivation). 
%%
%%\item Wald and Score are computationally attractive; both eliminate the need to compute one of $\tilde{\boldsymbol \beta}$ or $\hat{\boldsymbol \beta}$. Wald is good when $H_1$ has been fitted, score is good when $H_0$ has been fitted. 
%%
%%
%%\end{itemize}
%%
%%\end{frame}
%%
%%\begin{frame} {\S2.2.2 Likelihood ratio (LR), Wald and score tests (4)}
%%\begin{itemize}
%%\item
%%Under \[H_0: C{\boldmath \beta}={\boldmath \xi} \text{ with } \text{rank}(C)=s, \] $\lambda, W$ and
%%$U$ all follow a $\chi^2(s)$ distribution asymptotically under `regularity conditions'.
%%%\item
%%%$\lambda/\hat{\phi}$, $W$ and $U$ can be used for \textbf{analysis of deviance} (i.e. comparing nested
%%%models) as a special case of testing linear hypotheses.
%%
%%\end{itemize}
%%
%%
%%\end{frame}
%
%\subsection{\S2.2.3 Deviance, residuals and goodness-of-fit measure in GLM} \label{sec:2.2.3}
%
%\begin{frame} {Model adequacy measure in GLM}
%\begin{itemize}
%\item First of all, testing model fit only makes sense when over-dispersion is not assumed, since we don't assume the exponential model form is true anyway. Again a topic of quasi-likelihood covered later.   
%
%\item In the sequel, ``goodness-of-fit" and ``model adequacy" mean the same thing.
%\end{itemize}
%\end{frame}
%
%\begin{frame}{\S2.2.3 Model adequacy measure in GLM }
%\begin{itemize}
%\item
%In linear model, goodness-of-fit  is measured by the sum of squared residuals (or errors)
%$\mbox{SSE}\!=\!\sum_{i=1}^n(y_i-\hat{\mu}_i)^2$. SSE is no longer a reasonable measure if $\mbox{Var}(y)$ 
%varies with $\mu$.
%\item
%A generalisation of SSE to GLM  is the \textbf{Pearson $\chi^2$ statistic}, taking the form
%$$\chi^2=\sum_{i=1}^g \frac{(\bar{y}_i-\hat{\mu}_i)^2}{v(\hat{\mu}_i)/n_i}$$
%where $\bar{y}_i$ is the group mean, $v(\hat{\mu}_i)$ is the estimated variance function, and $n_i$ is the size of group $i$.
%\end{itemize}
%\end{frame}
%
%
%
%\begin{frame}{\S2.2.3 Model adequacy measure in GLM }
%\begin{itemize}
%%\item
%%Pearson $\chi^2$ statistic is often an appropriate model adequacy measure (or discrepancy measure, or goodness-of-fit
%%measure).
%\item
%Asymptotically (i.e. all $n_i$'s are large),
% \[\chi^2/\phi \sim_a \chi^2(g-p)\] if the GLM model is true. $p = $  number of coefficient parameters in the GLM.
%\item
%If $n$ is large but $n_i$'s are small (in particular when $n_i=1$), it is dangerous to use the Pearson
%$\chi^2$ statistics to test model adequacy.
%
%\end{itemize}
%\end{frame}
%
%\begin{frame}{\S2.2.3 Model adequacy measure in GLM}
%\begin{itemize}
%\item
%Another model adequacy measure is the {\bf deviance}. 
%\item
%It is twice the difference between the maximum log-likelihood achievable in a saturated model and that in
%the current model, i.e.
%$$D^*(\textbf{y}_n;\mbox{\boldmath $\hat{\mu}$})=2\left[\ln L(\mbox{\boldmath $\theta$}^*,\phi\mid\textbf{y}_n)
%-\ln L(\mbox{\boldmath $\hat{\theta}$},\phi\mid\textbf{y}_n)\right]$$
%which is called the \textbf{scaled deviance}. Just a special case the likelihood ratio statistic. 
%
%% Here {\small \vspace{-3pt}
%%\begin{equation*}
%%L(\mbox{\boldmath $\theta$},\phi\mid\textbf{y}_n)=\exp\left\{\sum_{i=1}^n\frac{y_i\theta_i-b(\theta_i)}{\phi}\cdot
%%\omega_i +\sum_{i=1}^nc(y_i,\phi, \omega_i)\right\} \vspace{-3pt}
%%\label{eq-8.4}
%%\end{equation*}}
%%is the likelihood function for the data $\textbf{y}_n=(y_1,\cdots,y_n)^T$.
%\item Remark: we are having grouped data in mind here
%
%\item
%%$\mbox{\boldmath $\theta$}\!=\!(\theta_1,\cdots,\theta_n)^T$, $\mbox{\boldmath $\mu$}\!=\!(\mu_1,\cdots,\mu_n)^T
%%\!=\!(b^\prime(\theta_1),\cdots, b^\prime(\theta_n))^T$,
%
% $\mbox{\boldmath $\hat{\theta}$}$ is the MLE of
%$\mbox{\boldmath $\theta$}$ under the GLM
%\item $\mbox{\boldmath $\theta$}^*=(\theta_1^*,\cdots, \theta_g^*)^T$ is such that 
%\[
%(\bar{y}_1,\cdots,\bar{y}_g)^T
%=(b^\prime(\theta_1^*),\cdots,b^\prime(\theta_g^*))^T = ({\mu}^*_1, \dots, {\mu}^*_g).\]
%\end{itemize}
%\end{frame}
%
%\begin{frame}{\S2.2.3 Deviance measure in GLM (2)}
%\begin{itemize}
%\item
%We call $\mbox{\boldmath $\theta$}^*$ the MLE of $\mbox{\boldmath $\theta$}$ in the \textbf{saturated model}, a full model with $g$ parameters (as many parameters as there are observsations).
%\item
%In a saturated model, from the likelihood expression we see that in this case, the MLE of 
%$\theta_i$ needs to satisfy $\bar{y}_i=b^\prime(\theta_i^*) = {\mu}^*_i$.
%\item
%In GLM, $\theta_1,\cdots,\theta_g$ are determined by the regression parameters $\beta_1,\cdots,\beta_p$
%and the independent variables $x_1,\cdots,x_p$. So the MLE of $\beta_1,\cdots,\beta_p$ are to be obtained first, 
%the MLE $\mbox{\boldmath $\hat{\theta}$}$ are then obtained.
%\item
%%The saturated model provides the best fit to the data because it has the same number of independent parameters
%%and the observations.
%Maximum likelihood from the saturated model must be at least as large as that from a usual GLM.
%Therefore, the scaled deviance $D^*(\textbf{y}_n;\mbox{\boldmath $\hat{\mu}$})\geq 0$. 
%%and measures
%%the adequacy of fit of a GLM.
%\end{itemize}
%\end{frame}
%
%\begin{frame}{\S2.2.3 Deviance measure in GLM (3)}
%\begin{itemize}
%\item
%%Consider a random sample $\textbf{y}_n=(y_1,\cdots,y_n)^T$ and a GLM as being formulated so far.
%By convention
%$$D(\textbf{y}_n;\mbox{\boldmath $\hat{\mu}$}):=\phi D^*(\textbf{y}_n;\mbox{\boldmath $\hat{\mu}$})$$
%%and call $D(\textbf{y}_n;\mbox{\boldmath $\hat{\mu}$})$ 
%as the \textbf{deviance} of the GLM under consideration.
%\item
%Asymptotically, $$\phi^{-1}D(\textbf{y}_n;\mbox{\boldmath $\hat{\mu}$})\sim_a\chi^2(g-p)$$ 
%asymptotically when the GLM  for $Y$ is correct  + when group sizes $n_i$ are all large enough + ``regularity conditions". 
%\item
%But when $n_i$ is small (e.g. $n_i = 1$), use of the chi-square limit is dangerous. 
%\end{itemize}
%\end{frame}
%
%\begin{frame}{\S2.2.3 Deviances for several important distributions}
%\textbf{Deviances for several important distributions in GLM.}
%\begin{eqnarray*}
%\mbox{Normal}\!\!\! & N(\mu,\sigma^2): & D(\textbf{y}_n;\mbox{\boldmath $\hat{\mu}$})=
%\sum_{i=1}^n (y_i-\hat{\mu}_i)^2 \\
%\mbox{Poisson}\!\!\! &\mbox{Poi}(\mu): &  D(\textbf{y}_n;\mbox{\boldmath $\hat{\mu}$})=2\sum_{i=1}^n
%\left\{y_i\ln\frac{y_i}{\hat{\mu}_i}-(y_i-\hat{\mu}_i)\right\} \\
%\mbox{binomial}\!\!\! & \mbox{Bin}(m,\mu): & D(\textbf{y}_n;\mbox{\boldmath $\hat{\mu}$})=2\sum_{i=1}^n
%\left\{y_i\ln\frac{y_i}{\hat{\mu}_i}\!+\!(\!m_i\!-\!y_i\!)\ln\frac{m_i\!-\!y_i}{m_i\!-\!\hat{\mu}_i}\right\} \\
%\mbox{gamma} \!\!\!& \mbox{Gam}(\mu,\nu): & D(\textbf{y}_n;\mbox{\boldmath $\hat{\mu}$})=2\sum_{i=1}^n
%\left\{-\ln\frac{y_i}{\hat{\mu}_i}+\frac{y_i-\hat{\mu}_i}{\hat{\mu}_i}\right\}
%\end{eqnarray*}
%\end{frame}
%
%\begin{frame}{\S2.2.3 Analysis of deviance (1)}
%\begin{itemize}
%\item
%An analysis of variance (ANOVA) table can be constructed for each linear regression model as we have seen
%in a linear model course. 
%\item
%Similar to this, an \textbf{analysis of deviance (ANODEV)} table can be constructed for each GLM.
%\item
%Generally speaking, tests performed on ANOVA and ANODEV tables are likelihood ratio tests.
%\item
%However, there are some differences in interpreting ANOVA and ADODEV tables.
%\item
%Consider a $3\times 3$ factorial design, of two treatment factors A and B, with 2 replicates at each combination
%level of A and B. So there will be in total 18 observations for the response variable $y$. ($3\times 3$ factorial design means that each of treatment $A$ or $B$ has 3 levels)
%\end{itemize}
%\end{frame}
%
%\begin{frame}{\S2.2.3 Analysis of deviance (2)}
%If it is appropriate to fit the 18 $y$ observations by linear regression models, an ANOVA table 
%would typically look like the following:
%
%\begin{center} Analysis of Variance (ANOVA)\end{center}
%
%{\small
%\begin{tabular}{|cccccc|}\hline
%Model & df & SSE & $\Delta$SS & $\Delta$df & Effect Source \\ \hline
%1 & 17 & 1000 & - & -  & - \\ \hline
%A & 15 & 400 & $600=1000-400$ & $2=17-15$ & A \\ \hline
%A+B & 13 & 300 & $100=400-300$ & $2=15-13$ & B \\ \hline
%A+B+A:B & 9 & 100 & $200=300-100$ & $4=13-9$ & A:B \\ \hline 
%\end{tabular}}
%\begin{itemize}
%\item
%$\Delta$SS is the reduction of SSE as an additional term is added to the model.
%$\Delta\mbox{SS}\stackrel{d}{=}\chi^2(\Delta\mbox{df})$. Effects are assessed by
%$F$ tests.
%\item
%$\Delta$df is the reduction of the model degrees of freedom.
%\item
%The effects of A, B and A:B on $Y$ are orthogonal.
%\end{itemize}
%\end{frame}
%
%\begin{frame}{\S2.2.3 Analysis of deviance (3)}
%If it is more appropriate to fit the 18 $y$ observations by GLM, an ANODEV table would typically look like the following:
%
%\begin{center} Analysis of Deviance (ANODEV)\end{center}
%
%{\footnotesize
%\begin{tabular}{|cccccc|}\hline
%Model & df & Deviance & $\Delta D$ & $\Delta$df & Effect Source \\ \hline
%1 & 17 & 1200 & - & - & - \\ \hline
%A & 15 & 800 & $400\!=\!1200\!-\!800$ & $2\!=\!17\!-\!15$ & A ignoring B \\ \hline
%A+B & 13 & 500 & $300\!=\!800\!-\!500$ & $2\!=\!15\!-\!13$ & B eliminating A\\ \hline
%A+B+A:B & 9 & 200 & $300\!=\!500\!-\!200$ & $4\!=\!13\!-\!9$ & A:B eliminating \\ 
%& & & &  & A \& B \\ \hline 
%\end{tabular}}
%\begin{itemize}
%\item
%$\Delta D=D_0-D_1$ is the reduction of deviance between two nested models $M_0\subseteq M_1$.
%$\Delta D\stackrel{d}{\approx}\chi^2(\Delta\mbox{df})$ if $M_0$ is adequate.
%\item
%$\Delta$df is the reduction of the model degrees of freedom.
%\item
%Effects of A, B and A:B on $Y$ are not orthogonal, depend on their positions in the model, and
%are assessed by $\chi^2$ or $F$ tests.
%\end{itemize}
%\end{frame}
%
%%\section{\S8.4 Residuals}
%%\label{sec:residuals}
%
%\begin{frame}{\S2.2.3 GLM residuals (1)}
%\begin{itemize}
%\item
%In a linear regression model $y=E(y)+\varepsilon =\mu+\varepsilon=\sum_{j=1}^p x_j\beta_j+\varepsilon$,
%the residual is defined as $r_i=y_i-\hat{\mu}_i$ for an observation $y_i$, with $\hat{\mu}_i=\hat{y}_i$ being the 
%fitted value of $y_i$.
%The residual $r_i$ measures the goodness of fit to $y_i$ by the model if $\mbox{Var}(y)$ is constant.
%\item
%In GLM the residual of $y_i$ has to be defined differently to accommodate the cases
%where the variance function $v(\mu)$ is no longer constant.
%\item
%At least four types of residual are defined for GLM: \textbf{Pearson}, \textbf{deviance}, \textbf{working}, and
%\textbf{response}. All these residuals can be computed using the \texttt{R} function \texttt{residuals.glm()}.
%\end{itemize}
%\end{frame}
%
%\begin{frame}{\S2.2.3 GLM residuals (2)}
%\begin{itemize}
%\item
%\textbf{Pearson residual}: $r_{_{Pi}}=\frac{\bar{y}_i-\hat{\mu}_i}{\sqrt{v(\hat{\mu})/n_i}}$.
%
%Therefore $\sum_{i=1}^g r^2_{_{Pi}}=\sum_{i=1}^g \frac{(y_i-\hat{\mu}_i)^2}{v(\hat{\mu}_i)/n_i}=\chi^2$, the
%Pearson $\chi^2$ statistic. 
%\item
%\textbf{Deviance residual}: $r_{_{Di}}=\mbox{sign}(y_i-\hat{\mu}_i)\sqrt{d_i}$, \\ where $d_i$'s are such that
%$\displaystyle \sum_{i=1}^n d_i=D(\textbf{y}_n;\mbox{\boldmath $\hat{\mu}$})$. In particular, deviance residuals for
%the following distributions are:
%{\footnotesize
%\begin{eqnarray*}
%\mbox{Normal}\!\!\! & N(\mu,\sigma^2): & r_{_{Di}}\!=\!\mbox{sign}(y_i-\hat{\mu}_i)
%|y_i-\hat{\mu}_i|\!=\!y_i-\hat{\mu}_i \\
%\mbox{Poisson}\!\!\! &\mbox{Poi}(\mu): &  r_{_{Di}}\!=\!\mbox{sign}(y_i-\hat{\mu}_i)\sqrt{2\left\{y_i
%\ln\frac{y_i}{\hat{\mu}_i}-(y_i-\hat{\mu}_i)\right\}} \\
%\mbox{binomial}\!\!\! & \mbox{Bin}(m,\mu): & r_{_{Di}}\!=\!\mbox{sign}(y_i-\hat{\mu}_i)\sqrt{2
%\left\{y_i\ln\frac{y_i}{\hat{\mu}_i}\!+\!(\!m_i\!-\!y_i\!)\ln\frac{m_i\!-\!y_i}{m_i\!-\!\hat{\mu}_i}\right\}} \\
%\mbox{gamma} \!\!\!& \mbox{Gam}(\mu,\nu): & r_{_{Di}}\!=\!\mbox{sign}(y_i-\hat{\mu}_i)\sqrt{2
%\left\{-\ln\frac{y_i}{\hat{\mu}_i}+\frac{y_i-\hat{\mu}_i}{\hat{\mu}_i}\right\}}
%\end{eqnarray*}}
%\end{itemize}
%\end{frame}
%
%\begin{frame}{\S2.2.3 GLM residuals (3)}
%\begin{itemize}
%\item
%\textbf{Working residual}: $r_{_{Wi}}=(y_i-\hat{\mu}_i)\frac{d\eta_i}{d\mu_i}\mid_{\mu_i=\hat{\mu}_i}
%\approx g(y_i)-g(\hat{\mu}_i)$, \\
%where $g(\mu_i)$ is the link function and $\eta_i=\textbf{z}_i^T\mbox{\boldmath $\beta$}$ is the linear predictor in GLM.
%
%The working residual is used during computing the MLE of $\mbox{\boldmath $\beta$}$.
%\item
%\textbf{response residual}: $r_{_{Ri}}=y_i-\hat{\mu}_i$, \\
%which is not a right measure except when the GLM reduces to a linear model.
%\end{itemize}
%\end{frame}
%
%\begin{frame}[fragile] \frametitle{\S2.2.3 Example: Caesarian birth study (1)}
%For the data of Caesarian birth study we combine categories I and II of the response variable into
%one category \texttt{Infection}. Then fit logistic models to the reorganized data where the new response
%variable has only two categories \texttt{Infection} and \texttt{Non-infection}.
%
%\begin{scriptsize}
%\begin{Schunk}
%\begin{Sinput}
%> Caes.dat <- read.csv("D:/MAST90084/Example2-Caes.csv")
%> is.data.frame(Caes.dat)
%\end{Sinput} 
%\begin{Soutput}
%[1] TRUE
%\end{Soutput}
%\begin{Sinput}
%> Caes.dat
%\end{Sinput} 
%\begin{Soutput}
%  Infection total Antib RiskF Plan
%1         1    18     1     1    1
%2         0     2     1     0    1
%3        28    58     0     1    1
%4         8    40     0     0    1
%5        11    98     1     1    0
%6         0     0     1     0    0
%7        23    26     0     1    0
%8         0     9     0     0    0
%\end{Soutput}
%\end{Schunk}
%\end{scriptsize}
%\end{frame}
%
%\begin{frame}[fragile] \frametitle{\S2.2.3 Example: Caesarian birth study (2)}
%\begin{scriptsize}
%\begin{Schunk}
%\begin{Sinput}
%> Caes1=glm(Infection/total~Antib+RiskF+Plan,family=binomial, weight=total, data=Caes.dat)
%> summary(Caes1)
%\end{Sinput} 
%\begin{Soutput}
%Deviance Residuals: 
%       1         2         3         4         5         7         8  
% 0.26470  -0.15231  -0.78520   1.21563  -0.07162   1.49623  -2.56229  
%
%Coefficients:
%            Estimate Std. Error z value Pr(>|z|)    
%(Intercept)  -0.8207     0.4947  -1.659   0.0971 .  
%Antib        -3.2544     0.4813  -6.761 1.37e-11 ***
%RiskF         2.0299     0.4553   4.459 8.25e-06 ***
%Plan         -1.0720     0.4254  -2.520   0.0117 *  
%---
%Signif. codes:  0 '***' 0.001 '**' 0.01 '*' 0.05 '.' 0.1  ' ' 1
%
%(Dispersion parameter for binomial family taken to be 1)
%
%    Null deviance: 83.491  on 6  degrees of freedom
%Residual deviance: 10.997  on 3  degrees of freedom
%  (1 observation deleted due to missingness)
%AIC: 36.178
%Number of Fisher Scoring iterations: 5
%\end{Soutput}
%\end{Schunk}
%\end{scriptsize}
%\end{frame}
%
%\begin{frame}[fragile] \frametitle{\S2.2.3 Example: Caesarian birth study (3)}
%\begin{scriptsize}
%\begin{Schunk}
%\begin{Sinput}
%> anova(Caes1, test="Chi")
%\end{Sinput} 
%\begin{Soutput}
%Analysis of Deviance Table; Model: binomial, link: logit; Response: Infection/total
%Terms added sequentially (first to last)
%
%      Df Deviance Resid. Df Resid. Dev  Pr(>Chi)    
%NULL                      6     83.491              
%Antib  1   38.736         5     44.756 4.852e-10 ***
%RiskF  1   26.881         4     17.874 2.164e-07 ***
%Plan   1    6.878         3     10.997  0.008728 ** 
%\end{Soutput}
%\begin{Sinput}
%> resid(Caes1,type="deviance")
%\end{Sinput}
%\begin{Soutput}
%       1        2        3        4        5        7        8 
% 0.26470 -0.15231 -0.78520  1.21563 -0.07162  1.49623 -2.56229
%\end{Soutput}
%\begin{Sinput}
%> resid(Caes1,type="pearson")
%\end{Sinput}
%\begin{Soutput}
%       1        2        3        4        5        7        8 
% 0.27684 -0.10786 -0.78635  1.29465 -0.07141  1.38710 -1.99030
%\end{Soutput}
%\begin{Sinput}
%> sum(resid(Caes1,type="deviance")^2)    #[1] 10.99674
%> sum(resid(Caes1,type="pearson")^2)/(7-4)    #[1] 2.757724
%> 1-pchisq(sum(resid(Caes1,type="pearson")^2),3)   #[1] 0.04069089
%> 1-pchisq(sum(resid(Caes1,type="deviance")^2),3)   #[1] 0.01174351
%\end{Sinput}
%\end{Schunk}
%\end{scriptsize}
%\end{frame}
%
%\begin{frame}[fragile] \frametitle{\S2.2.3 Example: Caesarian birth study (4)}
%\begin{scriptsize}
%\begin{Schunk}
%\begin{Sinput}
%Caes2=glm(Infection/total~Antib+RiskF+Plan+RiskF:Plan,family=binomial, 
%                         weight=total, data=Caes.dat)
%summary(Caes2);  anova(Caes2,test="Chi")
%sum(resid(Caes2,type="pearson")^2)/(7-5)
%1-pchisq(sum(resid(Caes2,type="pearson")^2),2)
%1-pchisq(sum(resid(Caes2,type="deviance")^2),2)
%
%Caes3=glm(Infection/total~Antib+RiskF+Plan+RiskF:Plan,
%            family=binomial(link="probit"), weight=total, data=Caes.dat)
%summary(Caes3); anova(Caes3, test="Chi")
%
%Caes4=glm(Infection/total~Antib+RiskF+Plan+RiskF:Plan,
%           family=binomial(link="probit"), weight=total, data=Caes.dat[-6,])
%summary(Caes4); anova(Caes3, test="Chi")
%
%Caes5=glm(Infection/total~Antib+RiskF+Plan+RiskF:Plan,family=quasibinomial, 
%              weight=total, data=Caes.dat[-6,])
%summary(Caes5); anova(Caes5, test="F")
%
%sum(resid(Caes5,type="pearson")^2)/(7-5)
%
%Caes6=glm(Infection/total~Antib+RiskF+Plan+RiskF:Plan,family=quasibinomial, 
%           weight=total, data=Caes.dat, subset=c(1,2,3,4,5,7,8))
%summary(Caes6);  anova(Caes6, test="F")
%\end{Sinput} 
%\end{Schunk}
%\end{scriptsize}
%\end{frame}
%
%\section{\S2.3 Some extensions of GLM: quasi-likelihood and Bayesian method}
%\subsection{\S2.3.1 Quasi-likelihood models}
%\begin{frame}{\S2.3.1 Quasi-likelihood models}
%\begin{wideitemize}
%\item In GLM, the mean  necessarily implies a certain variance structure. 
%
%\item e.g. Poisson model: we must have $\sigma^2 = \mu $ (dispersion parameter $\phi$ forced to be 1. )
% such as the Poisson model, 
%\[
%E(y|\textbf{x})=\mu=h(\textbf{z}^T\mbox{\boldmath $\beta$}), \ \ \mbox{var}(y|\textbf{x})=\phi v(\mu)/\omega, 
%\]
%so the aassociation between $E(y|\textbf{x})$ and  $\mbox{var}(y|\textbf{x})$ confined. May be insufficient for practical use.
%Assuming that the response variable $y$ in GLM belongs to a specific exponential family
%confines the association between $E(y|\textbf{x})=\mu=h(\textbf{z}^T\mbox{\boldmath $\beta$})$
%and $\mbox{var}(y|\textbf{x})=\phi v(\mu)/\omega$ to be of a specific type, which sometimes is
%insufficient for practical use.
%\item
%Quasi-likelihood models: drop the exponential family assumption, as such, separate the mean and
%variance structure.
%
%\item No full distributional assumptions; only the first and second moments have
%to be specified. 
%
%\item Under appropriate conditions, parameters can still be estimated consistently,
%and asymptotic inference is still possible.
%\end{wideitemize}
%\end{frame}


\begin{frame}{Quasi-likelihood inference (Ex. 2.7)}
\begin{wideitemize}
\item  Three working variance-mean relationships are explored:
\begin{itemize}
\item $\sigma^2 (\mu) = \phi \mu$
\item $\sigma^2 (\mu) = \phi \mu^2$
\item $\sigma^2 (\mu)  = \mu + \theta \mu^2 $
\end{itemize}


\item \text{CellDiff.R} works out the first two variance structures.

\item We can compute the sandwich covariance matrix ``by hand", using the formula $\hat{F}^{-1} \hat{V}\hat{F}^{-1}$, where
\[\hat{V}=\sum_{i=1}^n\textbf{z}_i\textbf{z}_i^T 
D_i^2(\hat{\mbox{\boldmath $\beta$}})\cdot\frac{[y_i-h(\textbf{z}_i^T\hat{\mbox{\boldmath $\beta$}})]^2}{
\sigma_i^4(\hat{\boldsymbol \beta})}\text{ and } D_i(\hat{\boldsymbol \beta}) = \frac{\partial h}{\partial \eta} ({\bf z}_i^T \hat{\boldsymbol \beta}).\]
This should be same as the \texttt{robust.variance} reported by the \texttt{gee} object.

\end{wideitemize}
\end{frame}



%
%\begin{frame}{The variance function $\sigma^2 (\mu)= \mu + \theta \mu^2 $}
%
%The third column in Table 2.7 of F \& T assumes 
%the log-linear mean structure (exp. resp. function)
%\[
%\mu = \exp(\eta)
%\]
%
%and
% the working quadratic variance structure
%\[
%var(y) = \mu + \theta \mu^2 
%\] 
%(So dispersion $\phi$ is fixed to be $1$ already )
%\end{frame}

\begin{frame} {The variance function $\sigma^2 (\mu)= \mu + \theta \mu^2 $}
\begin{wideitemize}
\item This is in fact the variance structure of the NB distribution
\[
f(y; \kappa, \mu) = \frac{\Gamma(\kappa + y)}{\Gamma(\kappa) y!}\frac{\mu^y \kappa^\kappa}{(\mu + \kappa)^{\kappa + y}},
\]
which has mean and the variance structure
\[
E[y] = \mu, \ \ var(y) = \mu + \mu^2/\kappa
\]
(Done in Assignment 1!)

\item So our variance structure $var(y_i | {\bf x}_i) = \mu_i + \theta \mu_i^2 $ is really the variance structure of a NB distribution with $\theta = 1/\kappa$


\item Moreover, the dispersion $\phi$ is fixed at 1 for this variance function.

\item {\bf You will work out the quasi-likelihood inference for this variance structure in Assignment 2}.

\end{wideitemize}


\end{frame}

\begin{frame} {The variance function $\sigma^2 (\mu)= \mu + \theta \mu^2 $}
\begin{wideitemize}
\item NB distribution is not of the exponential family. However, {\bf when $\kappa$ is fixed}, {\bf it does become an exponential family}. So we can have a GLM with NB-distributed $y$  when $\kappa$'s fixed.
\item V \& R has written a function \texttt{negative.binomial()} in the \texttt{MASS} package, which for a given $\kappa$ creates a negative binomial \texttt{family} object that can be submitted to the $\texttt{glm}$ function for GLM estimation. 

\item Hence, we can leverage this to perform \emph{quasi-likelihood} estimation of $\beta$ when $\theta$(or $\kappa$) is fixed; \textcolor{red}{\emph{recall that, up a multiplicative free dispersion parameter, the score function of a GLM is same as the quasi-score function defined by the same mean and variance structures of that GLM}}. 

\item {\bf CAVEAT}: Confusingly, V\&R denotes the $\kappa$ here by $\theta$ in their book and the R manual of \texttt{negative.binomial()}, be {\bf EXTREMELY careful}.
\end{wideitemize}



\end{frame}


\begin{frame}{The variance function $\sigma^2 (\mu)= \mu + \theta \mu^2 $}

\begin{wideitemize}

\item Of course, even $\phi = 1$ is known, both $\boldbeta$ and $\theta$ are to be estimated.

\begin{wideitemize}
\item This can be done by alternating between estimation of $\boldbeta$ given $\theta$ and estimation of $\theta$ given $\boldbeta$. Given $\theta$, an estimate of $\boldbeta$ can be obtained by a Newton-Rhapson-like method as before. Given $\boldbeta$, $\theta$ can be estimated by some ``methods of moments" 


\item (However, I am {\bf NOT} able to find an R package that readily implements this methodology. In particular, when defining a \texttt{family} object with the function \texttt{quasi()}, the only allowable \texttt{variance} forms are \texttt{constant}, \texttt{mu(1-mu)}, \texttt{mu}, 
\texttt{mu\^{}2} and \texttt{mu\^{}3}. Hence I cannot define a quasi family with the variance form that contains an extra parameter $\theta$ to plug into the \texttt{gee()} function. )

\item You will code this iterative algorithm.
%\item According to $V \& R$ (section 7.4), estimation of $\theta$ is via ``maximum likelihood method" provided by the function \texttt{glm.nb} in the \texttt{MASS} package.
\end{wideitemize}


\end{wideitemize}



\end{frame}

\begin{frame} {The variance function $\sigma^2 (\mu)= \mu + \theta \mu^2 $}
\begin{wideitemize}
\item  For a given $\hat{\boldbeta}$, method of moment should give

\[
\hat{\theta} = \frac{1}{n - p} \sum_{i = 1}^n \frac{[(y_i - \mu_i (\hat{\boldbeta}))^2 -  \mu_i (\hat{\boldbeta})]}{ \mu_i (\hat{\boldbeta})^2}
\]
(In the same spirit as the method of moments estimate for $\phi$ )
\item For given $\hat{\theta}$, we estimate $\hat{\boldbeta}$ by solving the score function of a NB GLM with the given $\hat{\theta}$; this is done by \texttt{glm()} and \texttt{negative.binomial()} in \texttt{R}. 

\item Iterate until $\hat{\theta}$ converges.

\end{wideitemize}

\end{frame}



\begin{frame} {The variance function $\sigma^2 (\mu)= \mu + \theta \mu^2 $}

\begin{wideitemize}
\item Of course, we need a reasonable initialization for, say,  $\theta$ (or equivalently $\kappa= 1/\theta$). 


\item Fortunately, V \& R has written another function \texttt{glm.nb} in the MASS package, a modification of \texttt{glm}, that produces MLE of $\boldbeta$ and $\kappa$ based on NB distributional assumption on $y$. We  can  take the reasonable $\kappa$ such produced to initialize our alternating algorithm.


\item  (Not sure \texttt{glm.nb} is a good name for such a function, because NB doesn't belong to the exponential family when $\theta$ is not fixed.)

\item {\bf CAVEAT}: We should NOT just treat the $(\boldbeta, \theta)$ produced by \texttt{glm.nb} as our final estimates; they are produced based on NB distributional assumption on $y$. Our alternating procedure, although leveraging the functions \texttt{glm} and  \texttt{negative.binomial}, only assumes the mean and variance structures, which is all moment-based in spirit. 

%\item ( Go to  \texttt{CellDiff.R} and  the documentation of \texttt{glm.nb} together)
\end{wideitemize}
\end{frame}

\begin{frame} {The variance function $\sigma^2 (\mu)= \mu + \theta \mu^2 $}

Also read p.58 of F \& T when in doubt!

\end{frame}


% bibliography
\begin{frame}[allowframebreaks]
  \frametitle{References}    
  \begin{thebibliography}{10}    


%  \beamertemplatearticlebibitems
%  \bibitem[Smyth, 2003]{smyth2003pearson}
%    Smyth, Gordon K
%    \newblock Pearson's goodness of fit statistic as a score test statistic
%    \newblock {\em Lecture notes-monograph series (2003): 115-126.}
    

    

       \beamertemplatearticlebibitems
  \bibitem[McCullagh and Nelder 1989]{GLM1989}
 P. McCullagh and J.A. Nelder
    \newblock  Generalized Linear Models
    \newblock {\em Chapman \& Hall (1989).}

      \beamertemplatearticlebibitems
  \bibitem[van der Vaart 1998]{VDV1998}
 A. W. van der Vaart
    \newblock  Asymptotic statistics
    \newblock {\em Cambridge University Press (1998).}
       
       \beamertemplatearticlebibitems
  \bibitem[Wedderburn 1974]{wedderburn1974}
R.W. Wedderburn
    \newblock  Quasi-likelihood functions, generalized linear models, and the Gauss—Newton method.
    \newblock {\em Biometrika 61, no. 3 (1974): 439-447.}
    
    
           \beamertemplatearticlebibitems
  \bibitem[Gourieroux, Monfort and Trognon 1984]{GMT1984}
Gourieroux C, Monfort A, Trognon A.
    \newblock  Pseudo maximum likelihood methods: Theory.
    \newblock {\em Econometrica: journal of the Econometric Society (1984): 681-700.}
    
  \end{thebibliography}
\end{frame}

\end{document}

